% Application Of Scalar And Vector Products — Further Mathematics Upper Sixth — Cours signé OnBuch+
% Official MINESEC syllabus. Compiler avec : tectonic FMA15-application-of-scalar-and-vector-products.tex
\newcommand{\DOCMATIERE}{Further Mathematics}
\newcommand{\DOCNIVEAU}{Upper Sixth}
\newcommand{\DOCMODULE}{Upper Sixth — Further mathematics}
\newcommand{\DOCLECON}{15}
\newcommand{\DOCTITRE}{Application Of Scalar And Vector Products}
% OnBuch+ — préambule « cours signés » ENRICHI (compilé avec tectonic / XeLaTeX)
% Système visuel « Pop » d'OnBuch+ : fond crème (jamais blanc), bordures encre
% épaisses, ombres « dures » décalées, titres Archivo Black en capitales.
% Version MATHÉMATIQUES : graphes (pgfplots/tikz), boîtes pédagogiques
% (définition, propriété, méthode, attention, le savais-tu, exercice résolu,
% exercice, TP, bilan), page de garde et signature OnBuch+ ancrée en bas.
%
% Macros attendues dans main.tex AVANT \input{preamble} :
%   \DOCMATIERE  \DOCNIVEAU  \DOCMODULE  \DOCLECON (numéro)  \DOCTITRE
\documentclass[11pt]{article}

\usepackage{fontspec}
\usepackage[english]{babel}
\usepackage[a4paper,margin=2cm,top=3.4cm,bottom=2.8cm,headheight=1.4cm,footskip=1.3cm]{geometry}
\usepackage{amsmath,amssymb,amsfonts}
\usepackage{mathtools} % \xmapsto, \coloneqq... souvent utilisés par les modèles
\usepackage{cancel} % simplifications barrées (bilans de forces)
% \abs{z} (module, valeur absolue) et \norm{u} (norme), utilisés par les modèles
\providecommand{\abs}{}\let\abs\relax\DeclarePairedDelimiter{\abs}{\lvert}{\rvert}
\providecommand{\norm}{}\let\norm\relax\DeclarePairedDelimiter{\norm}{\lVert}{\rVert}
\usepackage[table]{xcolor} % [table] : fournit aussi \rowcolors (alternance de lignes)
\usepackage{pagecolor}
\usepackage{tikz}
\usetikzlibrary{arrows.meta,positioning,calc,shapes.geometric,shapes.misc,shapes.arrows,decorations.pathmorphing,decorations.pathreplacing,decorations.markings,patterns,shadows,mindmap,trees,fit,backgrounds,matrix,intersections,angles,quotes}
\usepackage{pgfplots}
\usepackage[version=4]{mhchem} % équations nucléaires \ce{^{235}_{92}U}
\usepackage[european,siunitx]{circuitikz} % schémas électriques
\pgfplotsset{compat=1.18}
\usepgfplotslibrary{fillbetween,groupplots} % aires entre courbes (calcul intégral)
\usepackage{enumitem}
\usepackage{fancyhdr}
% [most] charge toutes les bibliothèques tcolorbox (listings, minted, raster,
% documentation...) et gonfle inutilement la mémoire TeX sur un document long
% (~300 boîtes) : on ne charge que ce qui est réellement utilisé.
\usepackage[skins,breakable]{tcolorbox}
\usepackage{graphicx}
\usepackage{colortbl}
\usepackage{booktabs}
\usepackage{tabularx}
\usepackage{diagbox} % case d'en-tête diagonale des tableaux à double entrée
% Colonnes extensibles centrée / gauche / droite (C, L, R), souvent utilisées par les modèles
\newcolumntype{C}{>{\centering\arraybackslash}X}
\newcolumntype{L}{>{\raggedright\arraybackslash}X}
\newcolumntype{R}{>{\raggedleft\arraybackslash}X}
\usepackage{array}
\usepackage{multicol}
\usepackage{siunitx}
% parse-numbers=false : accepte aussi \SI{2,5 \times 10^{-3}}{...}
\sisetup{locale=UK,output-decimal-marker={.},group-minimum-digits=4,parse-numbers=false}
\DeclareSIUnit\ohm{\ensuremath{\symup{\Omega}}} % U+2126 absent de la police
\DeclareSIUnit\division{div} % réglages d'oscilloscope (ms/div, V/div)
\DeclareSIUnit\tr{tr} % tours (vitesse de rotation, ex. \SI{1500}{\tr\per\minute})
\usepackage{qrcode}
% \degree : commande fréquemment utilisée par les modèles pour un angle ou une
% température en dehors de \SI{}{} (non fournie par les paquets chargés).
\providecommand{\degree}{^{\circ}}
% \qed (fin de démonstration), utilisé par les modèles sans amsthm
\providecommand{\qed}{\hfill$\square$}
% \barre : double barre des valeurs interdites dans un tableau de variations (tkz-tab)
\providecommand{\barre}{\ensuremath{\|}}
% environnement proof (amsthm non chargé) utilisé par les modèles
\ifdefined\proof\else\newenvironment{proof}[1][Proof]{\par\noindent\textit{#1.}\ }{\qed\par}\fi
\usepackage{lastpage}
\usepackage{hyperref}
\hypersetup{hidelinks}

% ============================================================
% Palette « Pop » OnBuch+ — mêmes hex que la classe P (plus_ui.dart)
% ============================================================
\definecolor{poporange}{HTML}{F59321}
\definecolor{popdark}{HTML}{E07A0C}
\definecolor{popcream}{HTML}{FFF6E8}
\definecolor{popink}{HTML}{1C1714}
\definecolor{poppurple}{HTML}{7A5AE0}
\definecolor{popgreen}{HTML}{2E9E6A}
\definecolor{poppink}{HTML}{E0457B}
\definecolor{popblue}{HTML}{2F7DE1}
\definecolor{popgold}{HTML}{C98A12}
\definecolor{popmuted}{HTML}{8A7F76}
\definecolor{popline}{HTML}{EADFD2}
\definecolor{popgray}{HTML}{8A7F76}
\colorlet{popgrey}{popgray}
% Alias tolérants (le modèle nomme parfois les couleurs différemment)
\colorlet{popwhite}{white}
\colorlet{popblack}{popink}
\colorlet{popred}{poppink}
\colorlet{popyellow}{popgold}
% Teintes claires pour fonds de tableaux / zones
\colorlet{popblueL}{popblue!10!white}
\colorlet{poporangeL}{poporange!14!white}
\colorlet{popgreenL}{popgreen!12!white}
\colorlet{poppurpleL}{poppurple!12!white}
\colorlet{poppinkL}{poppink!12!white}
\colorlet{popgoldL}{popgold!14!white}

\pagecolor{popcream}

% ============================================================
% Typographie
% ============================================================
\defaultfontfeatures{Ligatures=TeX}
\setmainfont{PlusJakartaSans-Medium.ttf}[Path=../../../fonts/,BoldFont=PlusJakartaSans-Bold.ttf]
% Maths en sans-serif (Fira Math) avec chiffres et lettres droites de Plus
% Jakarta, pour que les formules se fondent dans le texte.
\usepackage[math-style=ISO,bold-style=ISO]{unicode-math}
\setmathfont{FiraMath-Regular.otf}[Path=../../../fonts/]
\setmathfont{PlusJakartaSans-Medium.ttf}[Path=../../../fonts/,range={up/num,up/{Latin,latin},\mathup}]
% (icomma retiré : décimales avec point en anglais)
% Caractères mathématiques Unicode absents des polices de texte (Plus Jakarta,
% Archivo Black) : rendus par leur équivalent mathématique, en texte comme en
% formule — sinon ils s'affichent en carré vide (ex. « Arithmétique dans ℤ »).
\usepackage{newunicodechar}
\newunicodechar{ℝ}{\ensuremath{\mathbb{R}}}
\newunicodechar{ℤ}{\ensuremath{\mathbb{Z}}}
\newunicodechar{ℂ}{\ensuremath{\mathbb{C}}}
\newunicodechar{ℕ}{\ensuremath{\mathbb{N}}}
\newunicodechar{ℚ}{\ensuremath{\mathbb{Q}}}
\newunicodechar{𝔻}{\ensuremath{\mathbb{D}}}
\newunicodechar{α}{\ensuremath{\alpha}}
\newunicodechar{β}{\ensuremath{\beta}}
\newunicodechar{γ}{\ensuremath{\gamma}}
\newunicodechar{δ}{\ensuremath{\delta}}
\newunicodechar{ε}{\ensuremath{\varepsilon}}
\newunicodechar{θ}{\ensuremath{\theta}}
\newunicodechar{λ}{\ensuremath{\lambda}}
\newunicodechar{μ}{\ensuremath{\mu}}
\newunicodechar{σ}{\ensuremath{\sigma}}
\newunicodechar{τ}{\ensuremath{\tau}}
\newunicodechar{φ}{\ensuremath{\varphi}}
\newunicodechar{ψ}{\ensuremath{\psi}}
\newunicodechar{ω}{\ensuremath{\omega}}
\newunicodechar{Σ}{\ensuremath{\Sigma}}
% majuscules produites par \MakeUppercase dans les titres (α-aminés -> Α)
\newunicodechar{Α}{\ensuremath{\alpha}}
\newunicodechar{Β}{\ensuremath{\beta}}
\newunicodechar{Γ}{\ensuremath{\Gamma}}
\newunicodechar{Δ}{\ensuremath{\Delta}}
\newunicodechar{Ε}{\ensuremath{\varepsilon}}
\newunicodechar{Θ}{\ensuremath{\Theta}}
\newunicodechar{Λ}{\ensuremath{\Lambda}}
\newunicodechar{Μ}{\ensuremath{\mu}}
\newunicodechar{Τ}{\ensuremath{\tau}}
\newunicodechar{Φ}{\ensuremath{\Phi}}
\newunicodechar{Ψ}{\ensuremath{\Psi}}
\newunicodechar{Ω}{\ensuremath{\Omega}}
\newunicodechar{↦}{\ensuremath{\mapsto}}
\newunicodechar{∈}{\ensuremath{\in}}
\newunicodechar{∉}{\ensuremath{\notin}}
\newunicodechar{∧}{\ensuremath{\wedge}}
\newunicodechar{⇔}{\ensuremath{\Leftrightarrow}}
\newunicodechar{⟺}{\ensuremath{\Longleftrightarrow}}
\newunicodechar{⇒}{\ensuremath{\Rightarrow}}
\newunicodechar{⟹}{\ensuremath{\Longrightarrow}}
\newunicodechar{∘}{\ensuremath{\circ}}
\newunicodechar{∪}{\ensuremath{\cup}}
\newunicodechar{∩}{\ensuremath{\cap}}
\newunicodechar{≡}{\ensuremath{\equiv}}
\newunicodechar{⊂}{\ensuremath{\subset}}
\newunicodechar{⊥}{\ensuremath{\perp}}
\newunicodechar{∃}{\ensuremath{\exists}}
\newunicodechar{∀}{\ensuremath{\forall}}
\newunicodechar{∛}{\ensuremath{\sqrt[3]{\,}}}
\newunicodechar{∜}{\ensuremath{\sqrt[4]{\,}}}
\newunicodechar{⃗}{\ensuremath{{}^{\to}}}
\newunicodechar{ⁿ}{\ifmmode^{n}\else\textsuperscript{\ensuremath{n}}\fi}
\newunicodechar{ˣ}{\ifmmode^{x}\else\textsuperscript{\ensuremath{x}}\fi}
\newunicodechar{ᵇ}{\ifmmode^{b}\else\textsuperscript{\ensuremath{b}}\fi}
\newunicodechar{ʳ}{\ifmmode^{r}\else\textsuperscript{\ensuremath{r}}\fi}
\newunicodechar{ᵘ}{\ifmmode^{u}\else\textsuperscript{\ensuremath{u}}\fi}
\newunicodechar{ʸ}{\ifmmode^{y}\else\textsuperscript{\ensuremath{y}}\fi}
\newunicodechar{ᵃ}{\ifmmode^{a}\else\textsuperscript{\ensuremath{a}}\fi}
\newunicodechar{ᵉ}{\ifmmode^{e}\else\textsuperscript{\ensuremath{e}}\fi}
\newunicodechar{ᵏ}{\ifmmode^{k}\else\textsuperscript{\ensuremath{k}}\fi}
\newunicodechar{ᵖ}{\ifmmode^{p}\else\textsuperscript{\ensuremath{p}}\fi}
\newunicodechar{ᴺ}{\ifmmode^{N}\else\textsuperscript{\ensuremath{N}}\fi}
\newunicodechar{ᶜ}{\ifmmode^{c}\else\textsuperscript{\ensuremath{c}}\fi}
\newunicodechar{ʰ}{\ifmmode^{h}\else\textsuperscript{\ensuremath{h}}\fi}
\newunicodechar{ᵠ}{\ifmmode^{q}\else\textsuperscript{\ensuremath{q}}\fi}
\newunicodechar{ₙ}{\ifmmode_{n}\else\textsubscript{\ensuremath{n}}\fi}
\newunicodechar{ₐ}{\ifmmode_{a}\else\textsubscript{\ensuremath{a}}\fi}
\newunicodechar{ᵢ}{\ifmmode_{i}\else\textsubscript{\ensuremath{i}}\fi}
\newunicodechar{ᵣ}{\ifmmode_{r}\else\textsubscript{\ensuremath{r}}\fi}
\newunicodechar{ₖ}{\ifmmode_{k}\else\textsubscript{\ensuremath{k}}\fi}
\newunicodechar{ᵦ}{\ifmmode_{\beta}\else\textsubscript{\ensuremath{\beta}}\fi}
\newunicodechar{ₓ}{\ifmmode_{x}\else\textsubscript{\ensuremath{x}}\fi}
\newfontfamily\popdisplay{ArchivoBlack-Regular.ttf}[Path=../../../fonts/]
\newfontfamily\poplogo{SpaceGrotesk-Bold.ttf}[Path=../../../fonts/]
\newfontfamily\popsemi{PlusJakartaSans-SemiBold.ttf}[Path=../../../fonts/]
\newfontfamily\popxbold{PlusJakartaSans-ExtraBold.ttf}[Path=../../../fonts/]
\newfontfamily\popxboldls{PlusJakartaSans-ExtraBold.ttf}[Path=../../../fonts/,LetterSpace=3.0]

\setlist[enumerate]{leftmargin=1.7em,itemsep=5pt,topsep=4pt}
\setlist[itemize]{leftmargin=1.7em,itemsep=4pt,topsep=4pt,label={\color{poporange}\textbullet}}
\setlength{\parindent}{0pt}
\setlength{\parskip}{5pt}
\renewcommand{\arraystretch}{1.25}

% pgfplots : style Pop par défaut
\pgfplotsset{
  every axis/.append style={axis line style={popink,line width=1pt},tick style={popink},
    grid style={popline,line width=0.5pt},label style={font=\small},tick label style={font=\footnotesize}},
  popcurve/.style={poporange,line width=1.8pt,no markers,smooth},
}
% Même style disponible dans un \draw TikZ simple (hors pgfplots)
\tikzset{popcurve/.style={poporange,line width=1.8pt}}
\tikzset{popgrid/.style={popline,very thin}}
\colorlet{popcurve}{poporange} % le nom du style est parfois utilisé comme couleur

% ============================================================
% Pill / squiggle
% ============================================================
\newcommand{\poppill}[3]{%
  \tikz[baseline=-0.55ex]\node[draw=popink,line width=1.2pt,rounded corners=2.6pt,%
    fill=#2,inner xsep=6.5pt,inner ysep=3pt]{%
    \popxboldls\fontsize{7}{7}\selectfont\color{#3}\MakeUppercase{#1}};%
}
\newcommand{\popsquiggle}[1]{%
  \tikz[baseline=-0.4ex]{
    \draw[#1,line width=2.6pt,line cap=round]
      plot [smooth] coordinates {(0,0.09) (0.34,0.01) (0.68,0.21) (1.02,0.01) (1.36,0.10)};
  }%
}
% Mot-clé surligné (marqueur orange sous le texte)
\newcommand{\cle}[1]{{\popxbold\color{popdark}#1}}

% ============================================================
% En-tête / pied
% ============================================================
\pagestyle{fancy}
\fancyhf{}
\renewcommand{\headrulewidth}{0pt}
\renewcommand{\footrulewidth}{0pt}
\lhead{\raisebox{-0.15ex}{\poplogo\fontsize{13}{13}\selectfont\color{popink}OnBuch\color{poporange}+}}
\rhead{\poppill{\DOCMATIERE\ \textbullet\ \DOCNIVEAU\ \textbullet\ Chapter \DOCLECON}{white}{popink}}
\lfoot{\popxbold\fontsize{7.6}{7.6}\selectfont\color{popmuted}\MakeUppercase{Signed course OnBuch+}\ \textbullet\ {\popsemi\DOCTITRE}}
\rfoot{\tikz[baseline=-0.6ex]\node[rounded corners=8.5pt,fill=popink,minimum height=17pt,minimum width=17pt,inner xsep=5pt,inner sep=0]{\popxbold\fontsize{8}{8}\selectfont\color{popcream}\thepage\,/\,\pageref*{LastPage}};}

% ============================================================
% Titres
% ============================================================
\newcommand{\chaptitre}[2]{%
  \begin{tcolorbox}[enhanced,colback=poporange,colframe=popink,boxrule=2.6pt,arc=11pt,
      drop shadow=popink,left=15pt,right=15pt,top=12pt,bottom=11pt,before skip=3pt,after skip=18pt]
    \poppill{#2}{popink}{popcream}\par\vspace{9pt}
    {\raggedright\hyphenpenalty=10000\exhyphenpenalty=10000\popdisplay\fontsize{22}{24}\selectfont\color{popcream}\MakeUppercase{#1}\par}
  \end{tcolorbox}
}
% Section numérotée : pastille orange avec le numéro + titre Archivo
\newcounter{popsec}
\newcommand{\coursec}[1]{%
  \stepcounter{popsec}\setcounter{popsub}{0}%
  \par\vspace{16pt}\noindent
  \begin{minipage}{\linewidth}
  \tikz[baseline=-0.7ex]\node[draw=popink,line width=1.6pt,fill=poporange,rounded corners=4pt,minimum size=22pt,inner sep=0]{\popdisplay\fontsize{12}{12}\selectfont\color{popcream}\thepopsec};%
  \hspace{8pt}{\popdisplay\fontsize{15}{17}\selectfont\color{popink}\MakeUppercase{#1}}\par
  \vspace{4pt}\noindent{\color{poporange}\rule{36pt}{3.6pt}}
  \end{minipage}\par\nopagebreak\vspace{8pt}
}
\newcounter{popsub}
\newcommand{\courssub}[1]{%
  \stepcounter{popsub}%
  \par\vspace{9pt}\noindent{\popxbold\fontsize{11.5}{13}\selectfont\color{popink}{\color{poporange}\thepopsec.\thepopsub}\hspace{6pt}#1}\par\nopagebreak\vspace{4pt}
}

% ============================================================
% Boîtes pédagogiques — même recette : fond blanc, bordure épaisse colorée,
% ombre dure encre, badge plein. Titre optionnel [#1].
% ============================================================
\tcbset{popbase/.style={breakable,enhanced,colback=white,boxrule=2.3pt,arc=9pt,
  shadow={3pt}{-3pt}{0pt}{popink},left=14pt,right=14pt,top=11pt,bottom=11pt,before skip=11pt,after skip=15pt,
  fontupper=\popsemi\fontsize{10.6}{14.5}\selectfont\color{popink},
  fontlower=\popsemi\fontsize{10.6}{14.5}\selectfont\color{popink}}}
% Coche dessinée (le glyphe ✓ n'existe pas dans les polices du document)
\newcommand{\popcheck}[1]{\tikz[baseline=-0.5ex]\draw[#1,line width=1.6pt,line cap=round,line join=round](-0.13,0.0)--(-0.04,-0.09)--(0.13,0.1);}
\providecommand{\coche}{\popcheck{popgreen}}
\providecommand{\resultat}[1]{\par\smallskip\noindent\fcolorbox{popgreen}{popgreenL}{\parbox{\dimexpr\linewidth-2\fboxsep-2\fboxrule\relax}{\textbf{Result.} #1}}\par\smallskip}
\newcommand{\popboxhead}[3]{\poppill{#1}{#2}{white}\ifx\relax#3\relax\else\hspace{6pt}{\popxbold\fontsize{10.6}{12}\selectfont\color{popink}#3}\fi\par\vspace{7pt}}

\newtcolorbox{aretenir}[1][]{popbase,colframe=popgreen,before upper={\popboxhead{Key points}{popgreen}{#1}}}
\newtcolorbox{exemplebox}[1][]{popbase,colframe=poporange,before upper={\popboxhead{Exemple}{poporange}{#1}}}
\newtcolorbox{definition}[1][]{popbase,colframe=popblue,before upper={\popboxhead{Definition}{popblue}{#1}}}
\newtcolorbox{propriete}[1][]{popbase,colframe=poppurple,before upper={\popboxhead{Property}{poppurple}{#1}}}
\newtcolorbox{methode}[1][]{popbase,colframe=popgold,before upper={\popboxhead{Method}{popgold}{#1}}}
\newtcolorbox{attention}[1][]{popbase,colframe=poppink,before upper={\popboxhead{Warning}{poppink}{#1}}}
\newtcolorbox{experience}[1][]{popbase,colframe=popblue,colback=popblueL,before upper={\popboxhead{Experiment}{popblue}{#1}}}
% « Le savais-tu ? » : carte violette pleine (contexte, culture, Cameroun)
\newtcolorbox{savaistu}[1][]{popbase,colback=poppurple,colframe=popink,
  fontupper=\popsemi\fontsize{10.4}{14.2}\selectfont\color{white},
  before upper={\poppill{Did you know?}{white}{poppurple}\ifx\relax#1\relax\else\hspace{6pt}{\popxbold\color{white}#1}\fi\par\vspace{7pt}}}
% Exercice résolu : énoncé en haut, \tcblower puis la solution sur fond crème
\newcounter{popexo}\newcounter{popexr}
\newtcolorbox{exoresolu}[1][]{popbase,colframe=poporange,colbacklower=popcream,
  segmentation style={popink,line width=1.2pt,dashed},
  before upper={\stepcounter{popexr}\popboxhead{Worked example \thepopexr}{poporange}{#1}},
  before lower={\poppill{Solution}{popink}{popcream}\par\vspace{6pt}}}
% Exercice (non résolu en place) — la correction est dans la partie Corrigés
\newtcolorbox{exercice}[1][]{popbase,colframe=popink,
  before upper={\stepcounter{popexo}\popboxhead{Exercise \thepopexo}{popink}{#1}}}
% Corrigé
\newtcolorbox{corrige}[1][]{popbase,colframe=popgreen,colback=popgreenL,
  before upper={\popboxhead{Solution}{popgreen}{#1}}}
% Objectifs / prérequis / bilan
\newtcolorbox{objectifs}{popbase,colframe=popink,colback=poporangeL,
  before upper={\popboxhead{Chapter objectives}{popink}{}}}
\newtcolorbox{prerequis}{popbase,colframe=popink,colback=popblueL,
  before upper={\popboxhead{Prerequisites}{popblue}{}}}
\newtcolorbox{bilan}{popbase,colframe=popink,colback=white,boxrule=2.8pt,
  before upper={\popboxhead{Fiche bilan}{poporange}{L'essentiel à connaître pour l'examen}}}
% Figure encadrée avec légende
\newtcolorbox{popfigure}[1][]{enhanced,breakable=false,colback=white,colframe=popline,boxrule=1.4pt,arc=8pt,
  left=10pt,right=10pt,top=10pt,bottom=8pt,before skip=10pt,after skip=12pt,halign=center,
  lower separated=false,halign lower=center,
  fontlower=\popsemi\fontsize{9}{11.5}\selectfont\color{popmuted},
  #1}
\newcommand{\legende}[1]{\par\vspace{6pt}{\popsemi\fontsize{9}{11.5}\selectfont\color{popmuted}#1}}

% ============================================================
% Signature OnBuch+
% ============================================================
\newcommand{\poplogoinline}[1]{{\poplogo\fontsize{#1}{#1}\selectfont\color{popink}OnBuch\color{poporange}+}}
\newcommand{\signatureonbuch}{%
  \begin{tcolorbox}[enhanced,colback=popink,colframe=popink,boxrule=2.2pt,arc=10pt,
      drop shadow=poporange,left=14pt,right=14pt,top=10pt,bottom=10pt,before skip=0pt,after skip=0pt]
    \tikz[baseline=-0.7ex]\node[circle,fill=popgreen,draw=popcream,line width=1.3pt,minimum size=22pt,inner sep=0]{\popcheck{white}};%
    \hspace{9pt}\begin{minipage}[c]{0.62\linewidth}
      {\popxbold\fontsize{12}{13}\selectfont\color{popcream}Signed\ \textbullet\ {\poplogo OnBuch\color{poporange}+}}\par\vspace{2pt}
      {\raggedright\popsemi\fontsize{8}{10.5}\selectfont\color{popline}\MakeUppercase{OnBuch+ teaching team}\\\MakeUppercase{Follows the official MINESEC syllabus}\par}
    \end{minipage}\hfill
    {\poplogo\fontsize{22}{22}\selectfont\color{popcream}OnBuch\color{poporange}+}
  \end{tcolorbox}%
}

% ============================================================
% Page de garde
% #1 titre · #2 matière · #3 niveau · #4 module · #5 n° leçon · #6 durée ·
% #7 description / objectifs (texte ou itemize)
% ============================================================
\newcommand{\pagegarde}[7]{%
  \thispagestyle{empty}
  \newgeometry{a4paper,margin=1.7cm,top=1.6cm,bottom=1.4cm}%
  \begin{tikzpicture}[remember picture,overlay]
    \fill[poporange!18] ([xshift=-3.2cm,yshift=-3.6cm]current page.north east) circle (4.6cm);
    \fill[poppurple!14] ([xshift=2.4cm,yshift=7.6cm]current page.south west) circle (3.8cm);
  \end{tikzpicture}%
  {\poplogo\fontsize{30}{30}\selectfont\color{popink}OnBuch\color{poporange}+}\hfill
  \raisebox{0.6ex}{\poppill{Signed course \textbullet\ Premium}{popink}{popcream}}\par
  \vspace{1pt}\popsquiggle{poppurple}\par
  \vspace{8pt}
  \begin{tcolorbox}[enhanced,colback=poporange,colframe=popink,boxrule=2.8pt,arc=13pt,
      drop shadow=popink,left=18pt,right=18pt,top=12pt,bottom=13pt,after skip=10pt]
    \poppill{#2 \textbullet\ #3}{popink}{popcream}\hspace{5pt}\poppill{#4}{white}{popink}\par\vspace{8pt}
    {\popdisplay\fontsize{14}{14}\selectfont\color{popink}CHAPTER #5}\par\vspace{4pt}
    {\raggedright\hyphenpenalty=10000\exhyphenpenalty=10000\popdisplay\fontsize{25}{27}\selectfont\color{popcream}\MakeUppercase{#1}\par}
    \vspace{8pt}
    {\popxbold\fontsize{9.5}{9.5}\selectfont\color{popink}\MakeUppercase{Suggested time: #6}}
  \end{tcolorbox}
  \begin{tcolorbox}[enhanced,colback=white,colframe=popink,boxrule=2.2pt,arc=10pt,
      drop shadow=popink,left=16pt,right=16pt,top=10pt,bottom=10pt,after skip=8pt]
    \poppill{In this chapter}{poppurple}{white}\par\vspace{5pt}
    \popsemi\fontsize{9.3}{12.6}\selectfont\color{popink}#7
  \end{tcolorbox}
  \noindent\begin{minipage}[t]{0.487\linewidth}
    \begin{tcolorbox}[enhanced,colback=popgreen,colframe=popink,boxrule=2.2pt,arc=10pt,
        drop shadow=popink,left=11pt,right=11pt,top=10pt,bottom=10pt,height=3.1cm,valign=center]
      \tcbox[colback=white,colframe=popink,boxrule=1.3pt,arc=5pt,boxsep=0pt,nobeforeafter,
        left=3pt,right=3pt,top=3pt,bottom=3pt,valign=center]{\qrcode[height=1.9cm]{https://play.google.com/store/apps/details?id=com.luvvix.onbuch&hl=en}}%
      \hspace{8pt}\begin{minipage}[c]{0.5\linewidth}
        {\popdisplay\fontsize{11}{12.5}\selectfont\color{white}\MakeUppercase{Download OnBuch}}\par\vspace{4pt}
        {\raggedright\popsemi\fontsize{8.4}{11}\selectfont\color{white}Free \textbullet\ Google Play\\AI tutor, past papers, results\par}
      \end{minipage}
    \end{tcolorbox}
  \end{minipage}\hfill
  \begin{minipage}[t]{0.487\linewidth}
    \begin{tcolorbox}[enhanced,colback=poppurple,colframe=popink,boxrule=2.2pt,arc=10pt,
        drop shadow=popink,left=11pt,right=11pt,top=10pt,bottom=10pt,height=3.1cm,valign=center]
      \tikz[baseline=-0.7ex]\node[circle,fill=white,draw=popink,line width=1.3pt,minimum size=1.6cm,inner sep=0]{\popdisplay\fontsize{24}{24}\selectfont\color{poppurple}+};%
      \hspace{8pt}\begin{minipage}[c]{0.6\linewidth}
        {\popdisplay\fontsize{11}{12.5}\selectfont\color{white}\MakeUppercase{Go OnBuch+}}\par\vspace{4pt}
        {\raggedright\popsemi\fontsize{8.4}{11}\selectfont\color{white}All signed courses, unlimited AI tutor, PDF sheets — OnBuch+ tab in the app\par}
      \end{minipage}
    \end{tcolorbox}
  \end{minipage}
  \vspace{8pt}
  \popcontenu
  \vfill
  \signatureonbuch
  \restoregeometry
  \newpage
}

% Rangée « Ce cours contient » (page de garde)
\newcommand{\poptile}[3]{% #1 couleur #2 gros texte #3 légende
  \begin{tcolorbox}[enhanced,colback=#1,colframe=popink,boxrule=2pt,arc=9pt,shadow={2.5pt}{-2.5pt}{0pt}{popink},
      left=6pt,right=6pt,top=9pt,bottom=9pt,halign=center,valign=center,height=1.75cm,width=\linewidth,nobeforeafter]
    {\popdisplay\fontsize{12.5}{13}\selectfont\color{white}\MakeUppercase{#2}}\par\vspace{3pt}
    {\centering\popsemi\fontsize{7.8}{9.5}\selectfont\color{white}#3\par}
  \end{tcolorbox}}
\newcommand{\popcontenu}{%
  \par\noindent{\popxboldls\fontsize{7.5}{7.5}\selectfont\color{popmuted}THIS SIGNED COURSE CONTAINS}\par\vspace{7pt}
  \noindent\begin{minipage}[t]{0.235\linewidth}\poptile{popblue}{Course}{complete, illustrated, step by step}\end{minipage}\hfill
  \begin{minipage}[t]{0.235\linewidth}\poptile{poporange}{Methods}{and worked examples}\end{minipage}\hfill
  \begin{minipage}[t]{0.235\linewidth}\poptile{poppink}{Exercises}{exam style + solutions}\end{minipage}\hfill
  \begin{minipage}[t]{0.235\linewidth}\poptile{popgreen}{Summary sheet}{the essentials to remember}\end{minipage}\par}

% Sommaire
\newtcolorbox{sommaire}{enhanced,colback=white,colframe=popink,boxrule=2.2pt,arc=10pt,
  drop shadow=popink,left=18pt,right=18pt,top=13pt,bottom=13pt,before skip=6pt,after skip=20pt,
  before upper={\poppill{Contents}{poppurple}{white}\par\vspace{9pt}\popsemi\fontsize{10.6}{14.5}\selectfont\color{popink}}}

% Signature de fin de cours — poussée tout en bas de la dernière page
\newcommand{\findecours}{%
  \par\vfill\nopagebreak
  \begin{center}\popsquiggle{poporange}\end{center}\vspace{4pt}
  \signatureonbuch
}
\AtBeginDocument{\def\llbracket{\mathopen{[\mkern-3.5mu[}}\def\rrbracket{\mathclose{]\mkern-3.5mu]}}} % intervalles d'entiers (glyphe ⟦ absent de la police)
% Glyphes absents de Fira Math : redéfinis à partir de glyphes disponibles
\AtBeginDocument{%
  \def\setminus{\mathbin{\backslash}}\let\smallsetminus\setminus
  \def\vdots{\mathord{\vbox{\baselineskip3.2pt\lineskiplimit0pt\kern4pt\hbox{.}\hbox{.}\hbox{.}}}}%
  \def\ddots{\mathinner{\mkern1mu\raise6.5pt\hbox{.}\mkern2mu\raise3.6pt\hbox{.}\mkern2mu\raise0.7pt\hbox{.}\mkern1mu}}%
  \def\triangle{\mathord{\scalebox{0.9}{$\symup{\Delta}$}}}%
  \def\nabla{\mathord{\raisebox{\depth}{\scalebox{1}[-1]{$\symup{\Delta}$}}}}%
  \def\checkmark{\popcheck{popdark}}%
  \def\star{\mathbin{\ast}}\def\bigstar{\mathord{\ast}}%
  \def\diamond{\mathbin{\scalebox{0.6}{$\Diamond$}}}%
  \def\bigcup{\mathop{\mathchoice{\raisebox{-0.3em}{\scalebox{1.7}{$\cup$}}}{\scalebox{1.25}{$\cup$}}{\cup}{\cup}}\displaylimits}%
  \def\bigcap{\mathop{\mathchoice{\raisebox{-0.3em}{\scalebox{1.7}{$\cap$}}}{\scalebox{1.25}{$\cap$}}{\cap}{\cap}}\displaylimits}%
  \def\lesseqqgtr{\mathrel{\lessgtr}}\def\bigoplus{\mathop{\oplus}\displaylimits}%
}
\providecommand{\ssi}{if and only if} % abréviation fréquente
\AtBeginDocument{\providecommand{\wideparen}{\overparen}} % arc de cercle

% Fonctions usuelles que les modèles utilisent sans que amsmath les fournisse
\providecommand{\arsinh}{\operatorname{arsinh}}\providecommand{\arcosh}{\operatorname{arcosh}}
\providecommand{\artanh}{\operatorname{artanh}}\providecommand{\arsech}{\operatorname{arsech}}
\providecommand{\arcsch}{\operatorname{arcsch}}\providecommand{\arcoth}{\operatorname{arcoth}}
\providecommand{\arcsinh}{\operatorname{arsinh}}\providecommand{\arccosh}{\operatorname{arcosh}}
\providecommand{\arctanh}{\operatorname{artanh}}\providecommand{\sech}{\operatorname{sech}}
\providecommand{\csch}{\operatorname{csch}}\providecommand{\arcsec}{\operatorname{arcsec}}
\providecommand{\arccsc}{\operatorname{arccsc}}\providecommand{\arccot}{\operatorname{arccot}}
\providecommand{\sgn}{\operatorname{sgn}}\providecommand{\lcm}{\operatorname{lcm}}
\providecommand{\Var}{\operatorname{Var}}\providecommand{\Cov}{\operatorname{Cov}}

\begin{document}

\pagegarde{Application Of Scalar And Vector Products}{Further Mathematics}{Upper Sixth}{Upper Sixth — Further mathematics}{15}{6 periods}{This chapter applies the scalar (dot) and vector (cross) products to mechanics: resolving vectors along directions, computing work done by forces, moments of forces, couples, and the equilibrium of coplanar and three-dimensional force systems. It provides the essential toolkit for the GCE Advanced Level mechanics questions on statics.
\vspace{4pt}
\begin{multicols}{2}\small
\begin{itemize}
\item By the end of this chapter I can find the component (resolved part) of a vector in any given direction using the scalar product.
\item By the end of this chapter I can compute the work done by a constant force along a straight displacement, including forces at an angle.
\item By the end of this chapter I can calculate the moment of a force about a point in two and three dimensions using the vector product.
\item By the end of this chapter I can apply the principle of moments to solve equilibrium problems involving levers, beams and ladders.
\item By the end of this chapter I can identify and analyse a couple, compute its moment, and distinguish a couple from a single force.
\item By the end of this chapter I can reduce a system of coplanar forces to a single resultant force or to a force and a couple.
\end{itemize}\end{multicols}}

\begin{sommaire}
\begin{multicols}{2}
\begin{enumerate}
\item Section 1: Component of a Vector in a Given Direction (Lesson 102)
\item Section 2: Work Done by a Constant Force (Lesson 103)
\item Section 3: Moment of a Force and Equilibrium (Lesson 104)
\item Section 4: Single Forces and Couples (Lesson 105)
\item Section 5: Coplanar Forces and Simple 3D Systems (Lessons 106 \& 107)
\item Integration activity
\item Methods and worked examples
\item Exercises
\item Solutions to the exercises
\item Summary sheet
\end{enumerate}
\end{multicols}
\end{sommaire}

\begin{objectifs}
\begin{itemize}
\item By the end of this chapter I can find the component (resolved part) of a vector in any given direction using the scalar product.
\item By the end of this chapter I can compute the work done by a constant force along a straight displacement, including forces at an angle.
\item By the end of this chapter I can calculate the moment of a force about a point in two and three dimensions using the vector product.
\item By the end of this chapter I can apply the principle of moments to solve equilibrium problems involving levers, beams and ladders.
\item By the end of this chapter I can identify and analyse a couple, compute its moment, and distinguish a couple from a single force.
\item By the end of this chapter I can reduce a system of coplanar forces to a single resultant force or to a force and a couple.
\item By the end of this chapter I can solve equilibrium problems for bodies under coplanar forces using resolution and moments.
\item By the end of this chapter I can analyse simple three-dimensional force systems using vector methods (resultants and vector moments).
\end{itemize}
\end{objectifs}

\begin{prerequis}
\begin{itemize}
\item Scalar (dot) product and vector (cross) product of vectors, with their geometric interpretations (Lower Sixth).
\item Vector algebra in 2D and 3D: components, magnitude, unit vectors, position vectors.
\item Resolution of forces and basic notions of equilibrium from Lower Sixth mechanics.
\item Trigonometry: sine, cosine, angles between lines, and solving simple trigonometric equations.
\item Solving simultaneous linear equations in two or three unknowns.
\end{itemize}
\end{prerequis}

\newpage

\coursec{Section 1: Component of a Vector in a Given Direction (Lesson 102)}

On the Bamenda--Bafoussam road, a heavily loaded lorry stalls on a steep slope: the driver must know how much of the lorry's weight pulls it \emph{backwards along the road}, and how hard the handbrake must work to hold it. At the Douala port, a crane operator tilts a container and wonders at what angle it will topple. In a mechanic's workshop in Yaoundé, two spanners are used together to loosen a rusted nut that a single spanner cannot move. All these situations --- components of forces, work, turning effects, couples and equilibrium --- are governed by the \cle{scalar product} and the \cle{vector product} you met in Lower Sixth. This chapter turns those algebraic tools into the language of mechanics.

Our first problem is simple to state but fundamental: \emph{given a vector $\vec{a}$ (say a force) and a direction (say a slope), how much of $\vec{a}$ acts along that direction?} The answer is the \cle{component} of a vector in a given direction, and it is the key that unlocks work, moments and equilibrium in the rest of the chapter.

\courssub{Review of the Scalar Product}

\begin{definition}[Scalar product]
For two vectors $\vec{a}$ and $\vec{b}$ with angle $\theta$ ($0 \leq \theta \leq \pi$) between them, the \cle{scalar product} (or dot product) is
\[
\vec{a} \cdot \vec{b} = |\vec{a}|\,|\vec{b}|\cos\theta.
\]
If $\vec{a} = a_1\vec{i} + a_2\vec{j} + a_3\vec{k}$ and $\vec{b} = b_1\vec{i} + b_2\vec{j} + b_3\vec{k}$, then
\[
\vec{a} \cdot \vec{b} = a_1b_1 + a_2b_2 + a_3b_3, \qquad |\vec{a}| = \sqrt{a_1^2 + a_2^2 + a_3^2}.
\]
\end{definition}

\begin{aretenir}[Essential facts about the dot product]
\begin{itemize}
\item $\vec{a} \cdot \vec{b}$ is a \textbf{scalar}, not a vector.
\item $\vec{a} \cdot \vec{b} = 0 \iff \vec{a} \perp \vec{b}$ (for non-zero vectors).
\item $\vec{a} \cdot \vec{a} = |\vec{a}|^2$.
\item $\vec{i}\cdot\vec{i} = \vec{j}\cdot\vec{j} = \vec{k}\cdot\vec{k} = 1$ and $\vec{i}\cdot\vec{j} = \vec{j}\cdot\vec{k} = \vec{k}\cdot\vec{i} = 0$.
\item The angle between two vectors is obtained from $\cos\theta = \dfrac{\vec{a}\cdot\vec{b}}{|\vec{a}||\vec{b}|}$.
\end{itemize}
\end{aretenir}

\courssub{The Scalar Component (Resolved Part)}

Imagine the sun directly overhead. The \emph{shadow} that a stick $\vec{a}$ casts on the ground along a line in the direction of $\vec{b}$ has length $|\vec{a}|\cos\theta$. This ``shadow length'' is exactly what the driver on the Bamenda road needs: the amount of $\vec{a}$ that lies along the direction $\vec{b}$.

\begin{definition}[Scalar component of $\vec{a}$ in the direction of $\vec{b}$]
The \cle{scalar component} (or \cle{resolved part}) of $\vec{a}$ in the direction of $\vec{b}$ is
\[
\operatorname{comp}_{\vec{b}}(\vec{a}) = \vec{a}\cdot\hat{\vec{b}} = \frac{\vec{a}\cdot\vec{b}}{|\vec{b}|} = |\vec{a}|\cos\theta,
\]
where $\hat{\vec{b}} = \dfrac{\vec{b}}{|\vec{b}|}$ is the \cle{unit vector} in the direction of $\vec{b}$.
\end{definition}

The derivation is immediate: since $\hat{\vec{b}}$ has length $1$,
\[
\vec{a}\cdot\hat{\vec{b}} = |\vec{a}|\,|\hat{\vec{b}}|\cos\theta = |\vec{a}|\cos\theta.
\]

\textbf{Sign of the component.} Because $\cos\theta$ is positive for $\theta < 90^{\circ}$, zero for $\theta = 90^{\circ}$ and negative for $\theta > 90^{\circ}$, the scalar component tells us not only \emph{how much} of $\vec{a}$ lies along $\vec{b}$, but also \emph{in which sense}: a negative component means $\vec{a}$ points partly \emph{opposite} to $\vec{b}$.

\courssub{The Vector Projection}

The scalar component gives a \emph{length with a sign}. Often we want the actual \emph{vector} shadow of $\vec{a}$ along the line of $\vec{b}$: this is the \cle{vector projection}.

\begin{definition}[Vector projection of $\vec{a}$ on $\vec{b}$]
The \cle{vector projection} of $\vec{a}$ on $\vec{b}$ is the vector of magnitude $\operatorname{comp}_{\vec{b}}(\vec{a})$ pointing along $\hat{\vec{b}}$:
\[
\operatorname{proj}_{\vec{b}}(\vec{a}) = \left(\vec{a}\cdot\hat{\vec{b}}\right)\hat{\vec{b}} = \frac{\vec{a}\cdot\vec{b}}{|\vec{b}|^2}\,\vec{b}.
\]
\end{definition}

\textbf{Component perpendicular to a direction.} Since $\vec{a}$ is the sum of its parts along and perpendicular to $\vec{b}$, the part of $\vec{a}$ \cle{perpendicular} to $\vec{b}$ is
\[
\vec{a} - \operatorname{proj}_{\vec{b}}(\vec{a}).
\]
This vector is automatically perpendicular to $\vec{b}$: indeed $\bigl(\vec{a} - \operatorname{proj}_{\vec{b}}(\vec{a})\bigr)\cdot\vec{b} = \vec{a}\cdot\vec{b} - \dfrac{\vec{a}\cdot\vec{b}}{|\vec{b}|^2}|\vec{b}|^2 = 0$.

\begin{attention}[Scalar component versus vector projection]
The \textbf{component} $\operatorname{comp}_{\vec{b}}(\vec{a}) = \dfrac{\vec{a}\cdot\vec{b}}{|\vec{b}|}$ is a \textbf{scalar} (it can be negative). The \textbf{projection} $\operatorname{proj}_{\vec{b}}(\vec{a}) = \dfrac{\vec{a}\cdot\vec{b}}{|\vec{b}|^2}\vec{b}$ is a \textbf{vector}. Do not confuse them: examiners award separate marks for each; writing one when the question asks for the other loses them all.
\end{attention}

\begin{attention}[The most common exam error]
The formula $\operatorname{proj}_{\vec{b}}(\vec{a}) = (\vec{a}\cdot\hat{\vec{b}})\hat{\vec{b}}$ contains $|\vec{b}|$ \textbf{twice}: once from the dot product with $\hat{\vec{b}}$ and once from multiplying by $\hat{\vec{b}}$. Forgetting to divide by $|\vec{b}|$ (or dividing only once) when $\vec{b}$ is \emph{not} a unit vector is the single most frequent error at the GCE. Always check: is $\vec{b}$ already a unit vector? If not, divide by $|\vec{b}|$ for the component, and by $|\vec{b}|^2$ for the projection.
\end{attention}

\begin{popfigure}
\begin{center}
\begin{tikzpicture}[scale=1.1, >=stealth]
\coordinate (O) at (0,0);
\coordinate (A) at (4.2,2.6);
\coordinate (B) at (5.4,0);
\coordinate (P) at (4.2,0);
\draw[->, very thick, popblue] (O) -- (A) node[above left] {$\vec{a}$};
\draw[->, very thick, poporange] (O) -- (B) node[below right] {$\vec{b}$};
\draw[->, thick, popgreen] (O) -- (1.3,0) node[midway, below] {$\hat{\vec{b}}$};
\draw[dashed, popmuted] (A) -- (P);
\draw[popmuted] (P) ++(0,0.18) -- ++(-0.18,0.18) -- ++(-0.18,0);
\draw[->, very thick, poppurple] (O) -- (P) node[midway, above, yshift=2pt] {$\operatorname{proj}_{\vec{b}}(\vec{a})$};
\draw[<->, popdark] (0,-0.55) -- (4.2,-0.55) node[midway, below] {$\operatorname{comp}_{\vec{b}}(\vec{a}) = |\vec{a}|\cos\theta$};
\draw[popdark] (1.6,0) arc (0:31.7:1.6) node[midway, right, xshift=2pt] {$\theta$};
\end{tikzpicture}
\end{center}
\legende{The vector projection of $\vec{a}$ on $\vec{b}$ (purple) and its scalar component, the signed length of the shadow of $\vec{a}$ along $\vec{b}$.}
\end{popfigure}

\courssub{Components on the Coordinate Axes: Direction Cosines}

A beautiful special case: take $\vec{b}$ to be one of the unit vectors $\vec{i}, \vec{j}, \vec{k}$. If $\vec{a} = a_1\vec{i} + a_2\vec{j} + a_3\vec{k}$, then
\[
\vec{a}\cdot\vec{i} = a_1, \qquad \vec{a}\cdot\vec{j} = a_2, \qquad \vec{a}\cdot\vec{k} = a_3.
\]
So \emph{the Cartesian components of a vector are precisely its scalar components along the axes}. If $\alpha, \beta, \gamma$ are the angles that $\vec{a}$ makes with the $x$-, $y$- and $z$-axes, then
\[
\cos\alpha = \frac{a_1}{|\vec{a}|}, \qquad \cos\beta = \frac{a_2}{|\vec{a}|}, \qquad \cos\gamma = \frac{a_3}{|\vec{a}|}.
\]
These are the \cle{direction cosines} of $\vec{a}$, and they satisfy $\cos^2\alpha + \cos^2\beta + \cos^2\gamma = 1$ (since the unit vector $\hat{\vec{a}} = (\cos\alpha)\vec{i} + (\cos\beta)\vec{j} + (\cos\gamma)\vec{k}$ has length $1$).

\courssub{Resolving a Force Along a Given Line}

\begin{methode}[Resolving a force along a given line]
To find the resolved part of a force $\vec{F}$ along a line $\ell$:
\begin{enumerate}
\item Find any vector $\vec{d}$ along the line $\ell$ (e.g.\ from two points on the line, or from the direction of an inclined plane).
\item Compute the unit vector $\hat{\vec{d}} = \dfrac{\vec{d}}{|\vec{d}|}$.
\item The resolved part of $\vec{F}$ along $\ell$ is the scalar $\vec{F}\cdot\hat{\vec{d}}$.
\item The force vector along $\ell$ is the projection $(\vec{F}\cdot\hat{\vec{d}})\,\hat{\vec{d}}$; the remainder $\vec{F} - (\vec{F}\cdot\hat{\vec{d}})\hat{\vec{d}}$ acts perpendicular to $\ell$.
\item Interpret the sign: positive means $\vec{F}$ pulls along $\hat{\vec{d}}$; negative means it pulls against $\hat{\vec{d}}$.
\end{enumerate}
\end{methode}

\begin{exoresolu}[Component and projection with position vectors]
\textbf{Statement.} Given $\vec{a} = 2\vec{i} + 3\vec{j} + 6\vec{k}$ and $\vec{b} = 2\vec{i} - \vec{j} + 2\vec{k}$, find (a) the scalar component of $\vec{a}$ in the direction of $\vec{b}$; (b) the vector projection of $\vec{a}$ on $\vec{b}$; (c) the component of $\vec{a}$ perpendicular to $\vec{b}$; (d) the angle between $\vec{a}$ and $\vec{b}$.
\tcblower
\textbf{Solution.} First compute
\[
\vec{a}\cdot\vec{b} = (2)(2) + (3)(-1) + (6)(2) = 4 - 3 + 12 = 13,
\]
\[
|\vec{b}| = \sqrt{4 + 1 + 4} = 3, \qquad |\vec{a}| = \sqrt{4 + 9 + 36} = 7.
\]
(a) Scalar component:
\[
\operatorname{comp}_{\vec{b}}(\vec{a}) = \frac{\vec{a}\cdot\vec{b}}{|\vec{b}|} = \frac{13}{3}.
\]
(b) Vector projection:
\[
\operatorname{proj}_{\vec{b}}(\vec{a}) = \frac{13}{9}\,(2\vec{i} - \vec{j} + 2\vec{k}) = \frac{26}{9}\vec{i} - \frac{13}{9}\vec{j} + \frac{26}{9}\vec{k}.
\]
(c) Perpendicular component:
\[
\vec{a} - \operatorname{proj}_{\vec{b}}(\vec{a}) = \left(2 - \tfrac{26}{9}\right)\vec{i} + \left(3 + \tfrac{13}{9}\right)\vec{j} + \left(6 - \tfrac{26}{9}\right)\vec{k} = -\frac{8}{9}\vec{i} + \frac{40}{9}\vec{j} + \frac{28}{9}\vec{k}.
\]
Check: $\left(-\tfrac{8}{9}\right)(2) + \left(\tfrac{40}{9}\right)(-1) + \left(\tfrac{28}{9}\right)(2) = \tfrac{-16 - 40 + 56}{9} = 0$, so it is indeed perpendicular to $\vec{b}$. $\checkmark$

(d) Angle:
\[
\cos\theta = \frac{13}{7 \times 3} = \frac{13}{21} \implies \theta \approx 51.8^{\circ}.
\]
\end{exoresolu}

\courssub{Application to Forces on an Inclined Plane}

Return to our lorry on the Bamenda--Bafoussam road. Its weight $\vec{W}$ acts vertically downwards, but what matters for the handbrake is the resolved part of $\vec{W}$ \emph{along the slope}. If the road makes an angle $\theta$ with the horizontal, a unit vector down the slope is $\hat{\vec{d}} = (\cos\theta)\vec{i} - (\sin\theta)\vec{j}$ (taking $\vec{i}$ horizontal, $\vec{j}$ vertical), and with $\vec{W} = -W\vec{j}$:
\[
\vec{W}\cdot\hat{\vec{d}} = W\sin\theta.
\]
So a force $W\sin\theta$ pulls the lorry down the road, while $W\cos\theta$ presses it into the road surface. This is why steep slopes demand so much more from brakes and engines.

\begin{popfigure}
\begin{center}
\begin{tikzpicture}[scale=1.0, >=stealth]
\coordinate (O) at (0,0);
\draw[thick, popdark] (-0.5,0) -- (7,0);
\draw[thick, popdark] (0,0) -- (6,2.4);
\draw[popmuted] (2.2,0) arc (0:21.8:2.2) node[midway, right, xshift=3pt] {$\theta$};
\begin{scope}[rotate around={21.8:(3,1.2)}]
\draw[fill=popblueL, draw=popblue, thick] (2.4,1.2) rectangle (3.8,2.0);
\draw[fill=popink] (2.6,1.2) circle (0.14);
\draw[fill=popink] (3.6,1.2) circle (0.14);
\end{scope}
\coordinate (C) at (3.1,2.44);
\draw[->, very thick, poporange] (C) -- ++(0,-2.6) node[right] {$\vec{W}$};
\draw[->, very thick, poppurple] (C) -- ++(-1.66,-0.66) node[above left] {$W\sin\theta$};
\draw[->, very thick, popgreen] (C) -- ++(0.96,-2.4) node[below right] {$W\cos\theta$};
\draw[dashed, popmuted] (C) ++(-1.66,-0.66) -- ++(0.96,-2.4) -- ++(1.66,0.66);
\node[popdark] at (5.6,1.0) {slope};
\end{tikzpicture}
\end{center}
\legende{The weight $\vec{W}$ of a lorry on a slope resolved into a component $W\sin\theta$ along the plane (which the brakes must balance) and $W\cos\theta$ perpendicular to the plane.}
\end{popfigure}

\begin{exoresolu}[Resolved part of a force along a line of action]
\textbf{Statement.} A force $\vec{F} = (3\vec{i} + 4\vec{j})\ \mathrm{N}$ acts on a crate at a point $O$. A rope lies along the line from $O$ to the point $A(4, 1)$ (coordinates in metres). Find (a) the resolved part of $\vec{F}$ along $OA$; (b) the force vector of $\vec{F}$ along $OA$; (c) the magnitude of the force perpendicular to the rope.
\tcblower
\textbf{Solution.} (a) $\overrightarrow{OA} = 4\vec{i} + \vec{j}$, so $|\overrightarrow{OA}| = \sqrt{17}$ and
\[
\hat{\vec{d}} = \frac{1}{\sqrt{17}}(4\vec{i} + \vec{j}), \qquad \vec{F}\cdot\hat{\vec{d}} = \frac{(3)(4) + (4)(1)}{\sqrt{17}} = \frac{16}{\sqrt{17}} \approx 3.88\ \mathrm{N}.
\]
(b) Force vector along $OA$:
\[
\frac{16}{17}(4\vec{i} + \vec{j}) = \left(\frac{64}{17}\vec{i} + \frac{16}{17}\vec{j}\right)\ \mathrm{N}.
\]
(c) Perpendicular force: $\vec{F} - \operatorname{proj} = \left(3 - \tfrac{64}{17}\right)\vec{i} + \left(4 - \tfrac{16}{17}\right)\vec{j} = -\tfrac{13}{17}\vec{i} + \tfrac{52}{17}\vec{j}$, of magnitude
\[
\frac{1}{17}\sqrt{13^2 + 52^2} = \frac{13\sqrt{17}}{17} = \frac{13}{\sqrt{17}} \approx 3.15\ \mathrm{N}.
\]
Check: $\left(\tfrac{16}{\sqrt{17}}\right)^2 + \left(\tfrac{13}{\sqrt{17}}\right)^2 = \tfrac{256 + 169}{17} = 25 = |\vec{F}|^2$. $\checkmark$
\end{exoresolu}

\begin{savaistu}[Why this matters beyond the exam room]
Every time an engineer sizes a cable on a bridge, a physiotherapist estimates the force along a bone, or a pilot computes the headwind component on landing at Yaoundé--Nsimalen airport, the calculation is exactly $\vec{a}\cdot\hat{\vec{b}}$. The scalar component is arguably the most-used formula in all of applied vector mathematics.
\end{savaistu}

\begin{exercice}[Practice]
\begin{enumerate}
\item Given $\vec{a} = \vec{i} - 2\vec{j} + 2\vec{k}$ and $\vec{b} = 3\vec{i} + 4\vec{k}$, find the scalar component and the vector projection of $\vec{a}$ on $\vec{b}$.
\item A force $\vec{F} = (5\vec{i} - 2\vec{j} + \vec{k})\ \mathrm{N}$ acts along a line in the direction of $\vec{d} = \vec{i} + 2\vec{j} - 2\vec{k}$. Find the resolved part of $\vec{F}$ along $\vec{d}$ and state whether it acts with or against $\vec{d}$.
\item A lorry of weight $W$ stands on a road inclined at $12^{\circ}$ to the horizontal. Find, in terms of $W$, the force the handbrake must provide to hold it at rest.
\item Find the direction cosines of $\vec{a} = 2\vec{i} - 3\vec{j} + 6\vec{k}$ and verify that their squares sum to $1$.
\end{enumerate}
\end{exercice}

\begin{corrige}[Exercise 1]
\begin{enumerate}
\item $\vec{a}\cdot\vec{b} = 3 + 0 + 8 = 11$, $|\vec{b}| = 5$. Component: $\tfrac{11}{5}$. Projection: $\tfrac{11}{25}(3\vec{i} + 4\vec{k})$.
\item $\vec{F}\cdot\vec{d} = 5 - 4 - 2 = -1$, $|\vec{d}| = 3$, so the resolved part is $-\tfrac{1}{3}\ \mathrm{N}$: it acts \emph{against} $\vec{d}$.
\item $W\sin 12^{\circ} \approx 0.208\,W$.
\item $|\vec{a}| = 7$; direction cosines $\tfrac{2}{7}, -\tfrac{3}{7}, \tfrac{6}{7}$; squares: $\tfrac{4 + 9 + 36}{49} = 1$. $\checkmark$
\end{enumerate}
\end{corrige}

\begin{aretenir}[Key points of Lesson 102]
\begin{itemize}
\item Scalar component: $\operatorname{comp}_{\vec{b}}(\vec{a}) = \vec{a}\cdot\hat{\vec{b}} = \dfrac{\vec{a}\cdot\vec{b}}{|\vec{b}|}$ --- a signed scalar.
\item Vector projection: $\operatorname{proj}_{\vec{b}}(\vec{a}) = \dfrac{\vec{a}\cdot\vec{b}}{|\vec{b}|^2}\vec{b}$ --- a vector along $\vec{b}$.
\item Perpendicular part: $\vec{a} - \operatorname{proj}_{\vec{b}}(\vec{a})$.
\item Cartesian components are scalar components on $\vec{i}, \vec{j}, \vec{k}$; direction cosines are components of the unit vector $\hat{\vec{a}}$.
\item To resolve a force along a line: unit vector first, then dot product.
\end{itemize}
\end{aretenir}

\coursec{Section 2: Work Done by a Constant Force (Lesson 103)}

In Section 1 we learnt how to split a vector into its \cle{component} in a given direction and its component perpendicular to that direction. We saw that the component of a vector $\vec{F}$ along a direction $\vec{d}$ is $|\vec{F}|\cos\theta$, where $\theta$ is the angle between them. We now put this idea to work in one of the most important applications of the scalar product in mechanics: measuring the \cle{work} done by a force. The key observation is simple: only the component of the force \emph{along the displacement} actually does work. That component is exactly what the scalar product $\vec{F}\cdot\vec{d}$ captures.

\courssub{Intuition: what does "work" mean in mechanics?}

Imagine a porter at the Douala market pulling a loaded cart with a rope. The rope makes an angle $\theta$ with the horizontal ground, and the cart moves a distance $|\vec{d}|$ in a straight line. Everyday experience tells us two things:

\begin{itemize}
\item the \emph{harder} the porter pulls (bigger $|\vec{F}|$), the more work is done;
\item the \emph{farther} the cart moves (bigger $|\vec{d}|$), the more work is done;
\item pulling \emph{horizontally} is more effective than pulling steeply upwards: when $\theta$ is large, much of the effort is wasted lifting the rope rather than moving the cart.
\end{itemize}

The quantity that combines these three observations is precisely $|\vec{F}||\vec{d}|\cos\theta$ — the scalar product of force and displacement.

\begin{definition}[Work done by a constant force]
Let a \cle{constant force} $\vec{F}$ act on a body which undergoes a straight-line displacement $\vec{d}$, and let $\theta$ be the angle between $\vec{F}$ and $\vec{d}$. The \cle{work done} by the force is
\[
W \;=\; \vec{F}\cdot\vec{d} \;=\; |\vec{F}|\,|\vec{d}|\cos\theta.
\]
Work is a \emph{scalar} quantity. Its SI unit is the \cle{joule} (J): $1\ \mathrm{J} = 1\ \mathrm{N\,m}$.
\end{definition}

\begin{popfigure}
\begin{tikzpicture}[scale=1.0, >=stealth]
% ground
\draw[thick, popmuted] (-0.5,0) -- (8.5,0);
\foreach \x in {0,0.5,...,8} \draw[popmuted] (\x,0) -- (\x-0.25,-0.18);
% cart
\draw[fill=popblueL, draw=popdark] (1.2,0.25) rectangle (2.6,1.05);
\draw[fill=white, draw=popdark] (1.55,0.25) circle (0.22);
\draw[fill=white, draw=popdark] (2.25,0.25) circle (0.22);
% rope and force
\draw[very thick, poporange, ->] (2.6,0.9) -- (4.6,2.3) node[above left, poporange] {$\vec{F}$};
\draw[dashed, popmuted] (2.6,0.9) -- (5.1,0.9);
\draw[popdark] (3.9,0.9) arc (0:34:1.3) node[midway, right, xshift=2pt] {$\theta$};
% displacement
\draw[very thick, popblue, ->] (1.9,-0.7) -- (6.9,-0.7) node[midway, below, popblue] {displacement $\vec{d}$};
% component
\draw[dashed, popgreen, ->] (2.6,0.9) -- (4.24,0.9) node[midway, above, popgreen, yshift=1pt] {$|\vec{F}|\cos\theta$};
\end{tikzpicture}
\legende{A porter pulls a cart with a force $\vec{F}$ at angle $\theta$ to the horizontal. Only the component $|\vec{F}|\cos\theta$ along the displacement does work.}
\end{popfigure}

Since $\vec{F}\cdot\vec{d} = \big(|\vec{F}|\cos\theta\big)|\vec{d}|$, the work is the product of the \emph{component of the force in the direction of motion} (Section 1!) and the distance moved. Equivalently, $W = |\vec{F}|\big(|\vec{d}|\cos\theta\big)$: the force times the component of the displacement along the force. Both readings are useful.

\courssub{Sign of the work: positive, negative and zero work}

Because $W = |\vec{F}||\vec{d}|\cos\theta$, the sign of $W$ is the sign of $\cos\theta$:

\begin{center}
\begin{tabularx}{\linewidth}{|l|X|X|}
\hline
\rowcolor{popblueL} Angle $\theta$ & Sign of $W$ & Physical meaning \\
\hline
$0^{\circ} \le \theta < 90^{\circ}$ & $W > 0$ & The force \emph{helps} the motion (driving force). \\
\hline
$\theta = 90^{\circ}$ & $W = 0$ & The force neither helps nor hinders (e.g.\ normal reaction on a horizontal floor). \\
\hline
$90^{\circ} < \theta \le 180^{\circ}$ & $W < 0$ & The force \emph{opposes} the motion (e.g.\ friction, braking force). \\
\hline
\end{tabularx}
\end{center}

When a force does negative work, we often quote the \emph{work done \textbf{against} that force}, which is the corresponding positive number. For example, if friction does work $-150\ \mathrm{J}$ on a sliding crate, we say that $150\ \mathrm{J}$ of work is done against friction.

\begin{attention}[By a force versus against a force]
Work done \textbf{by} friction is always negative (friction opposes motion); work done \textbf{against} friction is quoted as a positive number. Read the GCE question carefully to see which one is required, and give the sign accordingly.
\end{attention}

\courssub{Computing work from components}

When the force and the displacement are given in component form, the dot product formula $W = \vec{F}\cdot\vec{d}$ is by far the fastest method — no angle is needed.

\begin{methode}[Work from components]
\begin{enumerate}
\item Write the force as $\vec{F} = F_1\vec{i} + F_2\vec{j}$ (or $F_1\vec{i}+F_2\vec{j}+F_3\vec{k}$ in 3D).
\item Form the displacement vector $\vec{d} = \overrightarrow{AB} = \vec{b} - \vec{a}$, where $\vec{a}, \vec{b}$ are the position vectors of the start and end points $A$ and $B$.
\item Compute $W = \vec{F}\cdot\vec{d} = F_1d_1 + F_2d_2\ (+F_3d_3)$.
\item Attach the unit (joules) and interpret the sign.
\end{enumerate}
\end{methode}

\begin{attention}[Displacement, not distance]
Work is defined using the \emph{displacement vector} $\vec{d} = \overrightarrow{AB}$. If the path is not straight, do not use the total distance travelled. For a constant force, only the straight-line displacement from start to finish matters.
\end{attention}

\begin{exoresolu}[Work done by a force in component form]
A force $\vec{F} = (3\vec{i} + 5\vec{j})\ \mathrm{N}$ moves a particle from the point $A(2,1)$ to the point $B(6,4)$, where coordinates are in metres. Find the work done by $\vec{F}$.
\tcblower
\textbf{Solution.} The displacement is
\[
\vec{d} = \overrightarrow{AB} = (6-2)\vec{i} + (4-1)\vec{j} = 4\vec{i} + 3\vec{j}\ \ (\mathrm{m}).
\]
Hence
\[
W = \vec{F}\cdot\vec{d} = (3)(4) + (5)(3) = 12 + 15 = 27\ \mathrm{J}.
\]
The work is positive: overall the force assists the motion.
\end{exoresolu}

\begin{exoresolu}[Porter pulling a cart on horizontal ground]
A porter pulls a cart along horizontal ground with a force of $80\ \mathrm{N}$ by means of a rope inclined at $30^{\circ}$ to the horizontal. Find the work done by the pull when the cart moves $25\ \mathrm{m}$ in a straight line.
\tcblower
\textbf{Solution.} Here $|\vec{F}| = 80\ \mathrm{N}$, $|\vec{d}| = 25\ \mathrm{m}$ and $\theta = 30^{\circ}$:
\[
W = |\vec{F}||\vec{d}|\cos\theta = 80 \times 25 \times \cos 30^{\circ} = 2000 \times \frac{\sqrt{3}}{2} = 1000\sqrt{3} \approx 1732\ \mathrm{J}.
\]
Note that the vertical component $80\sin 30^{\circ}$ does no work, since the cart does not move vertically.
\end{exoresolu}

\courssub{Work done by gravity}

Let a body of mass $m$ move through a \emph{vertical} displacement $h$ (taking $h > 0$ upwards). The weight $\vec{W}_g = m\vec{g}$ acts vertically downwards with magnitude $mg$.

\begin{itemize}
\item If the body moves \textbf{up} through height $h$: the angle between weight and displacement is $180^{\circ}$, so $W_{\text{gravity}} = -mgh$. We say that work $mgh$ is done \emph{against gravity}.
\item If the body moves \textbf{down} through height $h$: the angle is $0^{\circ}$, so $W_{\text{gravity}} = +mgh$.
\end{itemize}

A remarkable fact, consistent with the previous warning: since gravity is constant and vertical, the work done by gravity depends only on the \emph{vertical} change in height, not on the path. A climber who ascends $10\ \mathrm{m}$ vertically along a winding mountain path has $mgh$ of work done against gravity, exactly as if he had been lifted straight up.

\courssub{Total work of several forces}

\begin{propriete}[Work of a resultant]
Work is a scalar, and the scalar product is distributive over vector addition. Consequently, if forces $\vec{F}_1, \vec{F}_2, \ldots, \vec{F}_n$ all act on a body during the same displacement $\vec{d}$, then the total work is the sum of the individual works, and equals the work of the resultant:
\[
W_{\text{total}} = \sum_{i=1}^{n} \vec{F}_i \cdot \vec{d} = \Big(\sum_{i=1}^{n} \vec{F}_i\Big)\cdot\vec{d} = \vec{R}\cdot\vec{d}.
\]
\end{propriete}

This is extremely useful on an inclined plane, where applied force, weight, friction and normal reaction all act at once. The normal reaction, being perpendicular to the plane, does \emph{no} work.

\begin{popfigure}
\begin{tikzpicture}[scale=1.05, >=stealth]
% inclined plane
\draw[thick, popmuted] (0,0) -- (7.5,0);
\draw[thick, popdark] (0,0) -- (7,3.5);
\draw[popdark] (1.6,0) arc (0:26.565:1.6) node[midway, right, xshift=3pt] {$\alpha$};
% block (rotated frame)
\begin{scope}[rotate around={26.565:(3.2,1.6)}]
\draw[fill=popblueL, draw=popdark] (2.7,1.6) rectangle (3.7,2.4);
\end{scope}
% forces
\draw[very thick, poporange, ->] (3.2,2.35) -- (5.6,3.55) node[above, poporange] {$\vec{F}$ (applied)};
\draw[very thick, popgreen, ->] (3.2,2.0) -- (3.2,0.35) node[below, popgreen] {$m\vec{g}$};
\draw[very thick, poppink, ->] (2.95,2.1) -- (1.35,1.3) node[below left, poppink] {$\vec{f}$ (friction)};
\draw[very thick, popblue, ->] (3.45,2.45) -- (4.35,4.15) node[above right, popblue] {$\vec{R}$};
% displacement
\draw[very thick, popdark, ->] (4.6,1.55) -- (6.4,2.45) node[midway, below right, popdark] {$\vec{d}$};
\end{tikzpicture}
\legende{A block pulled up a rough inclined plane: applied force $\vec{F}$, weight $m\vec{g}$, friction $\vec{f}$ (down the slope), normal reaction $\vec{R}$ (no work), displacement $\vec{d}$ up the slope.}
\end{popfigure}

\begin{exoresolu}[Total work on a rough inclined plane]
A crate of mass $20\ \mathrm{kg}$ is pulled $8\ \mathrm{m}$ up a plane inclined at $25^{\circ}$ to the horizontal by a force of $150\ \mathrm{N}$ parallel to the plane. The frictional resistance is $40\ \mathrm{N}$. Take $g = 9.8\ \mathrm{m\,s^{-2}}$. Find (a) the work done by the applied force, (b) the work done by friction, (c) the work done by gravity, (d) the total work done on the crate.
\tcblower
\textbf{Solution.} Take the displacement $\vec{d}$ of magnitude $8\ \mathrm{m}$ up the slope.

\textbf{(a)} The applied force is parallel to $\vec{d}$:
\[
W_F = 150 \times 8 \times \cos 0^{\circ} = 1200\ \mathrm{J}.
\]

\textbf{(b)} Friction acts down the slope, opposite to $\vec{d}$:
\[
W_f = 40 \times 8 \times \cos 180^{\circ} = -320\ \mathrm{J}.
\]

\textbf{(c)} The crate rises vertically through $h = 8\sin 25^{\circ}\ \mathrm{m}$. Gravity acts downwards, so
\[
W_g = -mgh = -20 \times 9.8 \times 8\sin 25^{\circ} \approx -662.7\ \mathrm{J}.
\]
(Equivalently, the angle between the weight and $\vec{d}$ is $90^{\circ}+25^{\circ} = 115^{\circ}$, and $W_g = 196 \times 8 \times \cos 115^{\circ}$, giving the same result.)

\textbf{(d)} The normal reaction does no work, so
\[
W_{\text{total}} = 1200 - 320 - 662.7 = 217.3\ \mathrm{J}.
\]
The positive total work means the crate speeds up as it moves up the plane.
\end{exoresolu}

\courssub{Link with energy (exam-useful extra)}

\begin{savaistu}[The work--energy principle]
Although a formal proof belongs to mechanics, it is very useful at GCE to know the \cle{work--energy principle}: the total work done by all forces acting on a body equals the change in its \emph{kinetic energy}:
\[
W_{\text{total}} = \Delta\!\left(\tfrac{1}{2}mv^{2}\right) = \tfrac{1}{2}mv^{2} - \tfrac{1}{2}mu^{2}.
\]
In the previous example, the crate gained $\tfrac{1}{2}(20)(v^2-u^2) = 217.3\ \mathrm{J}$ of kinetic energy over the $8\ \mathrm{m}$. This gives a fast way of finding speeds without resolving forces and using kinematics.
\end{savaistu}

\begin{aretenir}[Key points]
\begin{itemize}
\item Work done by a constant force: $W = \vec{F}\cdot\vec{d} = |\vec{F}||\vec{d}|\cos\theta$, a scalar measured in joules.
\item $\theta < 90^{\circ}$: positive work; $\theta = 90^{\circ}$: zero work; $\theta > 90^{\circ}$: negative work.
\item With components, form $\vec{d} = \overrightarrow{AB}$ and take the dot product — no angle needed.
\item Work done by gravity: $\pm mgh$, sign depending on whether the body falls or rises; only the vertical height matters.
\item Total work = sum of individual works = work of the resultant. The normal reaction to a surface does no work.
\item Work done \emph{against} a resisting force is quoted as a positive number.
\end{itemize}
\end{aretenir}

\begin{exercice}[Practice]
\begin{enumerate}
\item A force $\vec{F} = (2\vec{i} - 3\vec{j} + \vec{k})\ \mathrm{N}$ acts on a particle that moves from $A(1,0,2)$ to $B(4,2,5)$, distances in metres. Find the work done.
\item A girl drags a box $12\ \mathrm{m}$ across a horizontal floor with a strap inclined at $40^{\circ}$ to the horizontal, exerting a force of $60\ \mathrm{N}$. Find the work done by the pull.
\item A load of mass $50\ \mathrm{kg}$ is raised vertically through $3\ \mathrm{m}$ at constant speed. Find the work done by the lifting force and the work done by gravity. ($g = 9.8\ \mathrm{m\,s^{-2}}$.)
\item A trolley of mass $10\ \mathrm{kg}$ is pushed $5\ \mathrm{m}$ up a smooth slope inclined at $30^{\circ}$ by a horizontal force. Using the work--energy principle, find the least work the pusher must supply if the trolley starts and ends at rest.
\end{enumerate}
\end{exercice}

\begin{corrige}[Exercise 1]
\textbf{1.} $\vec{d} = \overrightarrow{AB} = 3\vec{i} + 2\vec{j} + 3\vec{k}$, so $W = (2)(3) + (-3)(2) + (1)(3) = 6 - 6 + 3 = 3\ \mathrm{J}$.

\textbf{2.} $W = 60 \times 12 \times \cos 40^{\circ} \approx 720 \times 0.766 = 551.5\ \mathrm{J}$.

\textbf{3.} Constant speed means the lifting force equals the weight, $T = mg = 490\ \mathrm{N}$. Work done by the lifting force: $W_T = 490 \times 3 = 1470\ \mathrm{J}$. Work done by gravity: $W_g = -mgh = -1470\ \mathrm{J}$ (total work zero, consistent with constant speed).

\textbf{4.} Starting and ending at rest, $\Delta KE = 0$, so the pusher's work must exactly balance the work of gravity: $W = mgh = 10 \times 9.8 \times 5\sin 30^{\circ} = 245\ \mathrm{J}$.
\end{corrige}

\coursec{Section 3: Moment of a Force and Equilibrium (Lesson 104)}

In the previous section we saw that the scalar product measures the \emph{work} done by a force: it tells us how effectively a force moves a body along a straight line. But forces do not only translate bodies --- they also \emph{turn} them. When you push a door, tighten a nut with a spanner, or balance on a see-saw at the playground, you are producing a \cle{turning effect}. Measuring that turning effect is the business of the \emph{vector product}, and it leads us to one of the most powerful ideas in mechanics: the \cle{moment of a force}.

\courssub{The Turning Effect of a Force: Intuition First}

\begin{experience}[Three everyday observations]
\begin{itemize}
\item \textbf{The door.} Push a door near its hinges: it is very hard to open. Push it at the handle, far from the hinges: it opens easily. The same force has a greater turning effect when applied \emph{farther} from the pivot.
\item \textbf{The spanner.} A mechanic in Douala loosening a stubborn wheel nut uses a \emph{long} spanner, and pulls \emph{perpendicularly} to the handle. Pulling along the handle (towards the nut) produces no rotation at all.
\item \textbf{The see-saw.} A small child can balance a heavy adult on a see-saw, provided the child sits far from the pivot and the adult sits close to it.
\end{itemize}
\end{experience}

These observations tell us that the turning effect of a force about a point $O$ depends on \textbf{two} things: the \emph{magnitude} of the force, and the \emph{perpendicular distance} from $O$ to the line along which the force acts. This motivates our definition.

\begin{definition}[Moment of a force about a point (2D)]
Let a force $\vec{F}$ act in a plane, and let $O$ be a point of that plane. The \cle{moment} of $\vec{F}$ about $O$ is
\[ M_O = F \times d, \]
where $d$ is the \emph{perpendicular distance} from $O$ to the \cle{line of action} of $\vec{F}$. The moment is a signed quantity: we attach $+$ to one sense of rotation (usually anticlockwise) and $-$ to the other. The SI unit of moment is the newton-metre, \si{N.m}.
\end{definition}

\begin{attention}[Use the perpendicular distance!]
The distance $d$ is \emph{not} necessarily the distance from $O$ to the point of application of the force. It is the shortest (perpendicular) distance from $O$ to the \emph{line of action}. In particular, the moment of a force about any point \textbf{on its own line of action} is zero, because $d = 0$. This is why pushing a door directly towards its hinges never opens it.
\end{attention}

\begin{popfigure}
\begin{center}
\begin{tikzpicture}[scale=1.0, >=stealth]
% nut
\fill[gray!60] (0,0) circle (0.28);
\draw[thick, popdark] (0,0) circle (0.28);
\draw[thick, popdark] (0,0) circle (0.12);
\node[below] at (0,-0.35) {\small nut $O$};
% handle
\draw[line width=3pt, popblue] (0.28,0.12) -- (4.2,1.8);
\draw[line width=3pt, popblue] (0.28,-0.12) -- (4.2,1.56);
\draw[line width=3pt, popblue] (4.2,1.56) -- (4.2,1.8);
% force at end, perpendicular to handle
\draw[->, very thick, poporange] (4.2,1.68) -- ++(-0.62,1.44) node[above left] {$\vec{F}$};
% r along handle
\draw[<->, popgreen, thick] (0.15,-0.35) -- (4.05,1.28) node[midway, below, sloped] {$\vec{r}$};
% perpendicular distance dashed
\draw[dashed, poppink, thick] (0,0) -- (3.35,2.75);
\draw[popink] (3.35,2.75) -- ++(0.18,0.075) -- ++(0.075,-0.18);
\node[poppink, right] at (2.6,2.2) {$d = r\sin\theta$};
% angle theta
\draw[popdark] (1.3,0.52) arc (22.8:62:1.35);
\node[popdark] at (1.75,1.05) {$\theta$};
\end{tikzpicture}
\end{center}
\legende{A spanner turning a nut: the moment about $O$ is $M = Fd = Fr\sin\theta$.}
\end{popfigure}

\courssub{The Vector Definition in Three Dimensions: $\vec{M} = \vec{r} \times \vec{F}$}

In three dimensions the turning effect has not only a magnitude but also an \emph{axis} of rotation. Both are captured elegantly by the vector product.

\begin{definition}[Moment of a force about a point (3D)]
Let $\vec{F}$ be a force whose line of action passes through a point $A$, and let $\vec{r} = \overrightarrow{OA}$ be the position vector of $A$ relative to $O$. The \cle{moment} of $\vec{F}$ about $O$ is the vector
\[ \vec{M}_O = \vec{r} \times \vec{F}. \]
The result is unchanged if $A$ is replaced by \emph{any} other point of the line of action.
\end{definition}

Why is the choice of point on the line of action irrelevant? If $A'$ is another point of the line of action, then $\overrightarrow{OA'} = \vec{r} + \lambda \vec{F}$ for some scalar $\lambda$ (since $\overrightarrow{AA'}$ is parallel to $\vec{F}$). Hence
\[ \overrightarrow{OA'} \times \vec{F} = \vec{r} \times \vec{F} + \lambda\,(\vec{F} \times \vec{F}) = \vec{r} \times \vec{F}, \]
because the cross product of parallel vectors is zero.

\begin{propriete}[Magnitude and direction of the moment]
The magnitude of the moment is
\[ |\vec{M}_O| = |\vec{r}|\,|\vec{F}|\sin\theta = F\,d, \]
where $\theta$ is the angle between $\vec{r}$ and $\vec{F}$, and $d = r\sin\theta$ is the perpendicular distance from $O$ to the line of action. The direction of $\vec{M}_O$ is perpendicular to the plane containing $\vec{r}$ and $\vec{F}$, given by the \cle{right-hand rule}: curl the fingers of the right hand from $\vec{r}$ towards $\vec{F}$; the thumb points along $\vec{M}_O$, i.e.\ along the axis about which the force tends to rotate the body.
\end{propriete}

\begin{popfigure}
\begin{center}
\begin{tikzpicture}[scale=1.0, >=stealth]
% axes
\draw[->, gray] (0,0,0) -- (4.5,0,0) node[right] {$x$};
\draw[->, gray] (0,0,0) -- (0,3.2,0) node[above] {$y$};
\draw[->, gray] (0,0,0) -- (0,0,3.5) node[below left] {$z$};
% r vector
\draw[->, very thick, popblue] (0,0,0) -- (3,0,1.6) node[midway, below right] {$\vec{r}$};
% F vector
\draw[->, very thick, poporange] (3,0,1.6) -- (3,2.2,1.6) node[right] {$\vec{F}$};
% moment vector
\draw[->, very thick, poppurple] (0,0,0) -- (-0.9,1.8,1.35) node[above left] {$\vec{M} = \vec{r}\times\vec{F}$};
% right hand rule arc
\draw[->, popgreen, thick] (3.6,0.25,1.9) arc (40:130:0.55);
\node[popgreen] at (4.15,0.95) {\small right-hand rule};
% dashed helper
\draw[dashed, popmuted] (3,0,1.6) -- (3,2.2,1.6);
\node[popdark] at (0,0,0) [below left] {$O$};
\end{tikzpicture}
\end{center}
\legende{In 3D, $\vec{M} = \vec{r} \times \vec{F}$ is perpendicular to the plane of $\vec{r}$ and $\vec{F}$; its sense is given by the right-hand rule.}
\end{popfigure}

\begin{exemplebox}[Computing a moment in 3D]
A force $\vec{F} = (2\vec{i} + 3\vec{j} - \vec{k})\ \mathrm{N}$ acts at the point $A(1, 2, 0)$. Find its moment about the origin $O$.
\end{exemplebox}

We have $\vec{r} = \overrightarrow{OA} = \vec{i} + 2\vec{j}$, so
\[
\vec{M}_O = \vec{r} \times \vec{F} =
\begin{pmatrix} 1 \\ 2 \\ 0 \end{pmatrix} \times
\begin{pmatrix} 2 \\ 3 \\ -1 \end{pmatrix}
=
\begin{pmatrix} (2)(-1) - (0)(3) \\ (0)(2) - (1)(-1) \\ (1)(3) - (2)(2) \end{pmatrix}
=
\begin{pmatrix} -2 \\ 1 \\ -1 \end{pmatrix}\ \si{N.m}.
\]
The magnitude is $|\vec{M}_O| = \sqrt{4 + 1 + 1} = \sqrt{6}\ \si{N.m}$.

\courssub{Varignon's Theorem}

\begin{propriete}[Varignon's theorem]
The moment about any point $O$ of the resultant $\vec{R}$ of a system of forces $\vec{F}_1, \vec{F}_2, \dots, \vec{F}_n$ is equal to the sum of the moments of the individual forces about the same point $O$:
\[ \vec{M}_O(\vec{R}) = \sum_{i=1}^n \vec{M}_O(\vec{F}_i). \]
For a set of \emph{concurrent} forces (all passing through a point $A$ with position vector $\vec{r}$ relative to $O$), this reduces to $\vec{r} \times \vec{R} = \sum \vec{r} \times \vec{F}_i$, which follows immediately from the distributivity of the cross product over addition. A particularly useful corollary is that a single force may be replaced by its components acting at the same point when computing moments --- often the fastest technique in 2D problems.
\end{propriete}

\begin{exemplebox}[Using components (Varignon) in 2D]
A force of \SI{10}{N} acts at the point $A(3, 4)$ in the direction of the vector $(1, 1)$. Find its moment about $O$.
\end{exemplebox}

The unit vector along $(1,1)$ is $\tfrac{1}{\sqrt{2}}(1,1)$, so $\vec{F} = \tfrac{10}{\sqrt{2}}(\vec{i} + \vec{j}) = 5\sqrt{2}\,\vec{i} + 5\sqrt{2}\,\vec{j}$. Taking moments of the \emph{components} about $O$ (with anticlockwise positive):
\[ M_O = x F_y - y F_x = (3)(5\sqrt{2}) - (4)(5\sqrt{2}) = -5\sqrt{2}\ \si{N.m}. \]
The negative sign means the turning effect is \emph{clockwise}, of magnitude \SI{5\sqrt{2}}{N.m}.

\courssub{Equilibrium of a Rigid Body and the Principle of Moments}

A rigid body is in \cle{equilibrium} when it neither translates nor rotates. Both possibilities must be blocked:

\begin{aretenir}[Conditions for equilibrium of a rigid body]
A rigid body is in equilibrium if and only if
\begin{enumerate}
\item the \textbf{resultant force} is zero: $\sum \vec{F} = \vec{0}$ (no translation);
\item the \textbf{resultant moment} about \emph{any} point is zero: $\sum \vec{M}_O = \vec{0}$ (no rotation).
\end{enumerate}
In 2D problems, condition 1 gives two scalar equations (resolving horizontally and vertically) and condition 2 gives one scalar equation (moments).
\end{aretenir}

\begin{propriete}[Principle of moments]
When a body is in equilibrium, the sum of the \emph{clockwise} moments about any point equals the sum of the \emph{anticlockwise} moments about that same point:
\[ \sum M_{\text{clockwise}} = \sum M_{\text{anticlockwise}}. \]
This is simply condition 2 restated, and it is the physical law behind every balance, lever and scale.
\end{propriete}

\begin{attention}[Fix a sign convention before you start]
The most frequent source of error in GCE scripts is mixing signs: a candidate counts one clockwise moment as positive and another as negative in the same equation. \textbf{Before writing any moment equation}, decide: ``anticlockwise is positive'' (or the reverse), and apply it to \emph{every} term. Alternatively, collect all clockwise moments on one side and all anticlockwise moments on the other.
\end{attention}

\begin{methode}[Solving beam and support problems]
\begin{enumerate}
\item Draw a clear diagram of the beam; mark \emph{all} forces: loads, weights (acting at the centre of gravity of a uniform beam), and the unknown reactions at the supports.
\item Resolve forces vertically: $\sum F_y = 0$ (and horizontally if needed). This gives one equation linking the reactions.
\item Take moments about \emph{one of the supports}: the reaction at that support then has zero moment (its line of action passes through the point), so it disappears from the equation, leaving an equation in the \emph{other} reaction only.
\item Solve, then substitute back into the force equation to find the remaining reaction.
\item Check: take moments about the \emph{other} support and verify consistency.
\end{enumerate}
\end{methode}

\begin{popfigure}
\begin{center}
\begin{tikzpicture}[scale=1.0, >=stealth]
% beam
\draw[line width=3pt, popdark] (0,0) -- (8,0);
% supports
\draw[popblue, thick] (0.5,0) -- (0.2,-0.5) -- (0.8,-0.5) -- cycle;
\draw[popblue, thick] (6.5,0) -- (6.2,-0.5) -- (6.8,-0.5) -- cycle;
\node[below] at (0.5,-0.55) {\small $A$};
\node[below] at (6.5,-0.55) {\small $B$};
% reactions
\draw[->, very thick, popgreen] (0.5,-0.5) -- (0.5,0.9) node[above] {$R_A$};
\draw[->, very thick, popgreen] (6.5,-0.5) -- (6.5,0.9) node[above] {$R_B$};
% loads
\draw[->, very thick, poporange] (3,1.1) -- (3,0.05) node[above, yshift=32] {\SI{40}{N}};
\draw[->, very thick, poporange] (5.5,1.1) -- (5.5,0.05) node[above, yshift=32] {\SI{60}{N}};
% distances
\draw[<->, popmuted] (0.5,-1.1) -- (3,-1.1) node[midway, below] {\small \SI{2.5}{m}};
\draw[<->, popmuted] (3,-1.1) -- (5.5,-1.1) node[midway, below] {\small \SI{2.5}{m}};
\draw[<->, popmuted] (5.5,-1.1) -- (6.5,-1.1) node[midway, below] {\small \SI{1}{m}};
\draw[<->, popmuted] (6.5,-1.1) -- (8,-1.1) node[midway, below] {\small \SI{1.5}{m}};
\end{tikzpicture}
\end{center}
\legende{A beam on two supports $A$ and $B$ carrying two point loads. Take moments about $A$ to find $R_B$ directly.}
\end{popfigure}

\begin{exoresolu}[Reactions at the supports of a beam]
A light beam $AC$ of length \SI{7.5}{m} rests horizontally on supports at $A$ and at $B$, where $AB = \SI{6}{m}$ and $B$ is \SI{1.5}{m} from the end $C$. A load of \SI{40}{N} hangs \SI{2.5}{m} from $A$, and a load of \SI{60}{N} hangs \SI{5}{m} from $A$. Find the reactions $R_A$ and $R_B$.
\tcblower
\textbf{Step 1 --- Resolve vertically.} The beam is light (negligible weight), so
\[ R_A + R_B = 40 + 60 = \SI{100}{N}. \tag{1} \]
\textbf{Step 2 --- Moments about $A$} (anticlockwise positive; $R_A$ has zero moment about $A$):
\[ R_B \times 6 = 40 \times 2.5 + 60 \times 5 = 100 + 300 = 400, \]
\[ R_B = \frac{400}{6} = \frac{200}{3} \approx \SI{66.7}{N}. \]
\textbf{Step 3 --- Back-substitute in (1):}
\[ R_A = 100 - \frac{200}{3} = \frac{100}{3} \approx \SI{33.3}{N}. \]
\textbf{Check} (moments about $B$): anticlockwise moment of $R_A$: $\tfrac{100}{3}\times 6 = 200$; clockwise moments: $40 \times 3.5 + 60 \times 1 = 140 + 60 = 200$. \checkmark\ The two agree, so the beam is in equilibrium with $R_A = \tfrac{100}{3}\ \si{N}$ and $R_B = \tfrac{200}{3}\ \si{N}$.
\end{exoresolu}

\begin{exoresolu}[The see-saw: principle of moments]
A uniform see-saw plank of length \SI{4}{m} and weight \SI{100}{N} is pivoted at its centre. A child of weight \SI{250}{N} sits at one end. Where must a second child of weight \SI{400}{N} sit so that the see-saw balances horizontally?
\tcblower
The plank is uniform, so its weight acts at the pivot and has \emph{zero} moment about it. Take moments about the pivot, anticlockwise positive, and let $x$ be the distance of the second child from the pivot on the opposite side:
\[ 400\,x = 250 \times 2 \quad\Longrightarrow\quad x = \frac{500}{400} = \SI{1.25}{m}. \]
The second child must sit \SI{1.25}{m} from the pivot, on the side opposite to the first child. Note how the heavier child sits \emph{closer} to the pivot --- exactly what everyday experience predicts.
\end{exoresolu}

\begin{savaistu}[Moments in daily Cameroonian life]
The balance scale used by traders in the markets of Bamenda and Bafoussam is a direct application of the principle of moments: the unknown mass is placed on one pan and known masses on the other, and balance (zero resultant moment about the pivot) guarantees that the weights are equal. The same principle lets a carpenter use a long crowbar to lift a heavy stone with a modest effort: a small force with a large lever arm produces the same moment as a large force with a small lever arm. Archimedes summarised it long ago: ``Give me a place to stand, and I shall move the Earth.''
\end{savaistu}

\begin{attention}[Classic traps in moment problems]
\begin{itemize}
\item Forgetting the weight of a \emph{uniform} beam: it acts at the \textbf{midpoint} of the beam, not at a support.
\item Using the distance from $O$ to the point of application instead of the \emph{perpendicular} distance to the line of action when the force is oblique. When in doubt, resolve the force into components and use Varignon's theorem.
\item Taking moments about a point and \emph{still} including the moment of a force whose line of action passes through that point --- its moment is zero.
\item In 3D, computing $\vec{F} \times \vec{r}$ instead of $\vec{r} \times \vec{F}$: the order matters, since $\vec{F} \times \vec{r} = -\,\vec{r} \times \vec{F}$.
\end{itemize}
\end{attention}

\begin{exercice}[Ladder-type problem]
A uniform ladder of length \SI{5}{m} and weight \SI{200}{N} rests with its top against a smooth vertical wall and its foot on rough horizontal ground, making an angle of $60^{\circ}$ with the ground. The reaction of the smooth wall is horizontal; the ground exerts a vertical reaction and a friction force. By resolving forces and taking moments about the foot of the ladder, show that the wall reaction is $\dfrac{100}{\sqrt{3}}\ \si{N} \approx \SI{57.7}{N}$, and find the friction force and the vertical ground reaction.
\end{exercice}

\begin{corrige}[Exercise 1]
Let $N$ be the horizontal wall reaction, $F$ the friction and $R$ the vertical ground reaction. Resolving vertically: $R = \SI{200}{N}$. Resolving horizontally: $F = N$. The weight acts at the midpoint, \SI{2.5}{m} along the ladder; its perpendicular distance to the foot is $2.5\cos 60^{\circ} = \SI{1.25}{m}$. The wall reaction acts at the top, at perpendicular distance $5\sin 60^{\circ} = \tfrac{5\sqrt{3}}{2}\ \si{m}$ from the foot. Moments about the foot (anticlockwise positive):
\[ N \times \frac{5\sqrt{3}}{2} = 200 \times 1.25 = 250
\quad\Longrightarrow\quad
N = \frac{500}{5\sqrt{3}} = \frac{100}{\sqrt{3}} \approx \SI{57.7}{N}. \]
Hence $F = \dfrac{100}{\sqrt{3}}\ \si{N}$ and $R = \SI{200}{N}$.
\end{corrige}

\begin{aretenir}[Key points of Lesson 104]
\begin{itemize}
\item Moment of a force about $O$: $M = Fd$ in 2D (perpendicular distance), $\vec{M} = \vec{r} \times \vec{F}$ in 3D; unit \si{N.m}.
\item $|\vec{M}| = rF\sin\theta$; direction given by the right-hand rule; sense clockwise or anticlockwise in 2D.
\item Varignon: the moment of the resultant equals the sum of the moments of the forces --- resolve into components to simplify.
\item Equilibrium requires \textbf{both} $\sum \vec{F} = \vec{0}$ \textbf{and} $\sum \vec{M}_O = \vec{0}$; in equilibrium, clockwise moments equal anticlockwise moments about any point.
\item Strategy: resolve to zero, then take moments about a support to eliminate one unknown --- and always fix a sign convention first.
\end{itemize}
\end{aretenir}

With the moment of a single force mastered, we are ready for the next step: what happens when several forces act together --- and in particular the remarkable pair of forces called a \cle{couple}, which produces a pure turning effect with no translation at all. That is the subject of Section 4.

\coursec{Section 4: Single Forces and Couples (Lesson 105)}

In Section 3 we studied the moment of a single force and saw that a body is in equilibrium when both the resultant force and the resultant moment vanish. But a natural question arises: what happens when the resultant force is zero \emph{yet the body still rotates}? Anyone who has turned a steering wheel with both hands, or opened a water tap, already knows the answer. Two equal and opposite pushes applied at different points produce no translation at all, but a very definite rotation. This special pair of forces is called a \cle{couple}, and it is the subject of this lesson. Understanding couples will then allow us to answer a deeper question: can \emph{any} system of forces be reduced to something simple? The answer — a single force, or a force plus a couple — is one of the most powerful results in statics.

\courssub{Definition of a Couple}

\begin{experience}[Turning the steering wheel]
Sit in a parked car (engine off, handbrake on — ask permission first!). Place your left hand at the top of the steering wheel and your right hand at the bottom. Push the top to the right and pull the bottom to the left with forces of equal magnitude. The centre of the wheel does not move: the two forces cancel as far as translation is concerned. Yet the wheel turns. The rotation is produced by the \emph{pair} of forces, not by either force alone. The same effect occurs when you open a tap with two fingers, when a mechanic uses two spanners on a stubborn nut in a Douala workshop, or when you turn a screwdriver: the thumb pushes one way, the fingers push the other way.
\end{experience}

\begin{definition}[Couple]
A \cle{couple} is a system of two forces $\vec{F}$ and $-\vec{F}$ which are:
\begin{itemize}
\item equal in magnitude,
\item opposite in direction,
\item parallel,
\item and have \textbf{different lines of action}.
\end{itemize}
The plane containing the two lines of action is called the \emph{plane of the couple}, and the perpendicular distance $d$ between the two lines of action is called the \emph{arm} of the couple.
\end{definition}

\begin{attention}[A couple is not a balanced system]
Two equal and opposite forces acting along the \emph{same} line (like two teams pulling a rope with equal strength) cancel completely: no translation, no rotation. A couple is different precisely because the lines of action are \emph{distinct}. Never confuse ``equal and opposite forces'' with ``a couple'': the separation $d \neq 0$ is essential.
\end{attention}

\courssub{Moment of a Couple}

Let the force $\vec{F}$ act at point $B$ and the force $-\vec{F}$ act at point $A$, and let $O$ be \emph{any} point. The total moment about $O$ is
\[
\vec{G}_O = \vec{OB} \wedge \vec{F} + \vec{OA} \wedge (-\vec{F}) = (\vec{OB} - \vec{OA}) \wedge \vec{F} = \vec{AB} \wedge \vec{F}.
\]
The point $O$ has disappeared from the result! This proves the fundamental property:

\begin{propriete}[Moment of a couple — proof examinable]
The moment of a couple is the same about \emph{every} point. It is a \cle{free vector}:
\[
\vec{G} = \vec{r} \wedge \vec{F},
\]
where $\vec{r}$ is the vector joining the point of application of $-\vec{F}$ to the point of application of $\vec{F}$. Its magnitude is
\[
G = F\,d,
\]
where $F$ is the magnitude of either force and $d$ is the perpendicular distance between the two lines of action (the arm). The SI unit is the newton metre, $\mathrm{N\,m}$.
\end{propriete}

\begin{methode}[Computing the moment of a couple]
\begin{enumerate}
\item Identify the two equal, opposite, parallel forces and their lines of action.
\item Measure (or compute) the \emph{perpendicular} distance $d$ between the two lines of action.
\item Multiply: $G = F \times d$.
\item Give the sense of rotation (clockwise or anticlockwise), or the vector direction by the right-hand rule.
\end{enumerate}
\end{methode}

\begin{attention}[Use the perpendicular distance]
A very common mistake is to multiply $F$ by the distance $AB$ between the points of application instead of the \emph{perpendicular} distance between the lines of action. If $\vec{AB}$ makes an angle $\theta$ with $\vec{F}$, then $d = AB \sin\theta$ and $G = F \cdot AB \sin\theta$ — exactly what the cross product $\vec{AB} \wedge \vec{F}$ gives automatically.
\end{attention}

\begin{exoresolu}[Moment of a couple on a steering wheel]
A driver applies two forces of magnitude $F = 15\ \mathrm{N}$ tangentially at opposite ends of a diameter of a steering wheel of radius $20\ \mathrm{cm}$. Find the moment of the couple.
\tcblower
\textbf{Solution.} The two forces are equal, opposite and parallel, with lines of action separated by the diameter:
\[
d = 2 \times 0.20 = 0.40\ \mathrm{m}.
\]
Hence
\[
G = F\,d = 15 \times 0.40 = 6\ \mathrm{N\,m}.
\]
The moment is $6\ \mathrm{N\,m}$ about \emph{any} point — the centre of the wheel, the rim, or even a point outside the car.
\end{exoresolu}

\begin{popfigure}
\begin{center}
\begin{tikzpicture}[scale=1.1]
\draw[thick, popblue] (0,0) circle (2);
\draw[thick, popblue] (0,0) circle (0.35);
\draw[thick, popblue] (-2,0)--(-0.35,0); \draw[thick, popblue] (0.35,0)--(2,0);
\draw[thick, popblue] (0,-0.35)--(0,-2); \draw[thick, popblue] (0,0.35)--(0,2);
\fill (0,0) circle (0.05) node[below right=1pt, text=popink] {$O$};
\draw[-{Stealth}, very thick, poporange] (2,0) -- (2,1.6) node[right, text=poporange] {$\vec{F}$};
\draw[-{Stealth}, very thick, poporange] (-2,0) -- (-2,-1.6) node[left, text=poporange] {$-\vec{F}$};
\draw[{Stealth}-{Stealth}, thick, poppurple] (2,-0.35) -- (-2,-0.35) node[midway, below, text=poppurple] {$d$};
\draw[-{Stealth}, thick, popgreen] (0.9,1.15) arc (52:128:1.45) node[midway, above, text=popgreen] {$G = Fd$};
\end{tikzpicture}
\end{center}
\legende{A couple applied to a steering wheel: two equal opposite tangential forces at the rim, arm $d$ equal to the diameter, moment $G = Fd$.}
\end{popfigure}

\courssub{A Couple Produces Pure Rotation}

The vector sum of the two forces of a couple is
\[
\vec{F} + (-\vec{F}) = \vec{0},
\]
so a couple has \textbf{zero resultant force}: it cannot accelerate the centre of mass of a body. But its resultant moment $G = Fd \neq 0$, so it produces \cle{pure rotation}. A couple is the purest form of ``turning effect'': it is to rotation what a single force is to translation.

\begin{attention}[A couple cannot be balanced by a single force]
Since the resultant force of a couple is zero, no single force can ever cancel its turning effect. A couple can only be balanced by \textbf{another couple} of equal magnitude and opposite sense. In equilibrium problems, if the applied loads include a couple of moment $G$, the supports must provide a resisting couple $-G$; a reaction force alone will not do.
\end{attention}

\courssub{Equivalence of Couples}

Because the moment of a couple depends only on the product $Fd$ (and the sense of rotation), two couples are \cle{equivalent} whenever they have the same moment vector, whatever the individual forces and arms.

\begin{propriete}[Equivalence of couples]
Two couples are equivalent if and only if their moment vectors are equal. Consequently a couple may be replaced by any other couple with the same moment: for example, forces of $10\ \mathrm{N}$ with arm $0.6\ \mathrm{m}$ are equivalent to forces of $30\ \mathrm{N}$ with arm $0.2\ \mathrm{m}$, since both give $G = 6\ \mathrm{N\,m}$.
\end{propriete}

\begin{exemplebox}[Two spanners on a nut]
A mechanic loosens a nut using two spanners, pushing with forces of $40\ \mathrm{N}$ at a separation of $25\ \mathrm{cm}$: $G = 40 \times 0.25 = 10\ \mathrm{N\,m}$. His apprentice, with shorter spanners separated by $12.5\ \mathrm{cm}$, must push with forces of $\dfrac{10}{0.125} = 80\ \mathrm{N}$ to produce the same turning effect. Same couple, different forces.
\end{exemplebox}

\courssub{Translating a Force: Force Plus a Couple}

Here is a beautiful trick. Let a force $\vec{F}$ act at a point $A$, and let $O$ be any other point. Add at $O$ two forces $\vec{F}$ and $-\vec{F}$ (their sum is zero, so nothing has changed). Now look at what we have:
\begin{itemize}
\item the original $\vec{F}$ at $A$ together with $-\vec{F}$ at $O$ form a \textbf{couple} of moment $\vec{OA} \wedge \vec{F}$;
\item and there remains a force $\vec{F}$ acting at $O$.
\end{itemize}

\begin{propriete}[Force translation theorem — proof examinable]
A force $\vec{F}$ acting at a point $A$ is equivalent to an equal parallel force $\vec{F}$ acting at any other point $O$, \emph{together with a couple} of moment
\[
\vec{M} = \vec{OA} \wedge \vec{F}, \qquad M = Fd,
\]
where $d$ is the perpendicular distance from $O$ to the original line of action.
\end{propriete}

\begin{popfigure}
\begin{center}
\begin{tikzpicture}[scale=1.0]
% Left: original force at A
\fill (0,0) circle (0.06) node[below left, text=popink] {$O$};
\fill (3,0) circle (0.06) node[below right, text=popink] {$A$};
\draw[-{Stealth}, very thick, poporange] (3,0) -- (3,2) node[right, text=poporange] {$\vec{F}$};
\draw[dashed, popmuted] (0,0) -- (3,0);
\draw[{Stealth}-{Stealth}, thick, poppurple] (0,-0.5) -- (3,-0.5) node[midway, below, text=poppurple] {$d$};
\node[text=popink] at (1.5,-1.3) {Force at $A$};
% Arrow of equivalence
\draw[-{Stealth}, very thick, popink] (4.3,0.6) -- (5.5,0.6) node[midway, above, text=popink] {$\equiv$};
% Right: force at O plus couple
\begin{scope}[xshift=6.2cm]
\fill (0,0) circle (0.06) node[below left, text=popink] {$O$};
\fill (3,0) circle (0.06) node[below right, text=popink] {$A$};
\draw[-{Stealth}, very thick, poporange] (0,0) -- (0,2) node[left, text=poporange] {$\vec{F}$};
\draw[dashed, popmuted] (0,0) -- (3,0);
\draw[-{Stealth}, thick, popgreen] (2.2,0.9) arc (55:125:1.15) node[midway, above, text=popgreen] {$M = Fd$};
\node[text=popink] at (1.5,-1.3) {Force at $O$ $+$ couple};
\end{scope}
\end{tikzpicture}
\end{center}
\legende{Force translation theorem: $\vec{F}$ at $A$ is equivalent to the same force at $O$ plus a couple of moment $M = Fd$.}
\end{popfigure}

\begin{exoresolu}[Translating a force]
A force $\vec{F} = 50\ \mathrm{N}$ acts vertically upwards at point $A$, where $\vec{OA} = 0.3\ \mathrm{m}$ horizontally to the right of $O$. Replace it by a force at $O$ plus a couple.
\tcblower
\textbf{Solution.} By the force translation theorem, we place a force of $50\ \mathrm{N}$ vertically upwards at $O$, together with a couple of moment
\[
M = Fd = 50 \times 0.3 = 15\ \mathrm{N\,m}.
\]
The sense of the couple is anticlockwise (the original force tends to rotate the body about $O$ anticlockwise). Check with the cross product: with $\vec{OA} = 0.3\,\vec{i}$ and $\vec{F} = 50\,\vec{j}$,
\[
\vec{M} = \vec{OA} \wedge \vec{F} = 0.3 \times 50\, (\vec{i} \wedge \vec{j}) = 15\,\vec{k}\ \mathrm{N\,m},
\]
confirming the anticlockwise sense about $O$.
\end{exoresolu}

\courssub{Reduction of a System of Forces at a Point}

The force translation theorem gives us a complete recipe for simplifying \emph{any} system of forces. Take a system of forces $\vec{F}_1, \vec{F}_2, \ldots, \vec{F}_n$ acting at points $A_1, A_2, \ldots, A_n$, and choose a point $O$ (the \emph{centre of reduction}). Translate each force to $O$: each $\vec{F}_i$ contributes a force $\vec{F}_i$ at $O$ plus a couple $\vec{OA}_i \wedge \vec{F}_i$. Adding everything:

\begin{propriete}[Reduction of a force system]
Any system of forces is equivalent, at an arbitrary point $O$, to
\begin{itemize}
\item a \textbf{single resultant force} $\vec{R} = \displaystyle\sum_{i} \vec{F}_i$ acting through $O$,
\item together with a \textbf{single resultant couple} $\vec{G}_O = \displaystyle\sum_{i} \vec{OA}_i \wedge \vec{F}_i$.
\end{itemize}
Special cases:
\begin{itemize}
\item If $\vec{R} \neq \vec{0}$ and $\vec{G}_O$ is perpendicular to $\vec{R}$ (in particular for coplanar forces), the system can be reduced further to a single force $\vec{R}$ alone, acting along a line at perpendicular distance $d = \dfrac{G_O}{R}$ from $O$.
\item If $\vec{R} = \vec{0}$ and $\vec{G}_O \neq \vec{0}$, the system reduces to a \textbf{couple}.
\item If $\vec{R} = \vec{0}$ and $\vec{G}_O = \vec{0}$, the system is in \textbf{equilibrium}.
\item In general three-dimensional systems, the reduction yields a \cle{wrench}: a force $\vec{R}$ together with a couple $\vec{G}_\parallel$ parallel to $\vec{R}$ (the component of $\vec{G}_O$ along $\vec{R}$ cannot be eliminated). This is studied in Lesson 107.
\end{itemize}
\end{propriete}

\begin{exoresolu}[Reducing a system of two parallel forces]
Two parallel forces act on a horizontal bar: $\vec{F}_1 = 20\ \mathrm{N}$ downwards at $A_1$ with $OA_1 = 0.2\ \mathrm{m}$, and $\vec{F}_2 = 50\ \mathrm{N}$ downwards at $A_2$ with $OA_2 = 0.8\ \mathrm{m}$. Reduce the system at $O$, then find the single force to which it is equivalent.
\tcblower
\textbf{Solution.} \emph{Step 1: resultant force.}
\[
\vec{R} = \vec{F}_1 + \vec{F}_2 = (20 + 50) = 70\ \mathrm{N} \text{ downwards, acting at } O.
\]
\emph{Step 2: resultant couple at } $O$ (taking anticlockwise as positive, moments of downward forces to the right of $O$ are clockwise):
\[
G_O = -(20 \times 0.2) - (50 \times 0.8) = -4 - 40 = -44\ \mathrm{N\,m},
\]
i.e. $44\ \mathrm{N\,m}$ clockwise. \emph{Step 3: single force.} Since $\vec{R} \neq \vec{0}$ and the forces are coplanar, the system reduces to a single force of $70\ \mathrm{N}$ downwards whose line of action is at distance
\[
d = \frac{G_O}{R} = \frac{44}{70} \approx 0.629\ \mathrm{m}
\]
to the right of $O$. Check: $70 \times 0.629 \approx 44\ \mathrm{N\,m}$ clockwise. $\checkmark$
\end{exoresolu}

\begin{exercice}[Couples and reduction]
\begin{enumerate}
\item A couple consists of two forces of $25\ \mathrm{N}$ with arm $0.36\ \mathrm{m}$. Find its moment. What forces would give an equivalent couple with arm $0.60\ \mathrm{m}$?
\item A force of $80\ \mathrm{N}$ acts horizontally at $A$, with $O$ at perpendicular distance $0.45\ \mathrm{m}$ from its line of action. Replace it by a force at $O$ plus a couple.
\item Two forces of $30\ \mathrm{N}$ downwards act at distances $0.1\ \mathrm{m}$ and $0.5\ \mathrm{m}$ from $O$ on a bar. Reduce the system to a single force and give the position of its line of action.
\end{enumerate}
\end{exercice}

\begin{corrige}[Exercise 1]
\begin{enumerate}
\item $G = Fd = 25 \times 0.36 = 9\ \mathrm{N\,m}$. For an equivalent couple with arm $0.60\ \mathrm{m}$: $F = \dfrac{9}{0.60} = 15\ \mathrm{N}$.
\item Force of $80\ \mathrm{N}$ at $O$ (same direction) plus a couple $M = Fd = 80 \times 0.45 = 36\ \mathrm{N\,m}$, in the sense in which the original force turns about $O$.
\item $\vec{R} = 60\ \mathrm{N}$ downwards; $G_O = 30 \times 0.1 + 30 \times 0.5 = 18\ \mathrm{N\,m}$ clockwise. Single force of $60\ \mathrm{N}$ at distance $d = \dfrac{18}{60} = 0.3\ \mathrm{m}$ from $O$ (the midpoint, as symmetry suggests).
\end{enumerate}
\end{corrige}

\begin{aretenir}[Key points of Lesson 105]
\begin{itemize}
\item A \cle{couple} is a pair of equal, opposite, parallel forces with distinct lines of action.
\item Its moment $\vec{G} = \vec{r} \wedge \vec{F}$, of magnitude $G = Fd$, is a \cle{free vector}: the same about every point.
\item A couple has zero resultant force: it produces \cle{pure rotation} and can only be balanced by another couple.
\item Couples with the same moment are \cle{equivalent}.
\item A force at $A$ is equivalent to the same force at $O$ plus a couple $\vec{OA} \wedge \vec{F}$ (\cle{force translation theorem}).
\item Any force system reduces, at a point $O$, to a resultant force $\vec{R}$ plus a resultant couple $\vec{G}_O$; if $\vec{R} \neq \vec{0}$ and $\vec{G}_O \perp \vec{R}$ (coplanar case) it reduces further to a single force at distance $G_O/R$ from $O$.
\end{itemize}
\end{aretenir}

\begin{savaistu}[Couples all around us]
Every rotating machine obeys the physics of couples. The electric motors that power the fans in our homes deliver a couple (called \emph{torque}) to the shaft; the pedals of a bicycle work best when the rider pushes with two feet in opposition, forming a couple on the crank; and when carpenters in Bamenda tighten a vice, they instinctively push the two ends of the handle in opposite directions — applying a couple rather than a single off-centre force, which would bend the screw. Engineers design steering systems, taps and door handles precisely so that our two hands naturally form couples.
\end{savaistu}

\coursec{Section 5: Coplanar Forces and Simple 3D Systems (Lessons 106 \& 107)}

In Section 4 we saw that a system of forces acting on a rigid body can often be reduced either to a \cle{single force} (its resultant) or to a \cle{couple} (a pure turning effect with zero resultant force). We now put everything together. In this final section we answer three big questions that the GCE examiner loves: \emph{What is the resultant of a whole system of coplanar forces, and where does it act? When is a body in equilibrium under such a system? And what changes when the forces live in three dimensions, as with a particle hanging from cables?} Mastering this section is mastering the whole chapter.

\courssub{Resultant of a System of Coplanar Forces}

\textbf{Idea.} Suppose several forces $\vec{F}_1, \vec{F}_2, \dots, \vec{F}_n$ act on a rigid body in the same plane (the $xOy$ plane). The \cle{resultant} $\vec{R}$ is simply their vector sum, computed component by component:
\[
\vec{R} = \sum_{i=1}^{n} \vec{F}_i, \qquad R_x = \sum F_{ix}, \qquad R_y = \sum F_{iy}.
\]
Its magnitude and direction follow at once:
\[
|\vec{R}| = \sqrt{R_x^2 + R_y^2}, \qquad \tan\theta = \frac{R_y}{R_x},
\]
where $\theta$ is the angle the resultant makes with the positive $x$-axis (take care to place $\theta$ in the correct quadrant from the signs of $R_x$ and $R_y$).

\textbf{Where does the resultant act?} Knowing $\vec{R}$ is not enough: a force applied at different points of a rigid body produces different turning effects. The \cle{line of action} of the resultant is located by \emph{taking moments}. If the system is equivalent to the single force $\vec{R}$, then the moment of $\vec{R}$ about any point $O$ must equal the sum of the moments of the individual forces about $O$:
\[
\vec{R} \text{ acting along its line of action} \iff M_O(\vec{R}) = \sum_{i=1}^{n} M_O(\vec{F}_i).
\]
Recall from Section 3 that for a force $(F_x, F_y)$ applied at the point $(x, y)$, the moment about the origin is $M_O = xF_y - yF_x$. In practice, if the line of action of $\vec{R}$ cuts the $x$-axis at $(\bar{x}, 0)$, then $M_O(\vec{R}) = \bar{x}\, R_y$ (the component $R_x$ passes through the axis and gives no moment), so
\[
\bar{x} = \frac{\sum M_O(\vec{F}_i)}{R_y}.
\]
If $R_y = 0$ but $R_x \neq 0$, use instead the point $(0, \bar{y})$ where the line of action cuts the $y$-axis: $M_O(\vec{R}) = -\bar{y}\,R_x$, so $\bar{y} = -\dfrac{\sum M_O(\vec{F}_i)}{R_x}$.

\begin{methode}[Reducing a coplanar system to a single resultant]
\begin{enumerate}
\item Resolve every force into $x$ and $y$ components; sum to get $R_x$ and $R_y$.
\item Compute $|\vec{R}| = \sqrt{R_x^2+R_y^2}$ and $\theta = \tan^{-1}(R_y/R_x)$ in the correct quadrant.
\item Take moments of all the original forces about a convenient point $O$ (usually the origin), using $M_O = xF_y - yF_x$.
\item Locate the line of action from $\bar{x} = \dfrac{\sum M_O}{R_y}$ (or, if $R_y = 0$, use $\bar{y} = -\dfrac{\sum M_O}{R_x}$ on the $y$-axis).
\item If $\vec{R} = \vec{0}$ but $\sum M_O \neq 0$, the system reduces to a \textbf{couple} of moment $\sum M_O$.
\end{enumerate}
\end{methode}

\begin{exoresolu}[Resultant and its line of action]
Forces act on a lamina in the $xOy$ plane: $\vec{F}_1 = (3\vec{i} + 4\vec{j})\ \mathrm{N}$ at $A(0,0)$, $\vec{F}_2 = (-5\vec{i} + 2\vec{j})\ \mathrm{N}$ at $B(4,0)$, and $\vec{F}_3 = (4\vec{i} - 6\vec{j})\ \mathrm{N}$ at $C(2,3)$. Find the resultant and the equation of its line of action.
\tcblower
\textbf{Solution.} Components of the resultant:
\[
R_x = 3 - 5 + 4 = 2\ \mathrm{N}, \qquad R_y = 4 + 2 - 6 = 0\ \mathrm{N}.
\]
So $\vec{R} = 2\vec{i}\ \mathrm{N}$: a horizontal force of $2\ \mathrm{N}$. Since $R_y = 0$, we locate the line of action using the $y$-intercept instead. Moments about $O$ (using $M_O = xF_y - yF_x$):
\begin{align*}
M_O(\vec{F}_1) &= 0 \quad (\text{acts at } O),\\
M_O(\vec{F}_2) &= 4(2) - 0(-5) = 8\ \mathrm{N\,m},\\
M_O(\vec{F}_3) &= 2(-6) - 3(4) = -12 - 12 = -24\ \mathrm{N\,m}.
\end{align*}
Total: $\sum M_O = 8 - 24 = -16\ \mathrm{N\,m}$. If $\vec{R} = 2\vec{i}$ acts at $(0, \bar{y})$, its moment about $O$ is $-\bar{y} R_x = -2\bar{y}$. Hence $-2\bar{y} = -16$, giving $\bar{y} = 8\ \mathrm{m}$. The resultant is the force $2\vec{i}\ \mathrm{N}$ acting along the horizontal line $y = 8$.
\end{exoresolu}

\begin{popfigure}
\begin{center}
\begin{tikzpicture}[scale=0.9, >=stealth]
\draw[->, gray] (-0.5,0) -- (5.5,0) node[right] {$x$};
\draw[->, gray] (0,-0.5) -- (0,4.5) node[above] {$y$};
\draw[->, very thick, popblue] (0,0) -- (1.2,1.6) node[above left] {$\vec{F}_1$};
\fill[popblue] (0,0) circle (1.5pt) node[below left] {$A$};
\draw[->, very thick, poporange] (4,0) -- (2.5,0.8) node[above right] {$\vec{F}_2$};
\fill[poporange] (4,0) circle (1.5pt) node[below] {$B$};
\draw[->, very thick, popgreen] (2,3) -- (3.2,0.6) node[below right] {$\vec{F}_3$};
\fill[popgreen] (2,3) circle (1.5pt) node[above left] {$C$};
\draw[dashed, poppurple, thick] (-0.5,4) -- (5.5,4) node[right, poppurple] {line of action of $\vec{R}$};
\draw[->, very thick, poppurple] (1,4) -- (2.6,4) node[midway, above] {$\vec{R}=2\vec{i}$};
\end{tikzpicture}
\end{center}
\legende{A coplanar system reduced to a single resultant: the three forces are equivalent to $\vec{R} = 2\vec{i}\ \mathrm{N}$ acting along a horizontal line (schematic; the computation gives $y = 8$).}
\end{popfigure}

\courssub{Equilibrium of a Body Under Coplanar Forces}

A rigid body is in \cle{equilibrium} when it neither translates nor rotates. Translating is prevented by requiring the resultant force to vanish; rotating is prevented by requiring the total moment to vanish.

\begin{propriete}[Conditions of equilibrium under coplanar forces (proof examinable in outline)]
A rigid body is in equilibrium under a system of coplanar forces if and only if
\[
\sum F_x = 0, \qquad \sum F_y = 0, \qquad \sum M_O = 0 \quad \text{for any point } O.
\]
\emph{Justification.} If $\sum F_x = \sum F_y = 0$, the resultant force is zero, so the system reduces (by Section 4) either to nothing or to a couple. The couple, if any, has the same moment $\sum M_O$ about every point; imposing $\sum M_O = 0$ kills the couple too. Nothing remains: the body is in equilibrium. Conversely, if the body is in equilibrium, it has no tendency to translate (so $\sum F_x = \sum F_y = 0$) and no tendency to rotate about any point (so $\sum M_O = 0$).
\end{propriete}

\begin{attention}[Only three independent equations]
In the plane you have exactly \textbf{three} independent equilibrium equations, so you can determine at most \textbf{three} unknowns. Writing moments about a second point gives an equation that is a consequence of the first three — useful as a check, not as extra information.
\end{attention}

\begin{methode}[Choosing the moment centre wisely]
Take moments about the \textbf{point of intersection of two unknown forces}. Both unknowns then have zero moment arm and disappear from the moment equation, leaving an equation in the third unknown alone. This is the single most powerful time-saver in GCE statics problems.
\end{methode}

\begin{exoresolu}[Ladder against a wall]
A uniform ladder $AB$ of length $5\ \mathrm{m}$ and weight $W = 200\ \mathrm{N}$ leans against a smooth vertical wall at $A$ and rests on rough horizontal ground at $B$, with $AB$ making $60^{\circ}$ with the ground. Find the reaction at the wall and the frictional force at the ground.
\tcblower
\textbf{Solution.} The wall is smooth, so its reaction $R_A$ is horizontal (perpendicular to the wall); the ground reaction has a vertical component $N$ (the normal reaction) and a horizontal component $F$ (friction). Forces on the ladder: $W$ at the midpoint, $R_A$ at $A$, $N$ and $F$ at $B$.
\[
\sum F_y = 0:\quad N = W = 200\ \mathrm{N}; \qquad \sum F_x = 0:\quad F = R_A.
\]
Take moments about $B$ (this eliminates the two unknowns $N$ and $F$ acting at $B$). Perpendicular distances from $B$: the line of action of $W$ passes through the midpoint, at horizontal distance $\frac{5}{2}\cos 60^{\circ} = 1.25\ \mathrm{m}$ from $B$; the line of action of $R_A$ passes through $A$, at vertical height $5\sin 60^{\circ} = \frac{5\sqrt{3}}{2}\ \mathrm{m}$ above $B$. The weight tends to turn the ladder clockwise about $B$, the wall reaction anticlockwise:
\[
\sum M_B = 0:\quad R_A \cdot \frac{5\sqrt{3}}{2} - 200 \times 1.25 = 0
\;\Longrightarrow\; R_A = \frac{250}{4.330} \approx 57.7\ \mathrm{N}.
\]
Hence $F = R_A \approx 57.7\ \mathrm{N}$ and $N = 200\ \mathrm{N}$.
\end{exoresolu}

\courssub{Three Coplanar Forces in Equilibrium}

\begin{propriete}[Concurrency of three forces in equilibrium]
If a rigid body is in equilibrium under exactly \textbf{three non-parallel coplanar forces}, then the three forces are \cle{concurrent} (their lines of action meet at one point).
\end{propriete}

\textbf{Proof.} Suppose two of the forces, $\vec{F}_1$ and $\vec{F}_2$, meet at a point $P$ (they are not parallel, so they must meet). Take moments about $P$: $\vec{F}_1$ and $\vec{F}_2$ contribute nothing, since their lines of action pass through $P$. Equilibrium requires $\sum M_P = 0$, so the moment of $\vec{F}_3$ about $P$ is zero, which means the line of action of $\vec{F}_3$ also passes through $P$. $\blacksquare$

Consequently, since $\vec{F}_1 + \vec{F}_2 + \vec{F}_3 = \vec{0}$, the three force vectors placed head to tail form a closed triangle — the \cle{triangle of forces} — and from the sine rule in that triangle we obtain a flagged, exam-useful extra:

\begin{propriete}[Lami's theorem (useful extra, not a substitute for resolution)]
If three concurrent forces $\vec{P}, \vec{Q}, \vec{R}$ keep a body in equilibrium, and $\alpha, \beta, \gamma$ are the angles \emph{opposite} to them (i.e.\ the angles between the other two forces), then
\[
\frac{P}{\sin\alpha} = \frac{Q}{\sin\beta} = \frac{R}{\sin\gamma}.
\]
\end{propriete}

\begin{attention}[When Lami's theorem may NOT be used]
Do not apply Lami's theorem unless \textbf{exactly three} forces act and they are \textbf{concurrent}. A ladder on rough ground has four forces (weight, two reaction components at the ground, wall reaction): Lami is forbidden there. When in doubt, resolve and take moments — that method always works.
\end{attention}

\begin{popfigure}
\begin{center}
\begin{tikzpicture}[scale=1.1, >=stealth]
% concurrent forces
\fill[popdark] (0,0) circle (2pt);
\draw[->, very thick, popblue] (0,0) -- (2.2,0) node[right] {$\vec{F}_1$};
\draw[->, very thick, poporange] (0,0) -- (-1.1,1.9) node[above left] {$\vec{F}_2$};
\draw[->, very thick, popgreen] (0,0) -- (-1.1,-1.9) node[below left] {$\vec{F}_3$};
\node[popdark] at (0.25,-0.28) {$P$};
% triangle of forces
\begin{scope}[xshift=6.5cm, yshift=-1cm]
\draw[->, very thick, popblue] (0,0) -- (2.2,0) node[midway, below] {$\vec{F}_1$};
\draw[->, very thick, poporange] (2.2,0) -- (1.1,1.9) node[midway, right] {$\vec{F}_2$};
\draw[->, very thick, popgreen] (1.1,1.9) -- (0,0) node[midway, left] {$\vec{F}_3$};
\node at (1.1,-0.9) {triangle of forces};
\end{scope}
\end{tikzpicture}
\end{center}
\legende{Three concurrent coplanar forces in equilibrium (left); placed head to tail they form a closed triangle (right), from which Lami's theorem follows by the sine rule.}
\end{popfigure}

\begin{exemplebox}[Reduction to a couple]
Forces $5\ \mathrm{N}$ act along $AB$ and $5\ \mathrm{N}$ along $CD$ of a square $ABCD$ of side $2\ \mathrm{m}$, in opposite senses, with $AB \parallel CD$. Then $\sum \vec{F} = \vec{0}$, but moments about the centre give $5 \times 1 + 5 \times 1 = 10\ \mathrm{N\,m}$ (each force has moment arm $1\ \mathrm{m}$ from the centre, and both turn the same way). The system cannot be reduced to a single force: it is a \textbf{couple} of moment $10\ \mathrm{N\,m}$. Remember the test: $\vec{R} = \vec{0}$ and $\sum M \neq 0 \Rightarrow$ couple.
\end{exemplebox}

\courssub{Simple Systems of Forces in Three Dimensions}

Many real structures — a mast held by cables on a building site in Douala, a sign hanging from three ropes — are genuinely three-dimensional. Fortunately, the ideas extend naturally.

\textbf{Resultant of concurrent forces in 3D.} If forces $\vec{F}_1, \dots, \vec{F}_n$ act at the \emph{same particle}, their resultant is
\[
\vec{R} = \sum \vec{F}_i = \Big(\sum F_{ix}\Big)\vec{i} + \Big(\sum F_{iy}\Big)\vec{j} + \Big(\sum F_{iz}\Big)\vec{k},
\qquad |\vec{R}| = \sqrt{R_x^2 + R_y^2 + R_z^2}.
\]

\textbf{Equilibrium of a particle in 3D.} A particle is in equilibrium if and only if $\sum \vec{F} = \vec{0}$, i.e.\ the \textbf{three} scalar equations
\[
\sum F_x = 0, \qquad \sum F_y = 0, \qquad \sum F_z = 0.
\]

\begin{methode}[The unit-vector method for cable problems]
\begin{enumerate}
\item Write the position vector of each anchor point relative to the particle, e.g.\ $\vec{PA} = \vec{OA} - \vec{OP}$.
\item Compute its length $|\vec{PA}|$ and form the unit vector $\hat{u} = \dfrac{\vec{PA}}{|\vec{PA}|}$ along the cable.
\item Write each cable force as $\vec{T} = T\,\hat{u}$ (tension times unit vector).
\item Sum all forces, set each of the three components to zero, and solve the resulting system.
\end{enumerate}
\end{methode}

\begin{attention}[Three equations, not two]
In 3D, equilibrium requires \textbf{all three} component equations. Solving only $\sum F_x = 0$ and $\sum F_y = 0$ and forgetting $\sum F_z = 0$ is a frequent and costly error — especially when a weight acts in the $z$-direction.
\end{attention}

\begin{exoresolu}[Particle held by three cables]
A particle $P$ of weight $100\ \mathrm{N}$ hangs at $P(1,1,0)$, held by three cables anchored at $A(0,0,2)$, $B(2,0,2)$ and $C(1,3,2)$. Find the tensions $T_1, T_2, T_3$ in $PA$, $PB$, $PC$.
\tcblower
\textbf{Solution.} Position vectors from $P$:
\[
\vec{PA} = -\vec{i} - \vec{j} + 2\vec{k}, \quad
\vec{PB} = \vec{i} - \vec{j} + 2\vec{k}, \quad
\vec{PC} = 0\vec{i} + 2\vec{j} + 2\vec{k},
\]
with lengths $|\vec{PA}| = |\vec{PB}| = \sqrt{1+1+4} = \sqrt{6}$ and $|\vec{PC}| = \sqrt{0+4+4} = \sqrt{8} = 2\sqrt{2}$. The forces on $P$ are
\[
\frac{T_1}{\sqrt{6}}(-\vec{i} - \vec{j} + 2\vec{k}), \quad
\frac{T_2}{\sqrt{6}}(\vec{i} - \vec{j} + 2\vec{k}), \quad
\frac{T_3}{2\sqrt{2}}(2\vec{j} + 2\vec{k}), \quad -100\vec{k}.
\]
Component equations:
\begin{align*}
x:&\quad \frac{-T_1 + T_2}{\sqrt{6}} = 0 \;\Rightarrow\; T_1 = T_2,\\
y:&\quad \frac{-T_1 - T_2}{\sqrt{6}} + \frac{T_3}{\sqrt{2}} = 0,\\
z:&\quad \frac{2T_1 + 2T_2}{\sqrt{6}} + \frac{T_3}{\sqrt{2}} - 100 = 0.
\end{align*}
Setting $T_1 = T_2 = T$: the $y$-equation gives $\dfrac{-2T}{\sqrt{6}} + \dfrac{T_3}{\sqrt{2}} = 0$, so $T_3 = \dfrac{2T\sqrt{2}}{\sqrt{6}} = \dfrac{2T}{\sqrt{3}}$. Then $\dfrac{T_3}{\sqrt{2}} = \dfrac{2T}{\sqrt{6}}$, and substituting in the $z$-equation:
\[
\frac{4T}{\sqrt{6}} + \frac{2T}{\sqrt{6}} = 100 \;\Rightarrow\; \frac{6T}{\sqrt{6}} = 100 \;\Rightarrow\; T = \frac{100\sqrt{6}}{6} \approx 40.8\ \mathrm{N}.
\]
Hence $T_1 = T_2 \approx 40.8\ \mathrm{N}$ and $T_3 = \dfrac{2 \times 40.8}{\sqrt{3}} \approx 47.1\ \mathrm{N}$.
\end{exoresolu}

\begin{popfigure}
\begin{center}
\begin{tikzpicture}[scale=1.0, >=stealth]
% 3D axes
\draw[->, gray] (0,0) -- (4.2,0) node[right] {$x$};
\draw[->, gray] (0,0) -- (-1.8,-1.4) node[below left] {$y$};
\draw[->, gray] (0,0) -- (0,3.6) node[above] {$z$};
% anchors (z=2 plane, shifted up)
\coordinate (A) at (0,2.4);
\coordinate (B) at (3,2.4);
\coordinate (C) at (-0.9,1.7);
\coordinate (P) at (0.9,0.4);
\fill[popblue] (A) circle (1.8pt) node[left] {$A$};
\fill[popblue] (B) circle (1.8pt) node[right] {$B$};
\fill[popblue] (C) circle (1.8pt) node[left] {$C$};
\draw[thick, poporange] (P) -- (A) node[midway, left, poporange] {$\hat{u}_1$};
\draw[thick, poporange] (P) -- (B) node[midway, right, poporange] {$\hat{u}_2$};
\draw[thick, poporange] (P) -- (C) node[midway, above, poporange] {$\hat{u}_3$};
\fill[popdark] (P) circle (2.2pt) node[below right] {$P$};
\draw[->, very thick, popgreen] (P) -- (0.9,-1.1) node[below] {$\vec{W}$};
\end{tikzpicture}
\end{center}
\legende{A particle $P$ held by three cables anchored at $A$, $B$, $C$. Each tension acts along the unit vector from $P$ towards its anchor; the weight acts vertically downwards.}
\end{popfigure}

\textbf{Moments in 3D (extension for strong candidates).} The moment of a force $\vec{F}$ about a point $O$ is the vector
\[
\vec{M}_O = \vec{r} \times \vec{F},
\]
where $\vec{r}$ is the position vector of any point on the line of action of $\vec{F}$. A rigid body is in equilibrium in 3D when $\sum \vec{F} = \vec{0}$ \emph{and} $\sum \vec{M}_O = \vec{0}$ — six scalar equations in all. Only simple cases (e.g.\ a single moment computed with the cross product) are expected at this level.

\begin{exemplebox}[Moment via the cross product]
A force $\vec{F} = (2\vec{i} + \vec{j} - \vec{k})\ \mathrm{N}$ acts at the point $(1, 0, 2)$. Its moment about $O$ is
\[
\vec{M}_O = \begin{pmatrix} 1 \\ 0 \\ 2 \end{pmatrix} \times \begin{pmatrix} 2 \\ 1 \\ -1 \end{pmatrix}
= \begin{pmatrix} 0(-1) - 2(1) \\ 2(2) - 1(-1) \\ 1(1) - 0(2) \end{pmatrix}
= \begin{pmatrix} -2 \\ 5 \\ 1 \end{pmatrix}\ \mathrm{N\,m}.
\]
Check: the $y$-component is $r_z F_x - r_x F_z = 2(2) - 1(-1) = 5$, as required by the cross-product formula.
\end{exemplebox}

\courssub{Synthesis: GCE-Style Combined Problems}

Examination questions typically combine resolution, moments and (sometimes) couples in one situation. The strategy never changes: draw a clear force diagram, resolve, and take moments about the intersection of two unknowns.

\begin{exoresolu}[Hinged rod with a cable]
A uniform rod $AB$ of length $4\ \mathrm{m}$ and weight $80\ \mathrm{N}$ is hinged at $A$ to a wall and held horizontal by a cable from $B$ to a point $C$ on the wall $3\ \mathrm{m}$ above $A$. Find the tension in the cable and the horizontal and vertical components of the reaction at $A$.
\tcblower
\textbf{Solution.} The cable $BC$ has length $\sqrt{4^2+3^2} = 5\ \mathrm{m}$, so the tension $T$ at $B$ pulls along $BC$ with components $\big(-\tfrac{4}{5}T,\ \tfrac{3}{5}T\big)$. Let the hinge reaction be $(H, V)$.
Moments about $A$ (eliminating $H$ and $V$): the weight acts at the midpoint, $2\ \mathrm{m}$ from $A$; the horizontal part of $T$ acts along the rod and has no moment about $A$, so only the vertical part of $T$ contributes:
\[
\sum M_A = 0:\quad \frac{3}{5}T \times 4 - 80 \times 2 = 0 \;\Rightarrow\; T = \frac{160 \times 5}{12} \approx 66.7\ \mathrm{N}.
\]
Then
\[
\sum F_x = 0:\ H = \frac{4}{5}T \approx 53.3\ \mathrm{N}; \qquad
\sum F_y = 0:\ V = 80 - \frac{3}{5}T = 80 - 40 = 40\ \mathrm{N}.
\]
\end{exoresolu}

\begin{exercice}[Practice]
\begin{enumerate}
\item Forces $(2\vec{i}+\vec{j})\ \mathrm{N}$ at $(1,0)$, $(-\vec{i}+3\vec{j})\ \mathrm{N}$ at $(0,2)$ and $(-\vec{i}-4\vec{j})\ \mathrm{N}$ at $(3,1)$ act on a lamina. Show that the system reduces to a couple and find its moment.
\item A particle of weight $50\ \mathrm{N}$ hangs at $P(0,0,0)$ from two cables anchored at $A(0,3,4)$ and $B(0,-3,4)$. Find the two tensions.
\item Three concurrent forces keep a body in equilibrium; two of them are $10\ \mathrm{N}$ and $14\ \mathrm{N}$ with an angle of $120^{\circ}$ between them. Use the triangle of forces to find the third force.
\end{enumerate}
\end{exercice}

\begin{corrige}[Exercise 1]
$R_x = 2 - 1 - 1 = 0$ and $R_y = 1 + 3 - 4 = 0$, so $\vec{R} = \vec{0}$. Moments about $O$ (using $M_O = xF_y - yF_x$): $M_1 = 1(1) - 0(2) = 1$; $M_2 = 0(3) - 2(-1) = 2$; $M_3 = 3(-4) - 1(-1) = -12 + 1 = -11$. Total $\sum M_O = 1 + 2 - 11 = -8\ \mathrm{N\,m} \neq 0$: the system is a couple of moment $8\ \mathrm{N\,m}$ clockwise.
\end{corrige}

\begin{corrige}[Exercise 2]
$\vec{PA} = (0,3,4)$ and $\vec{PB} = (0,-3,4)$, with $|\vec{PA}| = |\vec{PB}| = \sqrt{9+16} = 5$. The forces on $P$ are $\dfrac{T_1}{5}(3\vec{j}+4\vec{k})$, $\dfrac{T_2}{5}(-3\vec{j}+4\vec{k})$ and $-50\vec{k}$.
\[
\sum F_y = 0:\ \frac{3T_1 - 3T_2}{5} = 0 \Rightarrow T_1 = T_2; \qquad
\sum F_z = 0:\ \frac{4T_1 + 4T_2}{5} = 50.
\]
With $T_1 = T_2 = T$: $\dfrac{8T}{5} = 50$, so $T = \dfrac{250}{8} = 31.25\ \mathrm{N}$. Hence $T_1 = T_2 = 31.25\ \mathrm{N}$. (The $x$-equation is satisfied identically: the geometry has no $x$-components.)
\end{corrige}

\begin{corrige}[Exercise 3]
In the triangle of forces, the angle between the sides representing the $10\ \mathrm{N}$ and $14\ \mathrm{N}$ forces is the supplement $180^{\circ} - 120^{\circ} = 60^{\circ}$ (head-to-tail placement reverses one force). By the cosine rule, the third force $R$ closes the triangle:
\[
R^2 = 10^2 + 14^2 - 2(10)(14)\cos 60^{\circ} = 100 + 196 - 140 = 156,
\]
so $R = \sqrt{156} \approx 12.5\ \mathrm{N}$.
\end{corrige}

\begin{aretenir}[Key points of Section 5]
\begin{itemize}
\item Resultant of coplanar forces: $\vec{R} = \sum \vec{F}_i$ by components; its line of action is fixed by $\sum M_O(\vec{F}_i) = M_O(\vec{R})$, with $M_O = xF_y - yF_x$.
\item Equilibrium in the plane: $\sum F_x = 0$, $\sum F_y = 0$, $\sum M_O = 0$ — three equations, three unknowns at most.
\item Take moments about the intersection of two unknown forces to eliminate them.
\item Three non-parallel coplanar forces in equilibrium are concurrent; then (and only then) the triangle of forces and Lami's theorem apply.
\item $\vec{R} = \vec{0}$ with $\sum M \neq 0$ means the system is a couple.
\item In 3D: write each cable force as $T\,\hat{u}$ with $\hat{u} = \vec{PA}/|\vec{PA}|$, and solve all \textbf{three} component equations; moments use $\vec{M}_O = \vec{r} \times \vec{F}$.
\end{itemize}
\end{aretenir}

\newpage

\coursec{Integration activity}

\coursec{Comprehensive Application: Static Analysis of the Mile 2 Market Awning}

\courssub{Aims of the Activity}
\begin{itemize}
    \item Consolidate the concepts of vector components, work done by a constant force, moment of a force, couples, and equilibrium of coplanar forces in a single authentic engineering task.
    \item Develop the ability to model a real‑world structure using rigid‑body mechanics: drawing a free‑body diagram, resolving forces, applying equilibrium conditions ($\sum \vec{F}=\vec{0}$, $\sum \vec{M}=\vec{0}$), and interpreting the physical meaning of the results.
    \item Practise systematic problem‑solving and clear technical communication.
\end{itemize}

\courssub{Problem Statement}
In the bustling \emph{Mile 2 Market} in Bamenda, a trader named Mama Ngo wants to install a sturdy horizontal awning to protect her goods. The awning is modelled as a \textbf{uniform plank} of weight $W = 200\,\mathrm{N}$ and length $L = 3\,\mathrm{m}$. One end $A$ is hinged to a vertical wall; the other end $B$ is supported by a light inextensible cable attached to a point $C$ on the wall, vertically $2\,\mathrm{m}$ above the hinge $A$. The cable makes an angle $\theta$ with the horizontal. The system is in static equilibrium.

\begin{popfigure}
\begin{tikzpicture}[scale=0.85, >=stealth]
    % Coordinate axes
    \draw[popmuted, thin, ->] (-1,0) -- (4,0) node[right] {$x$};
    \draw[popmuted, thin, ->] (0,-1) -- (0,3) node[above] {$y$};
    % Wall
    \draw[popdark, line width=3pt] (-0.3,0) -- (-0.3,3);
    % Hinge A
    \fill[popblue] (0,0) circle (4pt);
    \node[left] at (-0.3,0) {$A(0,0)$};
    % Plank AB
    \draw[poporange, line width=5pt] (0,0) -- (3,0);
    \node[below] at (1.5,-0.4) {$L=3\,\mathrm{m}$};
    % Point B
    \fill[popblue] (3,0) circle (4pt);
    \node[right] at (3,0) {$B(3,0)$};
    % Point C on wall
    \fill[popblue] (0,2) circle (4pt);
    \node[left] at (-0.3,2) {$C(0,2)$};
    % Cable BC
    \draw[popgreen, dashed, thick] (3,0) -- (0,2);
    % Centre of gravity G
    \fill[popred] (1.5,0) circle (3pt);
    \node[below] at (1.5,-0.2) {$G$};
    % Weight at G
    \draw[->, popred, thick] (1.5,0) -- (1.5,-1.2) node[below] {$\vec{W}=200\,\mathrm{N}$};
    % Tension at B
    \draw[->, poppurple, thick] (3,0) -- (2.1,0.7) node[above right] {$\vec{T}$};
    % Reaction at A
    \draw[->, poppink, thick] (0,0) -- (-1,0.6) node[above left] {$\vec{R}$};
    % Angle theta
    \draw[popmuted] (2.5,0) arc (0:33.69:0.6) node[midway, right, xshift=2pt] {$\theta$};
    % Dimensions
    \draw[<->, popmuted] (-1.2,0) -- (-1.2,2) node[midway, left] {$2\,\mathrm{m}$};
    \draw[<->, popmuted] (0,-0.7) -- (3,-0.7) node[midway, below] {$3\,\mathrm{m}$};
    % Force components (ghost)
    \draw[poppurple, thin, dashed] (3,0) -- (1.5,0) node[midway, below] {$T_x$};
    \draw[poppurple, thin, dashed] (1.5,0) -- (1.5,1) node[midway, right] {$T_y$};
    \draw[poppink, thin, dashed] (0,0) -- (1.5,0) node[midway, below] {$R_x$};
    \draw[poppink, thin, dashed] (1.5,0) -- (1.5,0.6) node[midway, right] {$R_y$};
\end{tikzpicture}
\legende{Free‑body diagram of the awning: hinge at $A$, cable $BC$, weight $\vec{W}$ at centre $G$, tension $\vec{T}$, reaction $\vec{R}$. Components shown dashed.}
\end{popfigure}

The trader asks a group of Upper Sixth Further Mathematics students to carry out a complete static analysis. The analysis must determine:
\begin{enumerate}
    \item The angle $\theta$ and the vector components of all forces.
    \item The tension $T$ in the cable using moment equilibrium about $A$.
    \item The reaction $\vec{R}$ at the hinge using force equilibrium.
    \item The work done by the tension if the awning is lowered slightly ($\delta y = 0.01\,\mathrm{m}$ at $B$).
    \item The effect of adding a second symmetric cable from $B$ to $D(0,-2)$: couple formation, concurrency condition, and viability.
    \item A safety assessment (cable max $500\,\mathrm{N}$, hinge max $400\,\mathrm{N}$).
\end{enumerate}

\courssub{Guided Solution}

\textbf{1. Geometry and Vector Components}\par
The cable runs from $B(3,0)$ to $C(0,2)$. The horizontal distance is $3\,\mathrm{m}$, the vertical distance is $2\,\mathrm{m}$.
\[
\tan\theta = \frac{2}{3} \quad\Rightarrow\quad \theta = \arctan\left(\frac{2}{3}\right) \approx 33.69^\circ.
\]
The length of the cable is $|\vec{BC}| = \sqrt{3^2+2^2} = \sqrt{13}\,\mathrm{m}$.
The unit vector along $\vec{BC}$ (from $B$ toward $C$) is
\[
\hat{u}_{BC} = \frac{\vec{BC}}{|\vec{BC}|} = \frac{-3\vec{i} + 2\vec{j}}{\sqrt{13}}.
\]
Since the tension pulls $B$ toward $C$, the tension vector is
\[
\vec{T} = T\,\hat{u}_{BC} = T\left(-\frac{3}{\sqrt{13}}\vec{i} + \frac{2}{\sqrt{13}}\vec{j}\right)
= -\frac{3T}{\sqrt{13}}\,\vec{i} + \frac{2T}{\sqrt{13}}\,\vec{j}.
\]
Thus the components are $T_x = -\dfrac{3T}{\sqrt{13}}$, $T_y = \dfrac{2T}{\sqrt{13}}$.

The weight acts vertically downward at the centre of gravity $G(1.5,0)$:
\[
\vec{W} = -200\,\vec{j}\ \mathrm{N}.
\]
The reaction at the hinge $A$ has unknown components:
\[
\vec{R} = R_x\,\vec{i} + R_y\,\vec{j}.
\]

\begin{aretenir}[Resolving Forces]
\begin{itemize}
\item The component of a vector $\vec{F}$ in the direction of a unit vector $\hat{u}$ is $\vec{F}\cdot\hat{u}$.
\item For a force $\vec{F}$ making an angle $\theta$ with the horizontal: $F_x = F\cos\theta$, $F_y = F\sin\theta$.
\item Always define your coordinate system ($\vec{i}$ horizontal $+$, $\vec{j}$ vertical $+$) and stick to it.
\end{itemize}
\end{aretenir}

\textbf{2. Moment Equilibrium about Hinge $A$}\par
We take moments about $A$ (clockwise positive). The forces acting are $\vec{W}$ at $G$, $\vec{T}$ at $B$, and $\vec{R}$ at $A$. The reaction $\vec{R}$ has zero moment about $A$ because its line of action passes through $A$.

\begin{itemize}
    \item \textbf{Weight $\vec{W}$:} Acts vertically downward at $G$, horizontal distance $AG = 1.5\,\mathrm{m}$ from $A$. The perpendicular distance from $A$ to the line of action of $\vec{W}$ is $1.5\,\mathrm{m}$. It tends to rotate the plank clockwise $\Rightarrow$ positive moment:
    \[
    M_W = +W \times AG = +200 \times 1.5 = +300\,\mathrm{N\,m}.
    \]
    \item \textbf{Tension $\vec{T}$:} Only the vertical component $T_y = \frac{2T}{\sqrt{13}}$ produces a moment about $A$ (the horizontal component $T_x$ passes through $A$). $T_y$ acts upward at $B$, distance $AB = 3\,\mathrm{m}$ from $A$, tending to rotate anticlockwise $\Rightarrow$ negative moment:
    \[
    M_T = -T_y \times AB = -\frac{2T}{\sqrt{13}} \times 3 = -\frac{6T}{\sqrt{13}}\,\mathrm{N\,m}.
    \]
\end{itemize}
Equilibrium condition $\sum M_A = 0$:
\[
300 - \frac{6T}{\sqrt{13}} = 0 \quad\Rightarrow\quad T = \frac{300\sqrt{13}}{6} = 50\sqrt{13}\,\mathrm{N}.
\]
Numerically, $\sqrt{13}\approx 3.60555$, so $T \approx 180.3\,\mathrm{N}$.

\begin{attention}[Moment Calculation]
\begin{itemize}
\item The moment of a force about a point is $M = F \times d_\perp$, where $d_\perp$ is the perpendicular distance from the point to the \emph{line of action} of the force.
\item Alternatively, $M = F_\perp \times d$, where $F_\perp$ is the component of the force perpendicular to the line joining the point to the point of application, and $d$ is that distance.
\item Choose the sign convention (clockwise $+$ or anticlockwise $+$) and apply it consistently.
\end{itemize}
\end{attention}

\textbf{3. Force Equilibrium and Hinge Reaction}\par
Apply $\sum F_x = 0$ and $\sum F_y = 0$.

$\sum F_x = 0$:
\[
R_x + T_x = 0 \quad\Rightarrow\quad R_x = -T_x = \frac{3T}{\sqrt{13}} = \frac{3 \times 50\sqrt{13}}{\sqrt{13}} = 150\,\mathrm{N}.
\]

$\sum F_y = 0$:
\[
R_y + T_y - W = 0 \quad\Rightarrow\quad R_y = W - T_y = 200 - \frac{2T}{\sqrt{13}} = 200 - \frac{2 \times 50\sqrt{13}}{\sqrt{13}} = 200 - 100 = 100\,\mathrm{N}.
\]

Hence the reaction vector is
\[
\vec{R} = 150\,\vec{i} + 100\,\vec{j}\ \mathrm{N}.
\]
Its magnitude:
\[
R = \sqrt{R_x^2 + R_y^2} = \sqrt{150^2 + 100^2} = \sqrt{32500} = 50\sqrt{13} \approx 180.3\,\mathrm{N}.
\]
Its direction (angle $\alpha$ above the horizontal):
\[
\tan\alpha = \frac{R_y}{R_x} = \frac{100}{150} = \frac{2}{3} \quad\Rightarrow\quad \alpha = \arctan\left(\frac{2}{3}\right) \approx 33.69^\circ.
\]

\begin{aretenir}[Equilibrium Conditions]
\begin{itemize}
\item For a rigid body in static equilibrium in a plane:
      $\sum F_x = 0$, $\sum F_y = 0$, $\sum M_{\text{any point}} = 0$.
\item Choosing the moment centre at a point where unknown forces act (like the hinge $A$) eliminates those unknowns from the moment equation.
\item The three forces $\vec{W}$, $\vec{T}$, $\vec{R}$ are \emph{not} concurrent (their lines of action do not meet at a single point). This is consistent because there are three forces but they are not the \emph{only} forces? Wait — they are the only forces. For a body in equilibrium under three non-parallel forces, the lines of action \textbf{must} be concurrent. Let's check: $\vec{W}$ is vertical through $G(1.5,0)$. $\vec{T}$ acts along $BC$. $\vec{R}$ acts along the line from $A$ at angle $\alpha$. The lines $BC$ and $AR$ intersect at some point. Does the vertical through $G$ pass through that intersection? In this problem, the three forces \textbf{are} concurrent (a classic result for a simply supported beam with a cable). The intersection of $BC$ and the line of $\vec{R}$ lies on the vertical through $G$. This can be verified geometrically.
\end{itemize}
\end{aretenir}

\textbf{4. Work Done by the Tension for a Small Displacement}\par
The awning is slowly lowered so that $B$ moves vertically downward by $\delta y = 0.01\,\mathrm{m}$ while $A$ remains fixed. The plank rotates about $A$ by a small angle $\delta\phi$. To first order, the displacement of $B$ is purely vertical:
\[
\delta \vec{r}_B = -0.01\,\vec{j}\ \mathrm{m}.
\]
The work done by a constant force $\vec{T}$ during a displacement $\delta \vec{r}$ is $\delta W = \vec{T} \cdot \delta \vec{r}$.
\[
\delta W = \left(-\frac{3T}{\sqrt{13}}\vec{i} + \frac{2T}{\sqrt{13}}\vec{j}\right) \cdot \left(-0.01\,\vec{j}\right)
= \frac{2T}{\sqrt{13}} \times (-0.01) = -\frac{0.02T}{\sqrt{13}}\,\mathrm{J}.
\]
Substitute $T = 50\sqrt{13}\,\mathrm{N}$:
\[
\delta W = -\frac{0.02 \times 50\sqrt{13}}{\sqrt{13}} = -1.0\,\mathrm{J}.
\]

\begin{aretenir}[Work Done by a Constant Force]
\begin{itemize}
\item $W = \vec{F} \cdot \vec{d} = F d \cos\phi$, where $\phi$ is the angle between $\vec{F}$ and $\vec{d}$.
\item In component form: $W = F_x d_x + F_y d_y$.
\item \textbf{Negative work} means the force component opposes the displacement; energy is transferred \emph{from} the body \emph{to} the force source (or stored as potential energy in the cable). Here, the tension pulls upward while $B$ moves downward, so the tension does negative work. An external agent must do positive work to lower the awning.
\end{itemize}
\end{aretenir}

\textbf{5. Adding a Second Symmetric Cable: Couple and Concurrency Analysis}\par
A second identical cable is attached from $B$ to $D(0,-2)$ (vertically $2\,\mathrm{m}$ below $A$). The two cables make equal angles $\theta = \arctan(2/3)$ with the horizontal, one above, one below. By symmetry, the tensions have equal magnitude $T'$ (to be determined). Their vectors are:
\[
\vec{T}_1 = T'\left(-\frac{3}{\sqrt{13}}\vec{i} + \frac{2}{\sqrt{13}}\vec{j}\right),\qquad
\vec{T}_2 = T'\left(-\frac{3}{\sqrt{13}}\vec{i} - \frac{2}{\sqrt{13}}\vec{j}\right).
\]

\begin{enumerate}
    \item \textbf{Do $\vec{T}_1$ and $\vec{T}_2$ form a couple?}
    A \textbf{couple} consists of two forces that are:
    \begin{itemize}
        \item equal in magnitude,
        \item opposite in direction (parallel lines of action, anti-parallel vectors),
        \item not collinear (different lines of action).
    \end{itemize}
    Here, $\vec{T}_1$ and $\vec{T}_2$ have the \emph{same} horizontal component (both leftward) and \emph{opposite} vertical components. They are \textbf{not} opposite vectors; their resultant is a pure horizontal force:
    \[
    \vec{T}_1 + \vec{T}_2 = -2T'\frac{3}{\sqrt{13}}\,\vec{i} = -\frac{6T'}{\sqrt{13}}\,\vec{i}.
    \]
    Since the resultant force is not zero, they do \textbf{not} form a couple. They are equivalent to a single horizontal force acting at $B$.

    \item \textbf{Three‑force concurrency condition.}
    The forces on the plank are now $\vec{W}$ (at $G$), $\vec{T}_1$ (along $BC$), $\vec{T}_2$ (along $BD$), and $\vec{R}$ (at $A$) — four forces. The three‑force concurrency theorem applies \emph{only} when exactly three forces act. If we consider the resultant of the two cables, $\vec{T}_{\text{res}} = \vec{T}_1+\vec{T}_2$ (horizontal at $B$), then we have three forces: $\vec{W}$, $\vec{T}_{\text{res}}$, $\vec{R}$. For equilibrium, their lines of action must be concurrent. $\vec{T}_{\text{res}}$ is horizontal through $B$; $\vec{W}$ is vertical through $G$; they intersect at $(3,0)$? No, $\vec{W}$ is at $x=1.5$, $\vec{T}_{\text{res}}$ is at $y=0$. They intersect at $(1.5,0)$? Wait: $\vec{T}_{\text{res}}$ acts at $B(3,0)$ horizontally. Its line of action is the $x$-axis ($y=0$). $\vec{W}$ acts at $G(1.5,0)$ vertically. Its line of action is $x=1.5$. These two lines intersect at $(1.5,0) = G$. The reaction $\vec{R}$ at $A(0,0)$ must also pass through this intersection for concurrency. But $\vec{R}$ passes through $A(0,0)$ and $G(1.5,0)$ only if it is horizontal. However, $\vec{R}$ must balance the vertical weight, so it has a vertical component. Thus the three lines are \textbf{not concurrent} — confirming that equilibrium is impossible with only these three forces (or four forces with a hinge that provides no moment). The hinge reaction $\vec{R}$ provides the necessary fourth force to satisfy equilibrium, but as we see next, even four forces cannot satisfy moment equilibrium with this geometry.

    \item \textbf{Moment equilibrium with two cables (viability check).}
    Take moments about $A$ again. The forces with moments about $A$ are $\vec{W}$ and the vertical components of $\vec{T}_1$ and $\vec{T}_2$ (horizontal components pass through $A$).
    \begin{align*}
        M_W &= +W \times 1.5 = +300\,\mathrm{N\,m} \quad\text{(clockwise)} \\
        M_{T_1} &= -T'_y \times 3 = -\frac{2T'}{\sqrt{13}} \times 3 = -\frac{6T'}{\sqrt{13}}\,\mathrm{N\,m} \quad\text{(anticlockwise)} \\
        M_{T_2} &= +T'_y \times 3 = +\frac{2T'}{\sqrt{13}} \times 3 = +\frac{6T'}{\sqrt{13}}\,\mathrm{N\,m} \quad\text{(clockwise, since $T_{2y}$ is downward)}
    \end{align*}
    The moments of the two vertical components \textbf{cancel exactly} ($M_{T_1}+M_{T_2}=0$). The weight moment $+300\,\mathrm{N\,m}$ remains unbalanced. The hinge $A$ cannot provide a moment (it is a simple hinge). Therefore:
    \[
    \sum M_A = 300 \neq 0.
    \]
    \textbf{Conclusion:} The symmetric two‑cable configuration \textbf{cannot be in static equilibrium} with a simple hinge at $A$. The awning would rotate clockwise. To make it work, one would need:
    \begin{itemize}
        \item A fixed support (built-in end) at $A$ that can exert a reaction moment, or
        \item The second cable attached at a different point on the plank (not at $B$), or
        \item The plank not horizontal (inclined).
    \end{itemize}
\end{enumerate}

\begin{savaistu}[Real‑World Design Insight]
In practice, market stall awnings in Cameroon often use a single inclined strut (a compression member) or a single cable with the hinge providing both horizontal and vertical reaction. Adding a second symmetric cable at the same point creates a "redundant" constraint that fights the weight's moment instead of helping it. Engineers use \emph{static determinacy} checks: a planar rigid body has 3 equilibrium equations. A hinge (2 reactions) + 1 cable (1 tension) = 3 unknowns $\rightarrow$ statically determinate (solvable). Hinge + 2 cables = 4 unknowns $\rightarrow$ statically indeterminate; but here the geometry makes it impossible to satisfy all equations simultaneously — a \emph{mechanism} or \emph{improper constraint}.
\end{savaistu}

\courssub{Safety Assessment and Conclusion}

\begin{center}
\begin{tabularx}{\linewidth}{|l|X|X|X|}
\hline
\rowcolor{popblueL}
\textbf{Component} & \textbf{Calculated Value} & \textbf{Maximum Allowable} & \textbf{Status} \\
\hline
Cable Tension $T$ & $50\sqrt{13} \approx 180.3\,\mathrm{N}$ & $500\,\mathrm{N}$ & \cellcolor{popgreenL} SAFE (Margin $\approx 2.8\times$) \\
\hline
Hinge Reaction $R$ & $50\sqrt{13} \approx 180.3\,\mathrm{N}$ & $400\,\mathrm{N}$ & \cellcolor{popgreenL} SAFE (Margin $\approx 2.2\times$) \\
\hline
\end{tabularx}
\end{center}

The single‑cable design meets the safety requirements comfortably. The proposed symmetric two‑cable modification is \textbf{not viable} with a simple hinge because it cannot satisfy moment equilibrium. A different arrangement (e.g., attaching the second cable to a point on the plank between $A$ and $B$, or using a fixed support at $A$) would be needed for a two‑cable system.

\courssub{Key Points for Examination Revision}
\begin{aretenir}[Essential Summary]
\begin{itemize}
\item \textbf{Free‑body diagram} is the starting point: identify all forces, their points of application, and directions.
\item \textbf{Vector components} simplify equilibrium: $\sum F_x=0$, $\sum F_y=0$, $\sum M=0$.
\item \textbf{Moment equilibrium} about a convenient point (often a hinge) eliminates unknown reactions.
\item \textbf{Work} $W = \vec{F}\cdot\vec{d}$; negative work means the force opposes the displacement.
\item \textbf{Couple}: two equal, opposite, parallel forces with different lines of action $\rightarrow$ pure moment, no resultant force.
\item \textbf{Three‑force concurrency}: if a rigid body is in equilibrium under three non‑parallel forces, their lines of action must meet at a single point.
\item \textbf{Statically determinate} systems have exactly as many unknowns as independent equilibrium equations.
\item Always \textbf{check safety margins} against material limits.
\end{itemize}
\end{aretenir}

\courssub{Common Pitfalls to Avoid}
\begin{attention}[Examiner's Traps]
\begin{itemize}
\item Forgetting that the weight of a uniform plank acts at its \emph{centre}, not at the end.
\item Using the wrong perpendicular distance when calculating moments (always the perpendicular distance from the pivot to the \emph{line of action} of the force).
\item Confusing the sign convention for moments (clockwise vs anticlockwise) and for work (positive when force and displacement are in the same general direction).
\item Assuming two symmetric forces automatically form a couple; they must be \emph{equal and opposite} (anti-parallel).
\item Neglecting that a hinge can exert a reaction force but \emph{no moment}.
\item In the two-cable scenario, assuming the tension magnitude remains the same as in the single-cable case. The equilibrium equations change, so the tensions change (or equilibrium becomes impossible).
\end{itemize}
\end{attention}

\begin{exercice}[Practice Problem]
A uniform rod $AB$ of length $4\,\mathrm{m}$ and weight $150\,\mathrm{N}$ is hinged at $A$ to a vertical wall. It is held horizontal by a light inextensible string attached at $B$ and to a point $C$ on the wall $3\,\mathrm{m}$ above $A$. A load of $100\,\mathrm{N}$ is suspended from the midpoint of the rod. Find:
\begin{enumerate}
    \item The tension in the string.
    \item The magnitude and direction of the reaction at the hinge.
    \item The work done by the tension if the rod is lowered until $B$ descends $0.02\,\mathrm{m}$.
\end{enumerate}
\end{exercice}

\begin{corrige}[Practice Problem]
\begin{enumerate}
    \item Geometry: $AB=4$, $AC=3$, so $BC=5$, $\sin\theta=3/5$, $\cos\theta=4/5$.
    Moments about $A$: $150\times2 + 100\times2 = T\sin\theta \times 4 \Rightarrow 500 = T\times\frac{3}{5}\times4 = \frac{12T}{5} \Rightarrow T = \frac{2500}{12} = \frac{625}{3} \approx 208.3\,\mathrm{N}$.
    \item $R_x = T\cos\theta = \frac{625}{3}\times\frac{4}{5} = \frac{500}{3} \approx 166.7\,\mathrm{N}$.
    $R_y = 150+100 - T\sin\theta = 250 - \frac{625}{3}\times\frac{3}{5} = 250 - 125 = 125\,\mathrm{N}$.
    $R = \sqrt{(500/3)^2 + 125^2} = \sqrt{250000/9 + 15625} = \sqrt{390625/9} = 625/3 \approx 208.3\,\mathrm{N}$.
    $\tan\alpha = 125/(500/3) = 3/4 \Rightarrow \alpha \approx 36.87^\circ$.
    \item $\delta \vec{r}_B = -0.02\,\vec{j}$. $\vec{T} = T(-\cos\theta\,\vec{i} + \sin\theta\,\vec{j}) = \frac{625}{3}(-\frac{4}{5}\vec{i} + \frac{3}{5}\vec{j})$.
    $\delta W = \vec{T}\cdot\delta\vec{r}_B = \frac{625}{3}\times\frac{3}{5}\times(-0.02) = 125\times(-0.02) = -2.5\,\mathrm{J}$.
\end{enumerate}
\end{corrige}

\newpage

\coursec{Methods and worked examples}

\coursec{Application Of Scalar And Vector Products}

\courssub{Vector Component of a Vector in a Given Direction}

\begin{definition}[Scalar and Vector Components]
Let $\vec{a}$ and $\vec{b}$ be two vectors, with $\vec{b}\neq\vec{0}$. The \cle{scalar component} (or \cle{projection}) of $\vec{a}$ onto the direction of $\vec{b}$ is the signed magnitude
\[
a_b = \vec{a}\cdot\hat{b} = \frac{\vec{a}\cdot\vec{b}}{|\vec{b}|},
\]
where $\hat{b}=\dfrac{\vec{b}}{|\vec{b}|}$ is the unit vector in the direction of $\vec{b}$. The \cle{vector component} of $\vec{a}$ parallel to $\vec{b}$ is
\[
\vec{a}_{\parallel} = a_b\,\hat{b} = \frac{\vec{a}\cdot\vec{b}}{|\vec{b}|^2}\,\vec{b}.
\]
The perpendicular component is $\vec{a}_{\perp} = \vec{a} - \vec{a}_{\parallel}$; it satisfies $\vec{a}_{\parallel}\cdot\vec{a}_{\perp}=0$ and $\vec{a}_{\parallel}+\vec{a}_{\perp}=\vec{a}$.
\end{definition}

\begin{propriete}[Geometric meaning]
If $\theta$ is the angle between $\vec{a}$ and $\vec{b}$, then $a_b = |\vec{a}|\cos\theta$. The sign of $a_b$ tells whether the projection points along $\vec{b}$ ($a_b>0$) or opposite to $\vec{b}$ ($a_b<0$). The magnitude of the perpendicular component is $|\vec{a}_{\perp}| = |\vec{a}|\sin\theta$.
\end{propriete}

\begin{methode}[Component of a vector along a given direction]
To find the vector component of $\vec{a}$ in the direction of $\vec{b}$:
\begin{enumerate}
\item Compute the unit vector $\hat{b} = \dfrac{\vec{b}}{|\vec{b}|}$.
\item Compute the scalar component (projection): $a_b = \vec{a}\cdot\hat{b} = \dfrac{\vec{a}\cdot\vec{b}}{|\vec{b}|}$.
\item Compute the vector component: $\vec{a}_{\parallel} = (\vec{a}\cdot\hat{b})\,\hat{b} = \dfrac{\vec{a}\cdot\vec{b}}{|\vec{b}|^2}\,\vec{b}$.
\item If the perpendicular component is needed: $\vec{a}_{\perp} = \vec{a} - \vec{a}_{\parallel}$.
\item Check: $\vec{a}_{\parallel}\cdot\vec{a}_{\perp} = 0$ and $\vec{a}_{\parallel}+\vec{a}_{\perp}=\vec{a}$.
\end{enumerate}
\textbf{Reflex:} if the scalar component is negative, the projection points \emph{opposite} to $\vec{b}$.
\end{methode}

\begin{popfigure}
\begin{tikzpicture}[scale=0.8, >=stealth]
\draw[->] (-1,0) -- (7,0) node[right] {$\vec{b}$ direction};
\draw[->] (0,-1) -- (0,4);
\draw[thick, popblue, ->] (0,0) -- (5,2.5) node[midway, above left] {$\vec{a}$};
\draw[thick, poporange, ->] (0,0) -- (5,0) node[midway, below] {$\vec{a}_{\parallel}$};
\draw[thick, popgreen, ->] (5,0) -- (5,2.5) node[midway, right] {$\vec{a}_{\perp}$};
\draw[dashed] (5,2.5) -- (5,0);
\draw (2.5,0) arc (0:26.5:2.5) node[midway, right] {$\theta$};
\node at (0,0) [below left] {$O$};
\end{tikzpicture}
\legende{Decomposition of $\vec{a}$ into components parallel and perpendicular to $\vec{b}$.}
\end{popfigure}

\begin{exoresolu}[Resolving a displacement along a road]
A surveyor in Bamenda moves from point $A$ to point $B$, with displacement $\vec{s} = \begin{pmatrix} 6 \\ 2 \\ -3 \end{pmatrix}$ km. A straight road runs in the direction $\vec{d} = \begin{pmatrix} 2 \\ -1 \\ 2 \end{pmatrix}$. Find (a) the scalar component of $\vec{s}$ along the road; (b) the vector component of $\vec{s}$ along the road; (c) the perpendicular distance from $B$ to the road direction through $A$.
\tcblower
\textbf{(a) Scalar component.} First $|\vec{d}| = \sqrt{4+1+4} = 3$, so $\hat{d} = \dfrac{1}{3}\begin{pmatrix} 2 \\ -1 \\ 2 \end{pmatrix}$.
\[
s_d = \vec{s}\cdot\hat{d} = \frac{1}{3}(12 - 2 - 6) = \frac{4}{3}\ \text{km}.
\]
\textbf{(b) Vector component.}
\[
\vec{s}_{\parallel} = \frac{4}{3}\hat{d} = \frac{4}{9}\begin{pmatrix} 2 \\ -1 \\ 2 \end{pmatrix} = \begin{pmatrix} 8/9 \\ -4/9 \\ 8/9 \end{pmatrix}\ \text{km}.
\]
\textbf{(c) Perpendicular component.}
\[
\vec{s}_{\perp} = \vec{s} - \vec{s}_{\parallel} = \begin{pmatrix} 6 - 8/9 \\ 2 + 4/9 \\ -3 - 8/9 \end{pmatrix} = \begin{pmatrix} 46/9 \\ 22/9 \\ -35/9 \end{pmatrix}.
\]
Its magnitude:
\[
|\vec{s}_{\perp}| = \frac{1}{9}\sqrt{46^2 + 22^2 + 35^2} = \frac{1}{9}\sqrt{2116 + 484 + 1225} = \frac{\sqrt{3825}}{9} = \frac{15\sqrt{17}}{9} = \frac{5\sqrt{17}}{3} \approx 6.87\ \text{km}.
\]
\textbf{Check:} $\vec{s}_{\parallel}\cdot\vec{s}_{\perp} = \dfrac{1}{81}(8\times 46 - 4\times 22 - 8\times 35) = \dfrac{368 - 88 - 280}{81} = 0$. \checkmark
\end{exoresolu}

\courssub{Work Done by a Constant Force}

\begin{definition}[Work Done by a Constant Force]
If a constant force $\vec{F}$ acts on a particle that undergoes a displacement $\vec{s}$, the \cle{work done} by the force is the scalar product
\[
W = \vec{F}\cdot\vec{s} = |\vec{F}|\,|\vec{s}|\cos\theta,
\]
where $\theta$ is the angle between $\vec{F}$ and $\vec{s}$. The SI unit is the joule (J): $1\ \text{J} = 1\ \text{N}\cdot\text{m}$.
\end{definition}

\begin{propriete}[Sign of work]
\begin{itemize}
\item $W>0$: the force has a component in the direction of displacement (force \emph{helps} the motion).
\item $W=0$: the force is perpendicular to the displacement (e.g.\ weight on horizontal ground, normal reaction).
\item $W<0$: the force has a component opposite to the displacement (force \emph{opposes} the motion, e.g.\ friction).
\end{itemize}
\end{propriete}

\begin{methode}[Computing work done]
\begin{enumerate}
\item Identify the force $\vec{F}$ (in newtons) and the displacement $\vec{s}$ (in metres) as vectors.
\item Work done: $W = \vec{F}\cdot\vec{s} = |\vec{F}||\vec{s}|\cos\theta$ (in joules).
\item If only magnitudes and the angle $\theta$ between force and displacement are known, use the trigonometric form directly.
\item Interpret the sign: $W>0$ (force helps motion), $W=0$ (force perpendicular to motion), $W<0$ (force opposes motion, e.g.\ friction).
\item For several forces, add the works algebraically, or sum the forces first: $W = \left(\sum \vec{F}_i\right)\cdot\vec{s}$.
\end{enumerate}
\textbf{Reflex:} weight does no work on horizontal ground; friction always does negative work.
\end{methode}

\begin{popfigure}
\begin{tikzpicture}[scale=0.8, >=stealth]
\draw[thick] (-1,0) -- (8,0);
\draw[fill=popgray] (1,0) rectangle (3,1);
\draw[popblue, thick, ->] (2,1) -- (4.5,2.5) node[midway, above left] {$\vec{F}=80\,\text{N}$};
\draw[poporange, thick, ->] (2,0.5) -- (6,0.5) node[midway, below] {$\vec{s}=25\,\text{m}$};
\draw[dashed] (2,1) -- (2,0.5);
\draw (2,0.5) arc (0:30:1.5) node[midway, right] {$30^{\circ}$};
\draw[popred, thick, ->] (3,0) -- (1,0) node[midway, below] {$\vec{F}_f=30\,\text{N}$};
\node at (2,1.2) [above] {crate};
\end{tikzpicture}
\legende{Force inclined at $30^{\circ}$ pulling a crate; friction opposes motion.}
\end{popfigure}

\begin{exoresolu}[Work done pulling a load]
A porter at the Douala market pulls a crate along horizontal ground with a rope inclined at $30^{\circ}$ to the horizontal, exerting a constant force of $80$ N. The crate moves $25$ m. (a) Find the work done by the porter. (b) A friction force of $30$ N opposes the motion; find the total work done on the crate. (c) Using unit vectors $\vec{i}$ (horizontal) and $\vec{j}$ (vertical), write $\vec{F}$ and $\vec{s}$ and verify part (a) with the scalar product.
\tcblower
\textbf{(a)} $W = |\vec{F}||\vec{s}|\cos\theta = 80 \times 25 \times \cos 30^{\circ} = 2000 \times \dfrac{\sqrt{3}}{2} = 1000\sqrt{3} \approx 1732$ J.
\textbf{(b)} Work done by friction: $W_f = 30 \times 25 \times \cos 180^{\circ} = -750$ J. Weight and normal reaction do no work (perpendicular to motion). Total work: $W_{\text{tot}} = 1000\sqrt{3} - 750 \approx 982$ J.
\textbf{(c)} $\vec{F} = 80\cos 30^{\circ}\,\vec{i} + 80\sin 30^{\circ}\,\vec{j} = 40\sqrt{3}\,\vec{i} + 40\,\vec{j}$ N; $\vec{s} = 25\,\vec{i}$ m.
\[
W = \vec{F}\cdot\vec{s} = 40\sqrt{3}\times 25 + 40\times 0 = 1000\sqrt{3}\ \text{J}. \checkmark
\]
\end{exoresolu}

\begin{exoresolu}[Work done in three dimensions]
A constant force $\vec{F} = (2\vec{i} - 3\vec{j} + 6\vec{k})$ N moves a particle from $P(1, 0, 2)$ to $Q(4, 5, -1)$, coordinates in metres. Find the work done.
\tcblower
Displacement: $\vec{s} = \overrightarrow{PQ} = \begin{pmatrix} 4-1 \\ 5-0 \\ -1-2 \end{pmatrix} = \begin{pmatrix} 3 \\ 5 \\ -3 \end{pmatrix}$ m.
\[
W = \vec{F}\cdot\vec{s} = (2)(3) + (-3)(5) + (6)(-3) = 6 - 15 - 18 = -27\ \text{J}.
\]
The negative sign means the force globally opposes the displacement: $27$ J of work must be supplied \emph{against} the force.
\end{exoresolu}

\courssub{Moment of a Force and Equilibrium}

\begin{definition}[Moment of a Force about a Point]
The \cle{moment} of a force $\vec{F}$ about a point $O$ is the vector
\[
\vec{M}_O = \vec{r} \times \vec{F},
\]
where $\vec{r}$ is the position vector from $O$ to \emph{any} point on the line of action of $\vec{F}$. Its magnitude is $M = Fd$, where $d$ is the perpendicular distance from $O$ to the line of action. In two dimensions, we use the scalar moment $M = Fd$ with a sign convention (usually anticlockwise positive).
\end{definition}

\begin{propriete}[Equilibrium of a Rigid Body]
A rigid body is in equilibrium under a system of forces if and only if
\[
\sum \vec{F} = \vec{0} \quad\text{and}\quad \sum \vec{M}_O = \vec{0} \text{ about any point } O.
\]
In two dimensions this gives three independent scalar equations: $\sum F_x=0$, $\sum F_y=0$, $\sum M=0$.
\end{propriete}

\begin{methode}[Moment about a point; equilibrium]
\begin{enumerate}
\item The moment of $\vec{F}$ about point $O$ is $\vec{M}_O = \vec{r} \times \vec{F}$, where $\vec{r}$ is the position vector from $O$ to any point of the line of action.
\item In 2D, use the scalar moment: $M = Fd$, where $d$ is the \cle{perpendicular distance} from $O$ to the line of action; fix a positive sense (usually anticlockwise).
\item For \cle{equilibrium} of a rigid body: $\sum \vec{F} = \vec{0}$ \textbf{and} $\sum \vec{M}_O = \vec{0}$ about \emph{any} point $O$.
\item Strategy: take moments about a point where unknown forces act, to eliminate them.
\item In 3D, compute $\vec{r}\times\vec{F}$ with the determinant formula and set each component of the total moment to zero.
\end{enumerate}
\textbf{Reflex:} a force whose line of action passes through $O$ has zero moment about $O$.
\end{methode}

\begin{popfigure}
\begin{tikzpicture}[scale=0.7, >=stealth]
\draw[thick] (0,0) -- (6,0);
\draw[fill=popgray] (0,-0.2) rectangle (0.4,0.2) node[midway, below] {$A$};
\draw[fill=popgray] (4,-0.2) rectangle (4.4,0.2) node[midway, below] {$C$};
\draw[->, thick] (3,0) -- (3,-1) node[midway, right] {$120\,\text{N}$};
\draw[->, thick] (6,0) -- (6,-1) node[midway, right] {$200\,\text{N}$};
\draw[->, thick, popblue] (0.2,-0.2) -- (0.2,1) node[midway, left] {$R_A$};
\draw[->, thick, popblue] (4.2,-0.2) -- (4.2,1) node[midway, left] {$R_C$};
\draw[<->] (0.2,1.2) -- (4.2,1.2) node[midway, above] {$4\,\text{m}$};
\draw[<->] (4.2,1.2) -- (6.2,1.2) node[midway, above] {$2\,\text{m}$};
\draw[<->] (0.2,1.5) -- (3.2,1.5) node[midway, above] {$3\,\text{m}$};
\end{tikzpicture}
\legende{Beam $AB$ with supports at $A$ and $C$, weight at midpoint, load at $B$.}
\end{popfigure}

\begin{exoresolu}[Equilibrium of a uniform beam]
A uniform beam $AB$ of length $6$ m and weight $120$ N rests horizontally on two supports at $A$ and at $C$, where $AC = 4$ m. A load of $200$ N hangs at $B$. Find the reactions at the supports.
\tcblower
Let $R_A$ and $R_C$ be the vertical reactions. The weight acts at the midpoint, $3$ m from $A$.
\textbf{Vertical equilibrium:} $R_A + R_C = 120 + 200 = 320$ N.
\textbf{Moments about $A$} (anticlockwise positive):
\[
R_C \times 4 - 120 \times 3 - 200 \times 6 = 0 \implies 4R_C = 360 + 1200 = 1560 \implies R_C = 390\ \text{N}.
\]
Then $R_A = 320 - 390 = -70$ N.
\textbf{Interpretation:} $R_A < 0$ means the support at $A$ would have to pull \emph{down}: the beam tips about $C$ unless it is bolted at $A$. The load is too far beyond $C$ for a simple resting contact at $A$.
\end{exoresolu}

\begin{popfigure}
\begin{tikzpicture}[scale=0.8, >=stealth]
\draw[->] (-1,0,0) -- (3,0,0) node[right] {$x$};
\draw[->] (0,-1,0) -- (0,3,0) node[above] {$y$};
\draw[->] (0,0,-1) -- (0,0,3) node[above] {$z$};
\draw[thick, popblue, ->] (0,0,0) -- (1,2,0) node[midway, above left] {$\vec{r}$};
\draw[thick, poporange, ->] (1,2,0) -- (4,4,-1) node[midway, above] {$\vec{F}$};
\draw[thick, popgreen, ->] (0,0,0) -- (-0.5,0.5,-1) node[midway, left] {$\vec{M}_O$};
\draw[dashed] (1,2,0) -- (1,0,0) -- (0,0,0);
\draw[dashed] (1,2,0) -- (0,2,0);
\node at (0,0,0) [below left] {$O$};
\node at (1,2,0) [above right] {$A$};
\end{tikzpicture}
\legende{Moment $\vec{M}_O = \vec{r}\times\vec{F}$ about the origin.}
\end{popfigure}

\begin{exoresolu}[Moment as a vector product]
A force $\vec{F} = (3\vec{i} + 2\vec{j} - \vec{k})$ N acts at the point $A(1, 2, 0)$, coordinates in metres. Find the moment of $\vec{F}$ about the origin $O$, and its magnitude.
\tcblower
$\vec{r} = \overrightarrow{OA} = \vec{i} + 2\vec{j}$.
\[
\vec{M}_O = \vec{r}\times\vec{F} = \begin{vmatrix} \vec{i} & \vec{j} & \vec{k} \\ 1 & 2 & 0 \\ 3 & 2 & -1 \end{vmatrix} = \vec{i}(2\times(-1) - 0\times 2) - \vec{j}(1\times(-1) - 0\times 3) + \vec{k}(1\times 2 - 2\times 3)
\]
\[
= -2\vec{i} + \vec{j} - 4\vec{k}\ \text{N\,m}.
\]
Magnitude: $|\vec{M}_O| = \sqrt{4 + 1 + 16} = \sqrt{21} \approx 4.58$ N\,m.
\end{exoresolu}

\courssub{Single Forces and Couples}

\begin{definition}[Couple]
A \cle{couple} consists of two forces equal in magnitude, opposite in direction, parallel but not collinear. The resultant force is zero, but the moment (torque) is non-zero and \emph{independent of the reference point}.
\end{definition}

\begin{propriete}[Moment of a couple]
The moment of a couple is $M = Fd$, where $F$ is the magnitude of one force and $d$ is the perpendicular distance between their lines of action. The moment vector $\vec{G}$ is the same about every point. A couple cannot be replaced by a single force.
\end{propriete}

\begin{methode}[Resultant force and couples]
\begin{enumerate}
\item A \cle{couple} is a pair of equal, opposite, parallel forces not in the same line: resultant force zero, but a non-zero moment.
\item Moment of a couple: $M = Fd$, where $d$ is the perpendicular distance between the two lines of action. It is the \emph{same about every point}.
\item To reduce a system to a single force at a point $O$: resultant $\vec{R} = \sum \vec{F}_i$ plus a couple of moment $\vec{G}_O = \sum \vec{M}_O(\vec{F}_i)$.
\item A single force $\vec{F}$ acting at $A$ is equivalent to the same force at $O$ \emph{plus} a couple of moment $\overrightarrow{OA}\times\vec{F}$.
\item Two systems are equivalent if and only if they have the same resultant force \emph{and} the same total moment about one (hence any) point.
\end{enumerate}
\textbf{Reflex:} a couple can never be replaced by a single force.
\end{methode}

\begin{popfigure}
\begin{tikzpicture}[scale=0.8, >=stealth]
\draw[thick] (0,0) circle (2);
\draw[popblue, thick, ->] (0,2) -- (0,3) node[above] {$50\,\text{N}$};
\draw[popblue, thick, ->] (0,-2) -- (0,-3) node[below] {$50\,\text{N}$};
\draw[dashed] (0,2) -- (0,-2);
\node at (0,0) {$O$};
\draw[<->] (2.2,2) -- (2.2,-2) node[midway, right] {$d=0.40\,\text{m}$};
\draw[->, thick] (1.5,1.5) arc (45:135:1.5) node[midway, above] {$M=20\,\text{N\,m}$};
\end{tikzpicture}
\legende{Couple on a steering wheel: equal opposite forces, moment $Fd$.}
\end{popfigure}

\begin{exoresolu}[Identifying a couple]
Two horizontal forces of $50$ N act on a steering wheel of diameter $40$ cm: one pushes north at the top of the wheel, the other pushes south at the bottom. (a) Show the resultant force is zero. (b) Find the moment of the couple. (c) Verify the moment is the same about the centre $O$ and about the top point $T$.
\tcblower
\textbf{(a)} The forces are equal ($50$ N), parallel and opposite: $\vec{R} = 50 - 50 = 0$ N.
\textbf{(b)} $M = Fd = 50 \times 0.40 = 20$ N\,m.
\textbf{(c)} About $O$: each force has lever arm $0.20$ m and both turn the wheel the same way: $M_O = 50(0.20) + 50(0.20) = 20$ N\,m. About $T$: the force at $T$ has zero moment; the force at the bottom has lever arm $0.40$ m: $M_T = 50 \times 0.40 = 20$ N\,m. Same value, as expected for a couple. \checkmark
\end{exoresolu}

\begin{exoresolu}[Replacing a force by a force and a couple]
A force $\vec{F} = (4\vec{i} - 3\vec{j})$ N acts at $A(2, 1)$. Replace it by an equal force through the origin $O$ together with a couple.
\tcblower
The force through $O$ is $\vec{F} = 4\vec{i} - 3\vec{j}$ N. The couple has moment
\[
\vec{G} = \overrightarrow{OA}\times\vec{F} = \begin{vmatrix} \vec{i} & \vec{j} & \vec{k} \\ 2 & 1 & 0 \\ 4 & -3 & 0 \end{vmatrix} = \vec{k}(2\times(-3) - 1\times 4) = -10\,\vec{k}\ \text{N\,m}.
\]
So the system is: force $(4\vec{i} - 3\vec{j})$ N at $O$, plus a clockwise couple of moment $10$ N\,m.
\end{exoresolu}

\courssub{Coplanar Forces and Equilibrium}

\begin{methode}[Solving coplanar equilibrium problems]
\begin{enumerate}
\item Draw a clear force diagram: mark \emph{all} forces (weights, reactions, tensions, friction).
\item Choose perpendicular axes (often horizontal/vertical or along/perpendicular to a plane) and resolve: $\sum F_x = 0$, $\sum F_y = 0$.
\item Take moments about a well-chosen point: $\sum M = 0$. Three independent equations for a rigid body in a plane.
\item For three forces only in equilibrium: they are concurrent or parallel — use the \cle{triangle of forces} or Lami's theorem: $\dfrac{F_1}{\sin\alpha_1} = \dfrac{F_2}{\sin\alpha_2} = \dfrac{F_3}{\sin\alpha_3}$, where $\alpha_i$ is the angle between the other two forces.
\item Solve the simultaneous equations and check signs physically.
\end{enumerate}
\textbf{Reflex:} with exactly three forces, Lami's theorem is usually the fastest route.
\end{methode}

\begin{popfigure}
\begin{tikzpicture}[scale=0.7, >=stealth]
\draw[thick] (0,0) -- (0,5) node[above] {wall};
\draw[thick] (0,0) -- (5,0) node[right] {ground};
\draw[thick, popblue] (0,0) -- (4,3) node[midway, above right] {ladder $5\,\text{m}$};
\draw[->, thick] (2,1.5) -- (2,0.5) node[midway, right] {$W=250\,\text{N}$};
\draw[->, thick, poporange] (4,3) -- (5,3) node[midway, above] {$R$};
\draw[->, thick, poporange] (0,0) -- (0,1) node[midway, left] {$N$};
\draw[->, thick, poporange] (0,0) -- (1,0) node[midway, below] {$F$};
\draw (1,0) arc (0:36.87:1) node[midway, right] {$60^{\circ}$};
\draw[dashed] (2,1.5) -- (2,0);
\draw[dashed] (4,3) -- (4,0);
\end{tikzpicture}
\legende{Ladder against smooth wall, rough ground.}
\end{popfigure}

\begin{exoresolu}[Ladder against a wall]
A uniform ladder of weight $W = 250$ N and length $5$ m leans against a smooth vertical wall, its foot on rough horizontal ground, making $60^{\circ}$ with the ground. Find the reaction of the wall, and the vertical and frictional components of the ground reaction.
\tcblower
Let $R$ = wall reaction (horizontal, wall smooth), $N$ = normal ground reaction, $F$ = friction at the ground. The weight acts at the midpoint.
\textbf{Resolve horizontally:} $F - R = 0 \implies F = R$.
\textbf{Resolve vertically:} $N - 250 = 0 \implies N = 250$ N.
\textbf{Moments about the foot} (lever arms: weight $\to 2.5\cos 60^{\circ}$; wall reaction $\to 5\sin 60^{\circ}$):
\[
R \times 5\sin 60^{\circ} - 250 \times 2.5\cos 60^{\circ} = 0
\]
\[
R = \frac{250 \times 2.5 \times 0.5}{5 \times \frac{\sqrt{3}}{2}} = \frac{312.5}{4.330} = \frac{125}{2\sqrt{3}} = \frac{125\sqrt{3}}{6} \approx 36.1\ \text{N}.
\]
Hence $F \approx 36.1$ N, $N = 250$ N, $R \approx 36.1$ N.
\end{exoresolu}

\begin{popfigure}
\begin{tikzpicture}[scale=0.8, >=stealth]
\draw[thick] (-3,0) -- (3,0) node[right] {ceiling};
\draw[popblue, thick, ->] (0,0) -- (-2,-1.5) node[midway, left] {$T_1$};
\draw[popblue, thick, ->] (0,0) -- (2,-1.5) node[midway, right] {$T_2$};
\draw[poporange, thick, ->] (0,0) -- (0,-2) node[midway, right] {$W=40\,\text{N}$};
\draw (-1.5,0) arc (180:130:1.5) node[midway, left] {$50^{\circ}$};
\draw (1.5,0) arc (0:50:1.5) node[midway, right] {$70^{\circ}$};
\draw (0,-0.5) arc (-90:-150:0.5) node[midway, left] {$140^{\circ}$};
\draw (0,-0.5) arc (-90:-30:0.5) node[midway, right] {$160^{\circ}$};
\node at (0,0) [above] {knot};
\end{tikzpicture}
\legende{Three forces in equilibrium at the knot: Lami's theorem applies.}
\end{popfigure}

\begin{exoresolu}[Three forces: Lami's theorem]
A picture of weight $40$ N hangs by two strings from two hooks. The strings make angles $50^{\circ}$ and $70^{\circ}$ with the horizontal ceiling. Find the tension in each string.
\tcblower
Three forces in equilibrium: $T_1$, $T_2$, $W = 40$ N. The angles between the forces at the knot: between $T_1$ and $T_2$: $180^{\circ} - 50^{\circ} - 70^{\circ} = 60^{\circ}$; between $T_1$ and the vertical weight: $90^{\circ} + 50^{\circ} = 140^{\circ}$; between $T_2$ and the weight: $90^{\circ} + 70^{\circ} = 160^{\circ}$.
Lami's theorem:
\[
\frac{T_1}{\sin 160^{\circ}} = \frac{T_2}{\sin 140^{\circ}} = \frac{40}{\sin 60^{\circ}}.
\]
\[
T_1 = \frac{40\sin 160^{\circ}}{\sin 60^{\circ}} = \frac{40 \times 0.3420}{0.8660} \approx 15.8\ \text{N}, \qquad T_2 = \frac{40\sin 140^{\circ}}{\sin 60^{\circ}} = \frac{40 \times 0.6428}{0.8660} \approx 29.7\ \text{N}.
\]
\textbf{Check (horizontal):} $T_2\cos 70^{\circ} \approx 29.7\times 0.3420 \approx 10.2$ N and $T_1\cos 50^{\circ} \approx 15.8\times 0.6428 \approx 10.2$ N. \checkmark
\end{exoresolu}

\courssub{Simple Systems of Forces in Three Dimensions}

\begin{methode}[Equilibrium in 3D]
\begin{enumerate}
\item Express every force in component form $\vec{F}_i = F_{ix}\vec{i} + F_{iy}\vec{j} + F_{iz}\vec{k}$; for a force along a line $AB$, use $\vec{F} = F\,\dfrac{\overrightarrow{AB}}{|\overrightarrow{AB}|}$.
\item Resultant: $\vec{R} = \sum \vec{F}_i$. Equilibrium of a particle: $\vec{R} = \vec{0}$, i.e.\ three scalar equations.
\item For a rigid body, add the moment condition: $\sum \vec{r}_i \times \vec{F}_i = \vec{0}$ (three more scalar equations).
\item Solve the system of up to six scalar equations for the unknowns.
\end{enumerate}
\textbf{Reflex:} write each force as magnitude $\times$ unit vector along its line — this is where most marks are won or lost.
\end{methode}

\begin{popfigure}
\begin{tikzpicture}[scale=0.7, >=stealth]
\draw[->] (-3,0,0) -- (3,0,0) node[right] {$x$};
\draw[->] (0,-3,0) -- (0,3,0) node[above] {$y$};
\draw[->] (0,0,-1) -- (0,0,4) node[above] {$z$};
\draw[thick, popblue, ->] (0,0,0) -- (2,0,3) node[midway, above right] {$T_A$};
\draw[thick, poporange, ->] (0,0,0) -- (-2,1,3) node[midway, above left] {$T_B$};
\draw[thick, popgreen, ->] (0,0,0) -- (0,-2,3) node[midway, right] {$T_C$};
\draw[thick, popred, ->] (0,0,0) -- (0,0,-1.5) node[midway, left] {$W=60\,\text{N}$};
\node at (0,0,0) [below left] {$P$};
\node at (2,0,3) [right] {$A$};
\node at (-2,1,3) [left] {$B$};
\node at (0,-2,3) [right] {$C$};
\end{tikzpicture}
\legende{Three cables supporting a lamp at $P(0,0,0)$; anchors at $z=3$ m.}
\end{popfigure}

\begin{exoresolu}[Particle held by three cables]
A lamp of weight $60$ N hangs at point $P(0,0,0)$ supported by three cables anchored at $A(2,0,3)$, $B(-2,1,3)$ and $C(0,-2,3)$ (units: metres). Find the tension in each cable, given symmetry is broken and the tensions are $T_A, T_B, T_C$.
\tcblower
Unit vectors from $P$ to each anchor:
\[
\overrightarrow{PA} = \begin{pmatrix} 2\\0\\3 \end{pmatrix},\ |\overrightarrow{PA}| = \sqrt{13}; \quad \overrightarrow{PB} = \begin{pmatrix} -2\\1\\3 \end{pmatrix},\ |\overrightarrow{PB}| = \sqrt{14}; \quad \overrightarrow{PC} = \begin{pmatrix} 0\\-2\\3 \end{pmatrix},\ |\overrightarrow{PC}| = \sqrt{13}.
\]
Equilibrium $\sum \vec{F} = \vec{0}$ gives:
\begin{align*}
x:&\quad \frac{2T_A}{\sqrt{13}} - \frac{2T_B}{\sqrt{14}} = 0,\\
y:&\quad \frac{T_B}{\sqrt{14}} - \frac{2T_C}{\sqrt{13}} = 0,\\
z:&\quad \frac{3T_A}{\sqrt{13}} + \frac{3T_B}{\sqrt{14}} + \frac{3T_C}{\sqrt{13}} = 60.
\end{align*}
From the $x$-equation: $T_B = \dfrac{\sqrt{14}}{\sqrt{13}}T_A$. From the $y$-equation: $T_C = \dfrac{\sqrt{13}}{2\sqrt{14}}T_B = \dfrac{T_A}{2}$.
Substitute into the $z$-equation:
\[
\frac{3T_A}{\sqrt{13}} + \frac{3T_A}{\sqrt{13}} + \frac{3T_A}{2\sqrt{13}} = 60 \implies \frac{3T_A}{\sqrt{13}}\left(1 + 1 + \tfrac{1}{2}\right) = 60 \implies \frac{15T_A}{2\sqrt{13}} = 60.
\]
\[
T_A = \frac{120\sqrt{13}}{15} = 8\sqrt{13} \approx 28.8\ \text{N}, \quad T_C = 4\sqrt{13} \approx 14.4\ \text{N}, \quad T_B = 8\sqrt{14} \approx 29.9\ \text{N}.
\]
\textbf{Check ($z$):} $\dfrac{3}{\sqrt{13}}(8\sqrt{13} + 4\sqrt{13}) + \dfrac{3}{\sqrt{14}}(8\sqrt{14}) = 36 + 24 = 60$ N. \checkmark
\end{exoresolu}

\begin{aretenir}[Key formulas for the chapter]
\begin{itemize}
\item Scalar component of $\vec{a}$ on $\vec{b}$: $a_b = \dfrac{\vec{a}\cdot\vec{b}}{|\vec{b}|}$.
\item Vector component: $\vec{a}_{\parallel} = \dfrac{\vec{a}\cdot\vec{b}}{|\vec{b}|^2}\,\vec{b}$; perpendicular: $\vec{a}_{\perp} = \vec{a} - \vec{a}_{\parallel}$.
\item Work: $W = \vec{F}\cdot\vec{s} = Fs\cos\theta$ (joules).
\item Moment about $O$: $\vec{M}_O = \vec{r}\times\vec{F}$; magnitude $M = Fd$ ($d$ = perpendicular distance).
\item Equilibrium: $\sum\vec{F}=\vec{0}$, $\sum\vec{M}_O=\vec{0}$.
\item Couple moment: $M = Fd$ (same about any point).
\item Lami's theorem (3 forces): $\dfrac{F_1}{\sin\alpha_1} = \dfrac{F_2}{\sin\alpha_2} = \dfrac{F_3}{\sin\alpha_3}$.
\item 3D force along $AB$: $\vec{F} = F\,\dfrac{\overrightarrow{AB}}{|\overrightarrow{AB}|}$.
\end{itemize}
\end{aretenir}

\begin{savaistu}[From Cameroon to the stars]
The principles of moments and equilibrium are used daily in Cameroonian engineering: from designing the supports of the \emph{Pont de la Réunification} in Douala to calculating the tension in the guy-wires of telecommunication masts on Mount Cameroon. In space, the same vector methods control the orientation of satellites — reaction wheels produce couples to turn the spacecraft without using fuel. The scalar product giving work is the foundation of the work–energy theorem, which explains why a car climbing the steep roads of Buea needs more fuel than on the flat Littoral plains.
\end{savaistu}

\begin{attention}[Frequent errors at the GCE]
\begin{itemize}
\item Confusing the \emph{scalar} component (a number) with the \emph{vector} component (a vector).
\item Using $\sin\theta$ instead of $\cos\theta$ in $W = Fs\cos\theta$ — $\theta$ is the angle \emph{between} force and displacement.
\item Taking moments with the distance along the beam instead of the \emph{perpendicular} distance to the line of action.
\item Forgetting that the moment of a couple is independent of the point chosen.
\item In 3D, forgetting to divide by the length when forming unit vectors along cables.
\item Misidentifying the angles in Lami's theorem: $\alpha_i$ is the angle between the \emph{other two} forces, not the angle the force makes with the horizontal.
\item Writing $\vec{r}\times\vec{F}$ with $\vec{r}$ from the force to the point instead of from the point to the force.
\end{itemize}
\end{attention}

\begin{exercice}[Practice problems]
\begin{enumerate}
\item Find the vector component of $\vec{a}=3\vec{i}-4\vec{j}+12\vec{k}$ parallel to $\vec{b}=2\vec{i}+\vec{j}-2\vec{k}$, and the perpendicular distance from the tip of $\vec{a}$ to the line of $\vec{b}$.
\item A force $\vec{F}=5\vec{i}+12\vec{j}$ N moves an object from $(1,2)$ to $(4,6)$ (metres). Find the work done. If a friction force $\vec{F}_f=-3\vec{i}-4\vec{j}$ N acts, find the net work.
\item A uniform rod $AB$ of weight $W$ and length $L$ rests with $A$ on rough horizontal ground and $B$ against a smooth vertical wall. The rod makes angle $\theta$ with the ground. Show that the minimum coefficient of friction for equilibrium is $\mu_{\min}=\frac{1}{2}\cot\theta$.
\item Two forces $\vec{F}_1=10\vec{j}$ N and $\vec{F}_2=-10\vec{j}$ N act at $(0,0,0)$ and $(0,0,0.5)$ respectively. Find the moment of this couple about the origin and about the point $(1,2,3)$.
\item A sign of weight $200$ N hangs from a horizontal beam $AB$ of length $2$ m, hinged at $A$ to a vertical wall. A cable from $B$ to a point $C$ on the wall $1.5$ m above $A$ holds the beam horizontal. Find the tension in the cable and the reaction at the hinge.
\item Four forces act on a particle at the origin: $\vec{F}_1=20\vec{i}$ N, $\vec{F}_2=15\vec{j}$ N, $\vec{F}_3$ along the line to $(3,4,0)$ with magnitude $25$ N, and $\vec{F}_4$ along the line to $(-1,2,2)$ with magnitude $3$ N. Is the particle in equilibrium? If not, find the additional force needed.
\end{enumerate}
\end{exercice}

\begin{corrige}[Exercise 1]
$\vec{a}\cdot\vec{b}=6-4-24=-22$, $|\vec{b}|=3$, $\hat{b}=\frac{1}{3}(2,1,-2)$. Scalar component $a_b=-22/3$. Vector component $\vec{a}_{\parallel}=\frac{-22}{9}(2,1,-2)=\left(-\frac{44}{9},-\frac{22}{9},\frac{44}{9}\right)$. $\vec{a}_{\perp}=\vec{a}-\vec{a}_{\parallel}=\left(\frac{71}{9},-\frac{14}{9},\frac{64}{9}\right)$. Distance $|\vec{a}_{\perp}|=\frac{1}{9}\sqrt{71^2+14^2+64^2}=\frac{\sqrt{9213}}{9}\approx 10.66$.
\end{corrige}

\begin{corrige}[Exercise 2]
$\vec{s}=(3,4)$, $W=\vec{F}\cdot\vec{s}=15+48=63$ J. $\vec{F}_{\text{net}}=(2,8)$, $W_{\text{net}}=6+32=38$ J. Alternatively $W_f=\vec{F}_f\cdot\vec{s}=-9-16=-25$ J, net $=63-25=38$ J.
\end{corrige}

\begin{corrige}[Exercise 3]
Forces: $N_A$ (up), $F_A$ (right), $N_B$ (left), $W$ (down at centre). Horizontal: $F_A=N_B$. Vertical: $N_A=W$. Moments about $A$: $N_B L\sin\theta - W\frac{L}{2}\cos\theta=0 \Rightarrow N_B=\frac{W}{2}\cot\theta$. Then $F_A=\frac{W}{2}\cot\theta$. Friction limit: $F_A\leq\mu N_A \Rightarrow \frac{W}{2}\cot\theta\leq\mu W \Rightarrow \mu\geq\frac{1}{2}\cot\theta$.
\end{corrige}

\begin{corrige}[Exercise 4]
Couple moment $\vec{G}=\vec{r}\times\vec{F}$ where $\vec{r}$ is from any point on $\vec{F}_1$ to any point on $\vec{F}_2$. Take $\vec{r}=(0,0,0.5)$. $\vec{G}=(0,0,0.5)\times(0,10,0)=(-5,0,0)$ N\,m. Magnitude $5$ N\,m. About $(1,2,3)$: $\vec{r}_1=(-1,-2,-3)$, $\vec{r}_2=(-1,-2,-2.5)$. $\vec{M}=\vec{r}_1\times\vec{F}_1+\vec{r}_2\times\vec{F}_2 = (-1,-2,-3)\times(0,10,0)+(-1,-2,-2.5)\times(0,-10,0) = (-30,0,-10)+( -25,0,-10) = (-5,0,0)$ N\,m. Same.
\end{corrige}

\begin{corrige}[Exercise 5]
Geometry: $AC=1.5$, $AB=2$, so $\sin\theta=1.5/\sqrt{2^2+1.5^2}=1.5/2.5=0.6$, $\cos\theta=0.8$, $\theta\approx36.87^\circ$. Forces: $T$ along $BC$, $R_A$ at hinge, $W=200$ N down at centre ($1$ m from $A$). Moments about $A$: $T\times1.5 - 200\times1 =0 \Rightarrow T=400/3\approx133.3$ N. Horizontal: $R_{Ax}=T\cos\theta=133.3\times0.8=106.7$ N. Vertical: $R_{Ay}+T\sin\theta=200 \Rightarrow R_{Ay}=200-133.3\times0.6=120$ N. $R_A=\sqrt{106.7^2+120^2}\approx160.5$ N at $\tan^{-1}(120/106.7)\approx48.4^\circ$ above horizontal.
\end{corrige}

\begin{corrige}[Exercise 6]
$\vec{F}_3=25\frac{(3,4,0)}{5}=(15,20,0)$. $\vec{F}_4=3\frac{(-1,2,2)}{3}=(-1,2,2)$. Sum: $\vec{R}=(20+15-1, 0+20+2, 0+0+2)=(34,22,2)\neq\vec{0}$. Not in equilibrium. Additional force needed: $\vec{F}_5=(-34,-22,-2)$ N.
\end{corrige}

\newpage

\coursec{Exercises}

\courssub{I check my knowledge}

\begin{exercice}[Scalar and vector components]
Given $\vec{a} = 3\mathbf{i} - 2\mathbf{j} + \mathbf{k}$ and $\vec{b} = \mathbf{i} + 4\mathbf{j} - 2\mathbf{k}$, find the scalar component of $\vec{a}$ in the direction of $\vec{b}$, the vector component of $\vec{a}$ in the direction of $\vec{b}$, and the angle between $\vec{a}$ and $\vec{b}$.
\end{exercice}

\begin{exercice}[Work done by a constant force]
A force $\vec{F} = (5\mathbf{i} + 2\mathbf{j})\,\mathrm{N}$ moves a crate from point $A(1,3)$ to point $B(7,9)$. Calculate the work done by $\vec{F}$. If a friction force of magnitude $4\,\mathrm{N}$ opposes the motion, find the total work done on the crate.
\end{exercice}

\begin{exercice}[Moment of a force]
A force $\vec{F} = 4\mathbf{i} - 3\mathbf{j}\,\mathrm{N}$ acts at the point $A(2,5)$. Replace it by an equal force through the origin together with a couple, and state the moment of the couple.
\end{exercice}

\begin{exercice}[Equilibrium of coplanar forces]
Three coplanar forces $P$, $Q$, $R$ act at a point and keep it in equilibrium. $P = 50\,\mathrm{N}$ acts along the positive $x$-axis and $Q = 40\,\mathrm{N}$ acts at $120^\circ$ to $P$. Find the magnitude and direction of $R$ using components, then verify your result with Lami's theorem.
\end{exercice}

\courssub{I practise}

\begin{exercice}[Work on an inclined plane]
A man pulls a $50\,\mathrm{kg}$ load up a plane inclined at $30^\circ$ to the horizontal with a rope parallel to the plane, over a distance of $8\,\mathrm{m}$. Find the work done against gravity and the work done by the man if friction is neglected. (Take $g = 9.8\,\mathrm{m\,s^{-2}}$.)
\end{exercice}

\begin{exercice}[Reactions on a beam]
A uniform beam $AB$ of length $6\,\mathrm{m}$ and weight $120\,\mathrm{N}$ rests on supports at $A$ and at $C$, where $C$ is $4\,\mathrm{m}$ from $A$. A load of $300\,\mathrm{N}$ hangs at $B$. Find the reactions at the supports by taking moments.
\end{exercice}

\begin{exercice}[Moment of a couple]
Two forces of $20\,\mathrm{N}$ act at the ends of a horizontal bar of length $0.5\,\mathrm{m}$, in opposite directions and perpendicular to the bar. Find the moment of the couple, and show that it is the same about any point of the bar.
\end{exercice}

\begin{exercice}[3D equilibrium of a particle]
A particle is held in equilibrium by three cables attached to points $A(2,0,0)$, $B(0,3,0)$ and $C(0,0,4)$, with the particle at $P(1,1,1)$. Given that the tensions in cables $PA$ and $PB$ are $T_A = 50\,\mathrm{N}$ and $T_B = 40\,\mathrm{N}$, set up the vector equilibrium equations and find the tension in cable $PC$.
\end{exercice}

\courssub{I prepare for the GCE}

\begin{exercice}[Ladder against a smooth wall] \textit{(12 marks)}
A uniform ladder of weight $W$ and length $L$ leans against a smooth vertical wall with its foot on rough horizontal ground. The ladder makes an angle $\theta$ with the horizontal. The coefficient of friction between the ladder and the ground is $\mu$. By resolving forces and taking moments, find the reaction at the wall and the friction force at the ground for equilibrium. Determine the minimum value of $\mu$ for which the ladder does not slip.
\end{exercice}

\begin{exercice}[Vector moment in 3D] \textit{(10 marks)}
A force $\vec{F} = 2\mathbf{i} + \mathbf{j} - 3\mathbf{k}\,\mathrm{N}$ acts at the point with position vector $\vec{r} = \mathbf{i} - 2\mathbf{j} + \mathbf{k}\,\mathrm{m}$ relative to the origin $O$. Find the vector moment of $\vec{F}$ about $O$ and its magnitude. Also find the perpendicular distance from $O$ to the line of action of $\vec{F}$.
\end{exercice}

\begin{exercice}[System of forces in 3D] \textit{(13 marks)}
A rigid body is subjected to the following forces:
\begin{itemize}
\item $\vec{F}_1 = 3\mathbf{i} - 2\mathbf{j} + 4\mathbf{k}\,\mathrm{N}$ acting at point $A(1,0,2)$,
\item $\vec{F}_2 = -2\mathbf{i} + 5\mathbf{j} - 3\mathbf{k}\,\mathrm{N}$ acting at point $B(0,2,1)$,
\item $\vec{F}_3 = \mathbf{i} + 3\mathbf{j} - 2\mathbf{k}\,\mathrm{N}$ acting at point $C(2,1,0)$.
\end{itemize}
Reduce this system to a single force at the origin together with a couple. Find the resultant force, the resultant couple, and state whether the system can be reduced to a single force. If so, find the line of action of that single force.
\end{exercice}

\newpage

\coursec{Solutions to the exercises}

\coursec{Application Of Scalar And Vector Products}

\courssub{I check my knowledge}

\begin{definition}[Scalar and Vector Components]
Let $\vec{a}$ and $\vec{b}$ be two vectors with $\vec{b} \neq \vec{0}$. The \cle{scalar component} of $\vec{a}$ in the direction of $\vec{b}$ is the signed magnitude of the orthogonal projection of $\vec{a}$ onto the line spanned by $\vec{b}$:
\[
\operatorname{comp}_{\vec{b}}\vec{a} = \frac{\vec{a}\cdot\vec{b}}{|\vec{b}|}.
\]
The \cle{vector component} (or orthogonal projection) of $\vec{a}$ onto $\vec{b}$ is the vector
\[
\operatorname{proj}_{\vec{b}}\vec{a} = \left(\frac{\vec{a}\cdot\vec{b}}{|\vec{b}|^2}\right)\vec{b} = \bigl(\operatorname{comp}_{\vec{b}}\vec{a}\bigr)\,\hat{\mathbf{b}},
\]
where $\hat{\mathbf{b}} = \vec{b}/|\vec{b}|$ is the unit vector in the direction of $\vec{b}$.
\end{definition}

\begin{propriete}[Geometric Interpretation]
If $\theta$ is the angle between $\vec{a}$ and $\vec{b}$ ($0 \le \theta \le \pi$), then $\vec{a}\cdot\vec{b} = |\vec{a}||\vec{b}|\cos\theta$. Hence
\[
\operatorname{comp}_{\vec{b}}\vec{a} = |\vec{a}|\cos\theta,\qquad
\operatorname{proj}_{\vec{b}}\vec{a} = (|\vec{a}|\cos\theta)\,\hat{\mathbf{b}}.
\]
The scalar component is positive for acute angles, zero for perpendicular vectors, and negative for obtuse angles.
\end{propriete}

\begin{exoresolu}[Exercise 1 — Scalar and vector components]
Given $\vec{a} = 3\mathbf{i} - 2\mathbf{j} + \mathbf{k}$ and $\vec{b} = \mathbf{i} + 4\mathbf{j} - 2\mathbf{k}$, find:
\begin{enumerate}
\item the scalar component of $\vec{a}$ in the direction of $\vec{b}$,
\item the vector component of $\vec{a}$ in the direction of $\vec{b}$,
\item the angle between $\vec{a}$ and $\vec{b}$.
\end{enumerate}
\tcblower
\textbf{Step 1: Magnitudes and scalar product.}
\begin{align*}
|\vec{a}| &= \sqrt{3^2 + (-2)^2 + 1^2} = \sqrt{14},\\
|\vec{b}| &= \sqrt{1^2 + 4^2 + (-2)^2} = \sqrt{21},\\
\vec{a}\cdot\vec{b} &= (3)(1) + (-2)(4) + (1)(-2) = 3 - 8 - 2 = -7.
\end{align*}

\textbf{Step 2: Scalar component of $\vec{a}$ in the direction of $\vec{b}$.}
\[
\operatorname{comp}_{\vec{b}}\vec{a} = \frac{\vec{a}\cdot\vec{b}}{|\vec{b}|} = \frac{-7}{\sqrt{21}} = -\frac{7\sqrt{21}}{21} = -\frac{\sqrt{21}}{3} \approx -1.53.
\]

\textbf{Step 3: Vector component of $\vec{a}$ in the direction of $\vec{b}$.}
The unit vector along $\vec{b}$ is $\hat{\mathbf{b}} = \dfrac{\vec{b}}{|\vec{b}|} = \dfrac{1}{\sqrt{21}}(\mathbf{i} + 4\mathbf{j} - 2\mathbf{k})$. Hence
\[
\operatorname{proj}_{\vec{b}}\vec{a} = \frac{\vec{a}\cdot\vec{b}}{|\vec{b}|^2}\,\vec{b} = \frac{-7}{21}(\mathbf{i} + 4\mathbf{j} - 2\mathbf{k}) = -\frac{1}{3}\mathbf{i} - \frac{4}{3}\mathbf{j} + \frac{2}{3}\mathbf{k}.
\]

\textbf{Step 4: Angle between $\vec{a}$ and $\vec{b}$.}
\[
\cos\theta = \frac{\vec{a}\cdot\vec{b}}{|\vec{a}||\vec{b}|} = \frac{-7}{\sqrt{14}\sqrt{21}} = \frac{-7}{\sqrt{294}} = \frac{-7}{7\sqrt{6}} = -\frac{1}{\sqrt{6}}.
\]
\[
\theta = \arccos\left(-\frac{1}{\sqrt{6}}\right) \approx 114.1^\circ.
\]

\textbf{Answers:} scalar component $= -\dfrac{\sqrt{21}}{3}$; vector component $= -\dfrac{1}{3}\mathbf{i} - \dfrac{4}{3}\mathbf{j} + \dfrac{2}{3}\mathbf{k}$; angle $\approx 114.1^\circ$. The negative scalar component correctly signals an obtuse angle.
\end{exoresolu}

\begin{definition}[Work Done by a Constant Force]
If a constant force $\vec{F}$ acts on a particle that undergoes a displacement $\vec{d}$, the \cle{work done} by the force is the scalar product
\[
W = \vec{F}\cdot\vec{d} = |\vec{F}|\,|\vec{d}|\cos\theta,
\]
where $\theta$ is the angle between $\vec{F}$ and $\vec{d}$. Work is a scalar quantity measured in joules (J). If the force varies or the path is curved, the work is computed by a line integral (not in this syllabus).
\end{definition}

\begin{attention}[Sign of Work]
Work is positive if the force has a component in the direction of displacement ($\theta < 90^\circ$), zero if perpendicular ($\theta = 90^\circ$), and negative if opposed to displacement ($\theta > 90^\circ$). Friction and weight (when lifting) typically do negative work.
\end{attention}

\begin{exoresolu}[Exercise 2 — Work done by a constant force]
A particle moves from $A(1,3)$ to $B(7,9)$ under the action of a constant force $\vec{F} = 5\mathbf{i} + 2\mathbf{j}\ \mathrm{N}$. A constant friction force of magnitude $4\ \mathrm{N}$ opposes the motion. Find:
\begin{enumerate}
\item the work done by $\vec{F}$,
\item the work done by friction,
\item the total work done on the particle.
\end{enumerate}
\tcblower
\textbf{Step 1: Displacement vector.}
\[
\overrightarrow{AB} = (7-1)\mathbf{i} + (9-3)\mathbf{j} = 6\mathbf{i} + 6\mathbf{j}\ \mathrm{m}.
\]

\textbf{Step 2: Work done by $\vec{F}$.}
\[
W_F = \vec{F}\cdot\overrightarrow{AB} = (5)(6) + (2)(6) = 30 + 12 = 42\ \mathrm{J}.
\]

\textbf{Step 3: Work done by friction.}
Friction opposes the motion, so it acts along $-\overrightarrow{AB}$. The distance moved is
\[
|\overrightarrow{AB}| = \sqrt{6^2 + 6^2} = 6\sqrt{2}\ \mathrm{m}.
\]
Hence
\[
W_{\text{friction}} = -4 \times 6\sqrt{2} = -24\sqrt{2} \approx -33.9\ \mathrm{J}.
\]

\textbf{Step 4: Total work.}
\[
W_{\text{total}} = 42 - 24\sqrt{2} \approx 42 - 33.9 = 8.1\ \mathrm{J}.
\]

\textbf{Answers:} $W_F = 42\ \mathrm{J}$; total work $= (42 - 24\sqrt{2})\ \mathrm{J} \approx 8.1\ \mathrm{J}$. Note that the work of a constant force depends only on the endpoints $A$ and $B$, not on the path taken.
\end{exoresolu}

\begin{definition}[Moment of a Force about a Point]
The \cle{moment} (or torque) of a force $\vec{F}$ acting at point $A$ about a point $O$ is the vector
\[
\vec{M}_O = \overrightarrow{OA} \times \vec{F}.
\]
Its magnitude is $|\vec{M}_O| = |\overrightarrow{OA}|\,|\vec{F}|\sin\theta = F \times d$, where $d$ is the perpendicular distance from $O$ to the line of action of $\vec{F}$. In two dimensions (forces in the $xy$-plane), the moment is a scalar $M_O = x F_y - y F_x$ (positive anticlockwise).
\end{definition}

\begin{propriete}[Equivalent Force–Couple System]
Any force $\vec{F}$ acting at $A$ can be replaced by an equal force $\vec{F}$ acting at $O$ together with a couple of moment $\vec{M}_O = \overrightarrow{OA} \times \vec{F}$. This is the \cle{principle of transmissibility} extended to a force–couple system.
\end{propriete}

\begin{exoresolu}[Exercise 3 — Moment of a force and equivalent system]
A force $\vec{F} = 4\mathbf{i} - 3\mathbf{j}\ \mathrm{N}$ acts at $A(2,5)$. Replace this force by an equivalent system consisting of a force through the origin $O$ and a couple.
\tcblower
\textbf{Principle.} Introduce at $O$ two equal and opposite forces $\vec{F}$ and $-\vec{F}$ (this changes nothing). The force $\vec{F}$ at $A$ and the force $-\vec{F}$ at $O$ form a couple; the remaining force $\vec{F}$ acts through $O$. Hence the system is equivalent to an equal force through $O$ together with a couple whose moment is the moment of the original force about $O$.

\textbf{Moment about $O$} (2D cross product, $M = xF_y - yF_x$):
\[
M_O = (2)(-3) - (5)(4) = -6 - 20 = -26\ \mathrm{N\,m}.
\]

\textbf{Answer:} the force is replaced by $\vec{F} = 4\mathbf{i} - 3\mathbf{j}\ \mathrm{N}$ acting through the origin, together with a couple of moment $-26\ \mathrm{N\,m}$, i.e. $26\ \mathrm{N\,m}$ clockwise.
\end{exoresolu}

\begin{propriete}[Equilibrium of Coplanar Forces]
A system of coplanar forces is in equilibrium if and only if:
\begin{enumerate}
\item The vector sum of the forces is zero: $\sum \vec{F} = \vec{0}$ (resolved components: $\sum F_x = 0$, $\sum F_y = 0$).
\item The sum of the moments about any point is zero: $\sum M_O = 0$.
\end{enumerate}
For three non-parallel forces in equilibrium, \cle{Lami's theorem} applies: each force is proportional to the sine of the angle between the other two.
\end{propriete}

\begin{exoresolu}[Exercise 4 — Equilibrium of coplanar forces]
Three forces $\vec{P}$, $\vec{Q}$, $\vec{R}$ act at a point. $\vec{P} = 50\ \mathrm{N}$ along the positive $x$-axis. $\vec{Q} = 40\ \mathrm{N}$ at $120^\circ$ to the positive $x$-axis. Find $\vec{R}$ for equilibrium.
\tcblower
\textbf{Step 1: Components of $P$ and $Q$.}
\[
\vec{P} = 50\mathbf{i},\qquad \vec{Q} = 40\cos 120^\circ\,\mathbf{i} + 40\sin 120^\circ\,\mathbf{j} = -20\mathbf{i} + 20\sqrt{3}\,\mathbf{j}.
\]

\textbf{Step 2: Equilibrium condition.} Since $\vec{P} + \vec{Q} + \vec{R} = \vec{0}$,
\[
\vec{R} = -(\vec{P} + \vec{Q}) = -(50 - 20)\mathbf{i} - 20\sqrt{3}\,\mathbf{j} = -30\mathbf{i} - 20\sqrt{3}\,\mathbf{j}.
\]

\textbf{Step 3: Magnitude and direction of $R$.}
\[
R = \sqrt{30^2 + (20\sqrt{3})^2} = \sqrt{900 + 1200} = \sqrt{2100} = 10\sqrt{21} \approx 45.8\ \mathrm{N}.
\]
Both components are negative, so $R$ points into the third quadrant. The angle below the negative $x$-axis is
\[
\tan\alpha = \frac{20\sqrt{3}}{30} = \frac{2\sqrt{3}}{3} \;\Longrightarrow\; \alpha \approx 49.1^\circ.
\]
So $R$ acts at $180^\circ + 49.1^\circ = 229.1^\circ$ from the positive $x$-axis (or $49.1^\circ$ below the negative $x$-axis).

\textbf{Step 4: Verification with Lami's theorem.} The angle between $P$ and $Q$ is $120^\circ$; the angle between $P$ and $R$ is $180^\circ - 49.1^\circ = 130.9^\circ$; the angle between $Q$ and $R$ is $360^\circ - 120^\circ - 130.9^\circ = 109.1^\circ$. Lami's theorem gives
\[
\frac{R}{\sin 120^\circ} = \frac{P}{\sin 109.1^\circ}
\;\Longrightarrow\;
R = \frac{50 \sin 120^\circ}{\sin 109.1^\circ} = \frac{50 \times 0.8660}{0.9455} \approx 45.8\ \mathrm{N}. \ \checkmark
\]

\textbf{Answer:} $R = 10\sqrt{21} \approx 45.8\ \mathrm{N}$ at $229.1^\circ$ to the positive $x$-axis; the two methods agree.
\end{exoresolu}

\courssub{I practise}

\begin{methode}[Work on an Inclined Plane]
For a body of mass $m$ on a plane inclined at angle $\alpha$:
\begin{enumerate}
\item Resolve weight: component down the plane $= mg\sin\alpha$, normal component $= mg\cos\alpha$.
\item If friction is present, $F_f = \mu R = \mu mg\cos\alpha$ (opposing motion).
\item Work done by a force $\vec{F}$ parallel to the plane over distance $d$: $W = F d$ (sign according to direction).
\item Work done against gravity: $W_g = mgh$ where $h = d\sin\alpha$ is the vertical rise.
\end{enumerate}
\end{methode}

\begin{exoresolu}[Exercise 5 — Work on an inclined plane]
A man pulls a $50\ \mathrm{kg}$ crate up a smooth plane inclined at $30^\circ$ over a distance of $8\ \mathrm{m}$. The pull is parallel to the plane. Find:
\begin{enumerate}
\item the force exerted by the man,
\item the work done against gravity,
\item the work done by the man.
\end{enumerate}
Take $g = 9.8\ \mathrm{m\,s^{-2}}$.
\tcblower
\textbf{Step 1: Force needed.} The component of the weight along the plane is
\[
mg\sin 30^\circ = 50 \times 9.8 \times 0.5 = 245\ \mathrm{N}.
\]
With friction neglected, the man must pull with $F = 245\ \mathrm{N}$ parallel to the plane.

\textbf{Step 2: Work done against gravity.} The vertical rise is $h = 8\sin 30^\circ = 4\ \mathrm{m}$:
\[
W_{\text{against gravity}} = mgh = 50 \times 9.8 \times 4 = 1960\ \mathrm{J}.
\]

\textbf{Step 3: Work done by the man.} The pull is parallel to the displacement, so
\[
W_{\text{man}} = F \times d = 245 \times 8 = 1960\ \mathrm{J}.
\]

\textbf{Answers:} work against gravity $= 1960\ \mathrm{J}$; work done by the man $= 1960\ \mathrm{J}$. The two are equal, as expected: with no friction and no change of speed, all the man's work becomes gravitational potential energy. This illustrates why the inclined plane is a machine — the force is halved ($245\ \mathrm{N}$ instead of $490\ \mathrm{N}$) but the distance is doubled.
\end{exoresolu}

\begin{methode}[Reactions on a Beam (Coplanar Forces)]
For a beam with vertical loads and supports:
\begin{enumerate}
\item Draw a free-body diagram showing all forces (weights, applied loads, reactions at supports).
\item Apply $\sum F_y = 0$ (vertical equilibrium).
\item Take moments about a convenient support to eliminate one reaction ($\sum M = 0$).
\item Solve for the reactions. A negative value means the reaction acts opposite to the assumed direction.
\end{enumerate}
\end{methode}

\begin{exoresolu}[Exercise 6 — Reactions on a beam]
A uniform beam $AB$ of length $6\ \mathrm{m}$ and weight $120\ \mathrm{N}$ rests on supports at $A$ and $C$, where $AC = 4\ \mathrm{m}$. A load of $300\ \mathrm{N}$ is attached at $B$. Find the reactions at $A$ and $C$.
\tcblower
Let $R_A$ and $R_C$ be the vertical reactions at $A$ and $C$ (assumed upward). The beam is uniform, so its weight $120\ \mathrm{N}$ acts at the midpoint, $3\ \mathrm{m}$ from $A$. The load $300\ \mathrm{N}$ acts at $B$, $6\ \mathrm{m}$ from $A$.

\textbf{Step 1: Vertical equilibrium.}
\[
R_A + R_C = 120 + 300 = 420\ \mathrm{N}. \qquad (1)
\]

\textbf{Step 2: Moments about $A$} (anticlockwise positive):
\[
R_C \times 4 - 120 \times 3 - 300 \times 6 = 0
\]
\[
4R_C = 360 + 1800 = 2160 \;\Longrightarrow\; R_C = 540\ \mathrm{N}.
\]

\textbf{Step 3: From (1),}
\[
R_A = 420 - 540 = -120\ \mathrm{N}.
\]

\textbf{Interpretation.} The negative sign means the support at $A$ must pull the beam \emph{down} with $120\ \mathrm{N}$ (a hold-down support): the heavy load at the overhanging end $B$ would otherwise tip the beam about $C$.

\textbf{Check — moments about $C$:} Taking anticlockwise positive about $C$:
\[
-R_A(4) + 120(1) - 300(2) = -(-120)(4) + 120 - 600 = 480 + 120 - 600 = 0. \ \checkmark
\]

\textbf{Answers:} $R_C = 540\ \mathrm{N}$ upwards; $R_A = 120\ \mathrm{N}$ downwards (the beam must be held down at $A$).
\end{exoresolu}

\begin{definition}[Couple]
A \cle{couple} consists of two equal, opposite, parallel forces $\vec{F}$ and $-\vec{F}$ separated by a perpendicular distance $d$. Its \cle{moment} is $M = F d$ (a scalar in 2D, a vector $\vec{M} = \vec{d} \times \vec{F}$ in 3D). The moment of a couple is a \cle{free vector}: it is the same about any point.
\end{definition}

\begin{propriete}[Moment of a Couple is Independent of Reference Point]
Let the forces $+\vec{F}$ and $-\vec{F}$ act at points $P$ and $Q$. For any point $X$,
\[
\vec{M}_X = \overrightarrow{XP} \times \vec{F} + \overrightarrow{XQ} \times (-\vec{F}) = (\overrightarrow{XP} - \overrightarrow{XQ}) \times \vec{F} = \overrightarrow{QP} \times \vec{F}.
\]
The point $X$ cancels out; the moment depends only on the forces and their separation.
\end{propriete}

\begin{exoresolu}[Exercise 7 — Moment of a couple]
A couple consists of two forces of magnitude $20\ \mathrm{N}$ acting on a light rod of length $0.5\ \mathrm{m}$. Find the moment of the couple and verify that it is the same about the midpoint and about one end of the rod.
\tcblower
\textbf{Step 1: Moment of the couple.}
\[
M = F \times d = 20 \times 0.5 = 10\ \mathrm{N\,m}.
\]

\textbf{Step 2: It is the same about any point.} (Proof given in the property above.)

\textbf{Numerical check.} Take the rod along the $x$-axis, $+\vec{F} = 20\mathbf{j}$ at $P(0.5, 0)$ and $-\vec{F}$ at $Q(0,0)$. About the midpoint $X(0.25, 0)$:
\[
M_X = (0.25)(20) + (0.25)(20) = 5 + 5 = 10\ \mathrm{N\,m}.
\]
About $Q$: $M_Q = (0.5)(20) + 0 = 10\ \mathrm{N\,m}$. Same value. $\checkmark$

\textbf{Answer:} $M = 10\ \mathrm{N\,m}$; the moment of a couple is a \emph{free vector}, independent of the reference point — this is the defining property of a couple.
\end{exoresolu}

\begin{propriete}[Equilibrium of a Particle in 3D]
A particle is in equilibrium under forces $\vec{F}_1, \vec{F}_2, \dots, \vec{F}_n$ iff the vector sum is zero:
\[
\sum_{i=1}^n \vec{F}_i = \vec{0} \quad\Longleftrightarrow\quad \sum F_{ix} = 0,\ \sum F_{iy} = 0,\ \sum F_{iz} = 0.
\]
Each force is usually given by its magnitude and direction (e.g., along a string). The direction is specified by a unit vector along the line of action.
\end{propriete}

\begin{exoresolu}[Exercise 8 — 3D equilibrium of a particle]
A particle at $P(1,1,1)$ is attached by light inextensible strings to fixed points $A(2,0,0)$, $B(0,3,0)$, $C(0,0,4)$. The tensions in $PA$ and $PB$ are $50\ \mathrm{N}$ and $40\ \mathrm{N}$ respectively. Find the tension in $PC$ for equilibrium, or show that equilibrium is impossible.
\tcblower
\textbf{Step 1: Direction vectors from $P$ to each anchor.}
\begin{align*}
\overrightarrow{PA} &= (2-1, 0-1, 0-1) = (1, -1, -1), & |\overrightarrow{PA}| &= \sqrt{3},\\
\overrightarrow{PB} &= (0-1, 3-1, 0-1) = (-1, 2, -1), & |\overrightarrow{PB}| &= \sqrt{6},\\
\overrightarrow{PC} &= (0-1, 0-1, 4-1) = (-1, -1, 3), & |\overrightarrow{PC}| &= \sqrt{11}.
\end{align*}

\textbf{Step 2: Tension vectors.}
\[
\vec{T}_A = \frac{50}{\sqrt{3}}(1, -1, -1),\qquad
\vec{T}_B = \frac{40}{\sqrt{6}}(-1, 2, -1),\qquad
\vec{T}_C = \frac{T_C}{\sqrt{11}}(-1, -1, 3).
\]

\textbf{Step 3: Equilibrium equations} $\vec{T}_A + \vec{T}_B + \vec{T}_C = \vec{0}$:
\[
\begin{cases}
\dfrac{50}{\sqrt{3}} - \dfrac{40}{\sqrt{6}} - \dfrac{T_C}{\sqrt{11}} = 0 & (x)\\[6pt]
-\dfrac{50}{\sqrt{3}} + \dfrac{80}{\sqrt{6}} - \dfrac{T_C}{\sqrt{11}} = 0 & (y)\\[6pt]
-\dfrac{50}{\sqrt{3}} - \dfrac{40}{\sqrt{6}} + \dfrac{3T_C}{\sqrt{11}} = 0 & (z)
\end{cases}
\]

\textbf{Step 4: Solve.} Numerically: $\dfrac{50}{\sqrt{3}} \approx 28.87$ and $\dfrac{40}{\sqrt{6}} \approx 16.33$.

From $(x)$: $\dfrac{T_C}{\sqrt{11}} = 28.87 - 16.33 = 12.54 \;\Longrightarrow\; T_C = 12.54\sqrt{11} \approx 41.6\ \mathrm{N}$.

Check $(y)$: $-28.87 + 32.66 - 12.54 = -8.75 \neq 0$. The three equations are inconsistent: with the given tensions the particle \emph{cannot} be in equilibrium.

\textbf{Consistency analysis.} Adding $(x)$ and $(y)$: $\dfrac{40}{\sqrt{6}} - \dfrac{2T_C}{\sqrt{11}} = 0 \Rightarrow \dfrac{T_C}{\sqrt{11}} = \dfrac{20}{\sqrt{6}} \approx 8.16$, which contradicts $(x)$. Subtracting, no value of $T_C$ satisfies all three equations.

\textbf{Answer:} the equilibrium equations are as above; solving them pairwise gives $T_C \approx 41.6\ \mathrm{N}$ from the $x$-equation, $T_C = 20\sqrt{11/6}\cdot\sqrt{...} \approx 27.1\ \mathrm{N}$ from the $y$-equation, and $T_C \approx 50.0\ \mathrm{N}$ from the $z$-equation. Since the three values differ, \textbf{no tension $T_C$ can produce equilibrium}: the data are inconsistent, and the particle would accelerate. (This is the correct rigorous conclusion; in an examination, showing the contradiction earns full credit.)
\end{exoresolu}

\courssub{I prepare for the GCE}

\begin{exoresolu}[Exercise 9 — Ladder against a smooth wall \textit{(12 marks)}]
A uniform ladder of length $L$ and weight $W$ rests with its foot on rough horizontal ground and its top against a smooth vertical wall. The ladder makes an angle $\theta$ with the horizontal. Find the minimum coefficient of friction $\mu$ between the ladder and the ground for the ladder not to slip.
\tcblower
\textbf{Setting up.} Forces on the ladder: weight $W$ at the midpoint ($L/2$ from the foot); normal reaction $S$ at the wall (horizontal, wall smooth); normal reaction $R$ and friction $F$ at the ground.

\begin{popfigure}
\begin{tikzpicture}[scale=0.8, >=stealth]
\draw[thick] (0,0) -- (4,3); % ladder
\draw[thick, pattern=north east lines] (-0.5,0) -- (5,0); % ground
\draw[thick] (4,0) -- (4,3.5); % wall
% forces
\draw[->, thick, popblue] (2,1.5) -- (2,0.5) node[midway, right] {$W$};
\draw[->, thick, poporange] (4,3) -- (3,3) node[midway, above] {$S$};
\draw[->, thick, popgreen] (0,0) -- (1,0) node[midway, below] {$F$};
\draw[->, thick, poppurple] (0,0) -- (0,1) node[midway, left] {$R$};
% angle
\draw (0.5,0) arc (0:36.87:0.5) node[midway, right] {$\theta$};
% dimensions
\draw[<->] (0,-0.5) -- (4,-0.5) node[midway, below] {$L\cos\theta$};
\draw[<->] (4.2,0) -- (4.2,3) node[midway, right] {$L\sin\theta$};
\end{tikzpicture}
\legende{Forces on a ladder against a smooth wall}
\end{popfigure}

\textbf{Resolving horizontally:}
\[
F = S. \qquad (1)
\]

\textbf{Resolving vertically:}
\[
R = W. \qquad (2)
\]

\textbf{Moments about the foot of the ladder} (eliminating $R$ and $F$):
\[
S \times L\sin\theta = W \times \frac{L}{2}\cos\theta
\;\Longrightarrow\;
S = \frac{W\cos\theta}{2\sin\theta} = \frac{W}{2}\cot\theta.
\]

Hence from (1): $F = \dfrac{W}{2}\cot\theta$.

\textbf{No-slip condition.} Limiting equilibrium requires $F \le \mu R$:
\[
\frac{W}{2}\cot\theta \le \mu W \;\Longrightarrow\; \mu \ge \frac{1}{2}\cot\theta.
\]
\[
\boxed{\mu_{\min} = \frac{1}{2}\cot\theta}
\]

\textbf{Indicative mark scheme (12 marks):}
\begin{center}
\begin{tabularx}{\linewidth}{|l|X|}\hline
\rowcolor{popblueL} Marks & Expected work \\ \hline
M1 & Correct force diagram: $W$ at midpoint, $S$ horizontal, $R$, $F$ at ground \\ \hline
M1 A1 & Resolving horizontally: $F = S$ \\ \hline
M1 A1 & Resolving vertically: $R = W$ \\ \hline
M1 & Taking moments about the foot (both terms, correct perpendicular distances) \\ \hline
A1 A1 & $S = \frac{W}{2}\cot\theta$ and $F = \frac{W}{2}\cot\theta$ \\ \hline
M1 & Using the friction condition $F \le \mu R$ \\ \hline
A1 & $\mu_{\min} = \frac{1}{2}\cot\theta$, correctly concluded \\ \hline
A1 & Clear statement of assumptions (smooth wall $\Rightarrow$ no friction at top; uniform ladder) \\ \hline
\end{tabularx}
\end{center}

\textbf{Examiner expectations:} moments must be taken about the foot to eliminate the two unknowns $R$ and $F$ in one step; perpendicular distances $L\sin\theta$ and $\frac{L}{2}\cos\theta$ must be justified; the inequality direction must be handled carefully. A common error is writing $F = \mu R$ from the start without stating that equilibrium requires $F \le \mu R$.
\end{exoresolu}

\begin{exoresolu}[Exercise 10 — Vector moment in 3D \textit{(10 marks)}]
A force $\vec{F} = 2\mathbf{i} + \mathbf{j} - 3\mathbf{k}\ \mathrm{N}$ acts at a point with position vector $\vec{r} = \mathbf{i} - 2\mathbf{j} + \mathbf{k}\ \mathrm{m}$ relative to the origin $O$. Find:
\begin{enumerate}
\item the vector moment of $\vec{F}$ about $O$,
\item the magnitude of this moment,
\item the perpendicular distance from $O$ to the line of action of $\vec{F}$.
\end{enumerate}
\tcblower
\textbf{Step 1: Vector moment about $O$.}
\[
\vec{M}_O = \vec{r} \times \vec{F} =
\begin{pmatrix} 1 \\ -2 \\ 1 \end{pmatrix}
\times
\begin{pmatrix} 2 \\ 1 \\ -3 \end{pmatrix}
=
\begin{pmatrix}
(-2)(-3) - (1)(1) \\
(1)(2) - (1)(-3) \\
(1)(1) - (-2)(2)
\end{pmatrix}
=
\begin{pmatrix}
6 - 1 \\
2 + 3 \\
1 + 4
\end{pmatrix}
=
\begin{pmatrix}
5 \\ 5 \\ 5
\end{pmatrix}\ \mathrm{N\,m}.
\]
So $\vec{M}_O = 5\mathbf{i} + 5\mathbf{j} + 5\mathbf{k}\ \mathrm{N\,m}$.

\textbf{Step 2: Magnitude.}
\[
|\vec{M}_O| = \sqrt{25 + 25 + 25} = 5\sqrt{3} \approx 8.66\ \mathrm{N\,m}.
\]

\textbf{Step 3: Perpendicular distance from $O$ to the line of action.} Since $|\vec{M}_O| = |\vec{F}|\,d$,
\[
|\vec{F}| = \sqrt{4 + 1 + 9} = \sqrt{14}\ \mathrm{N},
\qquad
d = \frac{|\vec{M}_O|}{|\vec{F}|} = \frac{5\sqrt{3}}{\sqrt{14}} = \frac{5\sqrt{42}}{14} \approx 2.32\ \mathrm{m}.
\]

\textbf{Indicative mark scheme (10 marks):}
\begin{center}
\begin{tabularx}{\linewidth}{|l|X|}\hline
\rowcolor{popblueL} Marks & Expected work \\ \hline
M1 & Setting up $\vec{M}_O = \vec{r} \times \vec{F}$ (correct order!) \\ \hline
M1 & Correct expansion of the cross product (determinant or component formula) \\ \hline
A2 & All three components correct: $(5, 5, 5)$ \\ \hline
M1 A1 & Magnitude $5\sqrt{3}\ \mathrm{N\,m}$ \\ \hline
M1 & Using $|\vec{M}_O| = |\vec{F}| d$ \\ \hline
A1 & $|\vec{F}| = \sqrt{14}$ \\ \hline
A2 & $d = \dfrac{5\sqrt{3}}{\sqrt{14}} \approx 2.32\ \mathrm{m}$ with units \\ \hline
\end{tabularx}
\end{center}

\textbf{Examiner expectations:} the order $\vec{r} \times \vec{F}$ (not $\vec{F} \times \vec{r}$) is essential — the wrong order loses the accuracy marks even if the magnitude is right; units ($\mathrm{N\,m}$, $\mathrm{m}$) must appear in the final answers.
\end{exoresolu}

\begin{definition}[Reduction of a System of Forces in 3D]
Any system of forces $\vec{F}_i$ acting at points $\vec{r}_i$ can be reduced to a single force $\vec{R} = \sum \vec{F}_i$ acting at a chosen point $O$ together with a couple $\vec{G}_O = \sum \vec{r}_i \times \vec{F}_i$. The scalar invariant $\vec{R}\cdot\vec{G}_O$ is independent of the choice of $O$. If $\vec{R}\cdot\vec{G}_O = 0$, the system reduces to a single force (the line of action is the \cle{central axis}); otherwise it reduces to a \cle{wrench} (a force plus a couple parallel to it).
\end{definition}

\begin{propriete}[Central Axis of a Wrench]
For a system with resultant force $\vec{R}$ and resultant moment $\vec{G}_O$ about $O$ (with $\vec{R}\cdot\vec{G}_O \neq 0$):
\begin{itemize}
\item The \cle{pitch} of the wrench is $p = \dfrac{\vec{R}\cdot\vec{G}_O}{|\vec{R}|^2}$.
\item The component of the couple parallel to $\vec{R}$ is $p\vec{R}$.
\item A point on the central axis is $\vec{r}_0 = \dfrac{\vec{R} \times \vec{G}_O}{|\vec{R}|^2}$.
\item The central axis is the line $\vec{r} = \vec{r}_0 + \lambda \vec{R}$.
\end{itemize}
\end{propriete}

\begin{exoresolu}[Exercise 11 — System of forces in 3D \textit{(13 marks)}]
Forces $\vec{F}_1 = 3\mathbf{i} - 2\mathbf{j} + 4\mathbf{k}\ \mathrm{N}$, $\vec{F}_2 = -2\mathbf{i} + 5\mathbf{j} - 3\mathbf{k}\ \mathrm{N}$, $\vec{F}_3 = \mathbf{i} + 3\mathbf{j} - 2\mathbf{k}\ \mathrm{N}$ act at points $A(1,0,2)$, $B(0,2,1)$, $C(2,1,0)$ respectively. Reduce the system to its simplest form at the origin $O$. Determine whether it can be reduced to a single force. If not, find the pitch and the vector equation of the central axis of the wrench.
\tcblower
\textbf{Step 1: Resultant force.}
\[
\vec{R} = \vec{F}_1 + \vec{F}_2 + \vec{F}_3 = (3 - 2 + 1,\; -2 + 5 + 3,\; 4 - 3 - 2) = (2,\; 6,\; -1)\ \mathrm{N}.
\]
\[
\vec{R} = 2\mathbf{i} + 6\mathbf{j} - \mathbf{k}\ \mathrm{N}.
\]

\textbf{Step 2: Resultant moment about $O$} — compute each $\vec{r}_i \times \vec{F}_i$.

For $\vec{F}_1$ at $A(1,0,2)$:
\[
\vec{r}_1 \times \vec{F}_1 =
\begin{pmatrix} 1 \\ 0 \\ 2 \end{pmatrix} \times
\begin{pmatrix} 3 \\ -2 \\ 4 \end{pmatrix} =
\begin{pmatrix} (0)(4)-(2)(-2) \\ (2)(3)-(1)(4) \\ (1)(-2)-(0)(3) \end{pmatrix} =
\begin{pmatrix} 4 \\ 2 \\ -2 \end{pmatrix}.
\]

For $\vec{F}_2$ at $B(0,2,1)$:
\[
\vec{r}_2 \times \vec{F}_2 =
\begin{pmatrix} 0 \\ 2 \\ 1 \end{pmatrix} \times
\begin{pmatrix} -2 \\ 5 \\ -3 \end{pmatrix} =
\begin{pmatrix} (2)(-3)-(1)(5) \\ (1)(-2)-(0)(-3) \\ (0)(5)-(2)(-2) \end{pmatrix} =
\begin{pmatrix} -11 \\ -2 \\ 4 \end{pmatrix}.
\]

For $\vec{F}_3$ at $C(2,1,0)$:
\[
\vec{r}_3 \times \vec{F}_3 =
\begin{pmatrix} 2 \\ 1 \\ 0 \end{pmatrix} \times
\begin{pmatrix} 1 \\ 3 \\ -2 \end{pmatrix} =
\begin{pmatrix} (1)(-2)-(0)(3) \\ (0)(1)-(2)(-2) \\ (2)(3)-(1)(1) \end{pmatrix} =
\begin{pmatrix} -2 \\ 4 \\ 5 \end{pmatrix}.
\]

Sum:
\[
\vec{G}_O = (4 - 11 - 2,\; 2 - 2 + 4,\; -2 + 4 + 5) = (-9,\; 4,\; 7)\ \mathrm{N\,m}.
\]

\textbf{Reduction at $O$:} a single force $\vec{R} = 2\mathbf{i} + 6\mathbf{j} - \mathbf{k}\ \mathrm{N}$ through $O$ together with a couple $\vec{G}_O = -9\mathbf{i} + 4\mathbf{j} + 7\mathbf{k}\ \mathrm{N\,m}$.

\textbf{Step 3: Can it reduce to a single force?} Compute the invariant $\vec{R}\cdot\vec{G}_O$:
\[
\vec{R}\cdot\vec{G}_O = (2)(-9) + (6)(4) + (-1)(7) = -18 + 24 - 7 = -1 \neq 0.
\]
Since $\vec{R}\cdot\vec{G}_O \neq 0$, the system \textbf{cannot} be reduced to a single force; its simplest form is a \textbf{wrench} (force plus a couple parallel to it).

\textbf{Step 4: Axis of the wrench (central axis).} The pitch of the wrench is
\[
p = \frac{\vec{R}\cdot\vec{G}_O}{|\vec{R}|^2} = \frac{-1}{4 + 36 + 1} = -\frac{1}{41}\ \mathrm{m}.
\]
The parallel component of the couple is $p\vec{R} = -\dfrac{1}{41}(2, 6, -1)\ \mathrm{N\,m}$. A point of the central axis is given by
\[
\vec{r}_0 = \frac{\vec{R} \times \vec{G}_O}{|\vec{R}|^2}.
\]
\[
\vec{R} \times \vec{G}_O =
\begin{pmatrix} 2 \\ 6 \\ -1 \end{pmatrix} \times
\begin{pmatrix} -9 \\ 4 \\ 7 \end{pmatrix} =
\begin{pmatrix} (6)(7)-(-1)(4) \\ (-1)(-9)-(2)(7) \\ (2)(4)-(6)(-9) \end{pmatrix} =
\begin{pmatrix} 46 \\ -5 \\ 62 \end{pmatrix},
\]
\[
\vec{r}_0 = \frac{1}{41}(46, -5, 62)\ \mathrm{m}.
\]
The central axis is the line through $\vec{r}_0$ in the direction of $\vec{R}$:
\[
\vec{r} = \frac{1}{41}\begin{pmatrix} 46 \\ -5 \\ 62 \end{pmatrix} + \lambda \begin{pmatrix} 2 \\ 6 \\ -1 \end{pmatrix}, \qquad \lambda \in \mathbb{R}.
\]

\textbf{Indicative mark scheme (13 marks):}
\begin{center}
\begin{tabularx}{\linewidth}{|l|X|}\hline
\rowcolor{popblueL} Marks & Expected work \\ \hline
M1 A1 & Resultant force $\vec{R} = (2, 6, -1)\ \mathrm{N}$ \\ \hline
M1 & Method: $\vec{G}_O = \sum \vec{r}_i \times \vec{F}_i$ \\ \hline
A3 & Each cross product correct (A1 each): $(4,2,-2)$, $(-11,-2,4)$, $(-2,4,5)$ \\ \hline
A1 & $\vec{G}_O = (-9, 4, 7)\ \mathrm{N\,m}$ \\ \hline
M1 & Testing the invariant $\vec{R}\cdot\vec{G}_O$ \\ \hline
A1 & $\vec{R}\cdot\vec{G}_O = -1 \neq 0$, correct conclusion: no single force, system is a wrench \\ \hline
M1 & Computing the pitch $p = -1/41\ \mathrm{m}$ \\ \hline
M1 A1 & Point of the axis $\vec{r}_0 = \frac{1}{41}(46, -5, 62)$ and vector equation of the central axis \\ \hline
\end{tabularx}
\end{center}

\textbf{Examiner expectations:} cross products must be shown with full working; the conclusion about reducibility must be \emph{justified by the invariant} $\vec{R}\cdot\vec{G}_O$, not guessed; when a single force is impossible, the candidate should name the wrench and give its axis. The most common errors are sign slips in cross products and forgetting that $\vec{R}\cdot\vec{G}_O$ is invariant under change of reduction point.
\end{exoresolu}

\newpage

\coursec{Summary sheet}

\begin{popfigure}
\begin{tikzpicture}[grow cyclic, mindmap, concept color=popblue,
  level 1 concept/.append style={level distance=4.4cm, sibling angle=60, font=\small\bfseries},
  level 2 concept/.append style={level distance=2.6cm, font=\scriptsize},
  every node/.style={concept, circular drop shadow}]
\node[concept, font=\bfseries, text=white, fill=popdark, minimum size=2.6cm]{Scalar \& Vector Products}
  child[concept color=poporange]{ node[concept]{Component in a Direction}
    child{ node[concept]{$\vec{a}\cdot\hat{u}$} }
    child{ node[concept]{Vector comp.\ $(\vec{a}\cdot\hat{u})\hat{u}$} }
  }
  child[concept color=popgreen]{ node[concept]{Work Done}
    child{ node[concept]{$W=\vec{F}\cdot\vec{d}$} }
    child{ node[concept]{$W=Fd\cos\theta$} }
  }
  child[concept color=poppurple]{ node[concept]{Moment of a Force}
    child{ node[concept]{$\vec{M}=\vec{r}\times\vec{F}$} }
    child{ node[concept]{$M=Fd$ perp.} }
  }
  child[concept color=popgold]{ node[concept]{Couples}
    child{ node[concept]{Moment $=Fd$} }
    child{ node[concept]{Zero resultant} }
  }
  child[concept color=poppink]{ node[concept]{Coplanar Forces}
    child{ node[concept]{$\sum F_x=0$} }
    child{ node[concept]{$\sum F_y=0$, $\sum M=0$} }
  }
  child[concept color=popblue]{ node[concept]{3D Systems}
    child{ node[concept]{$\sum\vec{F}=\vec{0}$} }
    child{ node[concept]{$\sum\vec{M}=\vec{0}$} }
  };
\end{tikzpicture}
\legende{Mind map of Chapter FMA15: applications of the scalar and vector products.}
\end{popfigure}

\begin{bilan}
\textbf{Key definitions}
\begin{itemize}
\item The \cle{scalar component} of $\vec{a}$ in the direction of $\vec{b}$ is $a\cos\theta=\dfrac{\vec{a}\cdot\vec{b}}{|\vec{b}|}=\vec{a}\cdot\hat{b}$; the \cle{vector component} is $\left(\vec{a}\cdot\hat{b}\right)\hat{b}$.
\item \cle{Work} done by a constant force $\vec{F}$ over displacement $\vec{d}$: $W=\vec{F}\cdot\vec{d}=Fd\cos\theta$ (joules).
\item The \cle{moment} of $\vec{F}$ about point $O$: $\vec{M}_O=\vec{r}\times\vec{F}$, where $\vec{r}$ is the position vector from $O$ to any point of the line of action; magnitude $M=Fd$ ($d$ = perpendicular distance), unit $\mathrm{N\,m}$.
\item A \cle{couple}: two equal, opposite, parallel forces with different lines of action; resultant zero, moment $M=Fd$ ($d$ = distance between the lines of action), the same about \emph{every} point.
\item A system of forces is in \cle{equilibrium} when $\sum\vec{F}=\vec{0}$ and $\sum\vec{M}_O=\vec{0}$ (about any point $O$).
\end{itemize}

\textbf{Formulas and theorems to know}
\begin{itemize}
\item Scalar product: $\vec{a}\cdot\vec{b}=a_1b_1+a_2b_2+a_3b_3=ab\cos\theta$; $\vec{a}\perp\vec{b}\iff\vec{a}\cdot\vec{b}=0$.
\item Vector product: $\vec{a}\times\vec{b}=(a_2b_3-a_3b_2,\ a_3b_1-a_1b_3,\ a_1b_2-a_2b_1)$; $|\vec{a}\times\vec{b}|=ab\sin\theta$; anti-commutative: $\vec{b}\times\vec{a}=-\vec{a}\times\vec{b}$.
\item Component perpendicular to $\vec{b}$: $\vec{a}-(\vec{a}\cdot\hat{b})\hat{b}$.
\item Work: $W>0$ (force helps motion), $W=0$ if $\vec{F}\perp\vec{d}$, $W<0$ (force opposes motion, e.g.\ friction).
\item Varignon's theorem: the moment of the resultant equals the sum of the moments of the components.
\item A single force $\vec{F}$ acting at $A$ is equivalent to the same force at $O$ \textbf{plus} a couple of moment $\vec{r}_{OA}\times\vec{F}$.
\item Equilibrium of a rigid body under \textbf{three} coplanar forces: the three lines of action are concurrent (or parallel).
\item Lami's theorem (three concurrent forces in equilibrium): $\dfrac{F_1}{\sin\alpha_1}=\dfrac{F_2}{\sin\alpha_2}=\dfrac{F_3}{\sin\alpha_3}$.
\item In 3D equilibrium, resolve along three axes: $\sum F_x=\sum F_y=\sum F_z=0$ and $\sum M_x=\sum M_y=\sum M_z=0$.
\end{itemize}

\textbf{Methods}
\begin{itemize}
\item \textbf{Component along a direction:} form the unit vector $\hat{u}=\vec{u}/|\vec{u}|$, then compute $\vec{a}\cdot\hat{u}$.
\item \textbf{Work:} identify $\vec{F}$ and the displacement $\vec{d}$ (not the path), then dot them; check the sign.
\item \textbf{Moment in 2D:} choose the pivot, draw the perpendicular distance $d$, then $M=Fd$ with sign (anticlockwise positive by convention).
\item \textbf{Moment in 3D:} write $\vec{r}$ from the pivot to a point of the line of action and compute the cross product.
\item \textbf{Equilibrium problems:} (1) draw a clear force diagram; (2) resolve in two (or three) perpendicular directions; (3) take moments about a point where unknown forces pass to eliminate them; (4) solve the system.
\item \textbf{Reducing a system:} resultant $\vec{R}=\sum\vec{F}_i$; resultant moment $\vec{M}_O=\sum\vec{r}_i\times\vec{F}_i$. If $\vec{R}=\vec{0}$ and $\vec{M}\neq\vec{0}$, the system reduces to a couple.
\end{itemize}

\textbf{Common mistakes}
\begin{itemize}
\item Confusing scalar component (a number) with vector component (a vector).
\item Using the distance along the force instead of the \emph{perpendicular} distance in $M=Fd$.
\item Forgetting that $\vec{a}\times\vec{b}\neq\vec{b}\times\vec{a}$ (order matters).
\item Taking moments about different points in the same equation, or mixing sign conventions.
\item For a couple, using the distance to the pivot instead of the distance between the two lines of action.
\item Forgetting that work is zero when the force is perpendicular to the displacement (e.g.\ normal reaction on a horizontal road).
\item Omitting units: work in joules (J), moment in $\mathrm{N\,m}$ — these are \emph{not} the same physical quantity.
\end{itemize}

\textbf{GCE tips}
\begin{itemize}
\item Always state your sign convention for moments before writing equations.
\item Choose the pivot at the point of application of an unknown reaction: it disappears from the moment equation.
\item Show the force diagram: method marks are awarded for it at the GCE.
\item For three-force equilibrium, Lami's theorem is often the fastest route — but check the forces are concurrent first.
\item In 3D problems, work column by column with components $(x,y,z)$ to avoid sign slips in cross products.
\item Give final answers with correct units and sensible accuracy (2 or 3 significant figures).
\end{itemize}
\end{bilan}

\findecours

\end{document}
